% Standalone source: no external figures, bibliography, or input files.
% Compile with: latexmk -pdf A93_A95_E94_E95_neutral_ABL_EP_AWF_combined.tex
\documentclass[11pt,a4paper]{article}
\usepackage[T1]{fontenc}
\usepackage[utf8]{inputenc}
\usepackage{lmodern}
\usepackage[margin=25mm,headheight=14pt]{geometry}
\usepackage{amsmath,amssymb,amsthm,mathtools,bm}
\usepackage{booktabs,array,tabularx,longtable}
\usepackage{enumitem,microtype,xcolor,listings,fancyhdr}
\usepackage[numbers,sort&compress]{natbib}
\usepackage[colorlinks=true,linkcolor=blue!45!black,citecolor=blue!45!black,
urlcolor=blue!45!black]{hyperref}
\usepackage[nameinlink,noabbrev]{cleveref}
\hypersetup{
pdftitle={Neutral Atmospheric Surface-Layer Equilibria and Equilibrium-Preserving Analytical Wall Functions for A93, A95, and E94/E95},
pdfsubject={Original turbulence closures, minimal outer changes, finite-viscosity wall elements, kappa consistency, spectral/hp implementation, stability and positivity},
pdfauthor={Technical research note}}
\numberwithin{equation}{section}
\newtheorem{proposition}{Proposition}[section]
\newtheorem{theorem}[proposition]{Theorem}
\newtheorem{lemma}[proposition]{Lemma}
\newtheorem{corollary}[proposition]{Corollary}
\newtheorem{remark}[proposition]{Remark}
\newcommand{\eps}{\varepsilon}
\newcommand{\ust}{u_*}
\newcommand{\nut}{\nu_t}
\newcommand{\dd}{\mathrm d}
\newcommand{\qs}{q_*}
\newcommand{\cmus}{C_{\mu,*}}
\newcommand{\Id}{\boldsymbol I}
\newcommand{\R}{\boldsymbol R}
\newcommand{\Q}{\boldsymbol Q}
\newcommand{\Amat}{\boldsymbol A}
\newcommand{\Smat}{\boldsymbol S}
\newcommand{\Omat}{\boldsymbol\Omega}
\DeclareMathOperator{\dev}{dev}
\DeclareMathOperator{\tr}{tr}
\DeclareMathOperator{\diag}{diag}
\setlength{\parindent}{0pt}
\setlength{\parskip}{4pt}
\setlength{\emergencystretch}{2em}
\setlist[itemize]{topsep=3pt,itemsep=2pt}
\setlist[enumerate]{topsep=3pt,itemsep=3pt}
\renewcommand{\arraystretch}{1.15}
\newcolumntype{Y}{>{\raggedright\arraybackslash}X}
\pagestyle{fancy}\fancyhf{}
\fancyhead[L]{\small Neutral ABL equilibria and analytical wall elements}
\fancyhead[R]{\small A93, A95, E94/E95}
\fancyfoot[C]{\thepage}
\lstset{language=Python,basicstyle=\ttfamily\scriptsize,
keywordstyle=\color{blue!55!black},commentstyle=\color{black!55},
showstringspaces=false,breaklines=true,columns=fullflexible,
keepspaces=true,frame=single,rulecolor=\color{black!20}}
\title{\Large\bfseries Neutral Atmospheric Surface-Layer Equilibria\\[3pt]
and Equilibrium-Preserving Analytical Wall Functions\\[3pt]
for A93, A95, and E94/E95\\[7pt]
\large Original formulations, minimal changes, finite-viscosity wall elements,\\
spectral/$hp$ implementation, stability, and positivity}
\author{Standalone mathematical research note}
\date{4 November 2026}
\begin{document}
\maketitle
\begin{abstract}
This note derives exact neutral atmospheric surface-layer equilibria for
the Shih--Zhu--Lumley algebraic-stress models of 1993 and 1995 (A93 and A95)
and the Shih--Liou--Shabbir--Yang--Zhu eddy-viscosity model reported in 1994
and published in 1995 (E94/E95). The original stress laws determine three
constant-energy branches. Their original dissipation coefficients imply
the compatible logarithmic slopes
$0.42832363$, $0.40406100$, and $0.41208747$.
A prescribed common slope is obtained by changing one outer
dissipation-diffusivity coefficient for each model.

The outer equations use their stated high-Reynolds-number approximation;
the analytical wall element retains physical molecular viscosity.
Starting from the original Suga--Craft--Iacovides rough-wall kernel, the
note derives no-slip reconstruction, signed forcing response, full
traction matching, aerodynamic-roughness compatibility, a positive
dissipation continuation, and the conservative redistribution needed for
every spectral/$hp$ kinetic-energy test mode. Exact matching at a chosen
height is distinguished from a common logarithmic overlap. The selected
outer $\kappa$ enters all equilibrium-derived wall quantities; an
additional one-coefficient change of the auxiliary viscosity slope
produces the matching leading overlap. Neither option removes the
finite molecular/roughness intercept.

Conservative curved-wall equations, interface pressure and mass-flux
coupling, spatial transition, and a Nektar++ implementation specification
are developed explicitly. All three original constitutive laws have a
positive incompressible momentum principal bound at their neutral
reference, yet each admits an explicit negative-diffusion state away
from it. Covariance realizability, coefficient-domain restrictions,
transported-turbulence positivity, and discrete stationarity are assessed
separately. A bounded-rate inner destruction remedy preserves the
reference while requiring a consistent dependent-dissipation trace.
Executable checks and a staged validation programme distinguish the
proved continuum identities and local certificates from numerical and
physical validation still to be performed.
\end{abstract}
\newpage
{\small\tableofcontents}
\newpage

\section{Problem, assumptions, and notation}
\label{sec:scope}

\subsection{Objective and model identities}

The objective is a turbulence and wall-element formulation that admits
the analytical neutral surface-layer reference
\begin{equation}
 U_{o,0}(z)=\frac{\ust}{\kappa}\ln\frac{z}{z_0},
 \qquad H\le z\le Z,
 \label{scope:target}
\end{equation}
with constant reference turbulent kinetic energy, while enforcing
numerical no slip at the ground $z=0$. Here $H$ is a physical matching
height and $Z>H$ is the upper boundary. The wall-layer reconstruction
occupies $0\le z\le H$; the outer logarithm is not extended to the
numerical ground.

Four primary publications define three closures:
A93 is the quadratic algebraic-stress model in NASA TM-105993
\citep{Shih1993}; A95 is the later, distinct quadratic algebraic-stress
model \citep{Shih1995ASM}; E94/E95 is the variable-coefficient linear
eddy-viscosity model in NASA TM-106721 and its journal presentation
\citep{Shih1994EVM,Shih1995EVM}. The scalar analytical wall-function
(AWF) baseline is the rough-wall formulation of Suga, Craft, and
Iacovides \citep{SCI2006}. EP-AWF denotes the
equilibrium-preserving wall-element construction derived here.
Its dissipation redistribution, spectral/$hp$ adaptation, curved
completion, and optional coefficient recalibrations are identified
as additional model choices rather than attributed to those papers.

\subsection{Reference assumptions and the two viscosity regimes}

The reference is stationary, incompressible, constant-density, neutral,
nonrotating, horizontally homogeneous, and unidirectional. There is no
streamwise pressure gradient or Coriolis forcing, the Reynolds shear
stress is constant, and $k_0$ is independent of height. Hydrostatic
pressure may be absorbed into the pressure variable provided its body
force is removed consistently. These assumptions define a surface-layer
reference, not a complete rotating atmospheric boundary layer.

The selected \emph{outer} equations omit molecular diffusion and use the
high-Reynolds-number form of the dissipation sink. The \emph{inner}
wall element retains the same fixed physical kinematic viscosity
$\nu>0$ in momentum, scalar diffusion, wall units, and the original
dissipation baseline. This is a declared composite approximation.
It does not assert an exact logarithmic solution of the full
finite-Reynolds-number outer equations.

The numerical ground is stationary, impermeable, and no slip for mean
velocity. The modeled ground $k$ flux is zero. Imposing a wall-resolved
condition $k=0$ would exclude the constant-$k_0$ reference and would
define a different wall model. Inner dissipation is a dependent
profile; only outer dissipation is transported independently.

\subsection{Minimal changes and levels of exactness}

For a prescribed $\kappa$, the outer modification is minimized by
counting changed constant coefficients while holding the original
constitutive law and the canonical constant-$k$ family fixed. Changing
one dissipation-diffusivity coefficient is sufficient; retaining the
original compatible $\kappa_m$ needs no outer change. This restricted
minimality does not order every possible turbulence-model modification.

The main EP-AWF retains the original auxiliary viscosity and fixed-input
momentum response. Its additional dissipation and conservative
redistribution are derived from physical wall-layer balances.
A separate optional change matches the auxiliary leading logarithmic
slope. All coefficient changes are distinguished from a one-time
calibration of a wall roughness input.

Exact continuum equilibrium, exact finite-dimensional representation,
zero assembled spatial residual, and a stationary fully discrete
timestep are different properties. A positive Reynolds covariance is
also distinct from positive production, positive incremental momentum
diffusion, and positivity of transported $k$ and $\eps$.
The note gives separate tests and qualifications for these claims.

\subsection{Coordinates, indices, and notation}

Vectors and tensors are expressed in physical Cartesian coordinates
$(x,y,z)$, with $x$ streamwise, $y$ spanwise, and $z$ vertical on the
flat reference. Repeated Cartesian indices are summed. For curved
geometry, a normal distance $r$ and a surface parametrization are used
only to describe the mapping; vector components remain Cartesian.
The subscript $o$ denotes the outer region, $w$ the wall element,
$A$ the original auxiliary AWF quantity, and $c$ a corrected
dissipation continuation. A subscript $0$ means evaluation on the
selected equilibrium; it never means evaluation at the ground.

{\small
\renewcommand{\arraystretch}{1.05}
\begin{longtable}{@{}p{0.24\textwidth}p{0.71\textwidth}@{}}
\caption{Principal notation. Additional local symbols are defined where used.}
\label{tab:notation}\\
\toprule Symbol & Meaning\\ \midrule
\endfirsthead
\toprule Symbol & Meaning\\ \midrule
\endhead
$\boldsymbol U,k,\eps,\pi$ &
Mean velocity, turbulent kinetic energy, dissipation rate, and
kinematic pressure with a declared hydrostatic convention.\\
$R_{ij}=\overline{u_i'u_j'}$ &
Reynolds covariance; its negative is the modeled turbulent momentum stress.\\
$A_{ij}=\partial_jU_i$, $S,\Omega$ &
Velocity gradient, its symmetric strain, and antisymmetric rotation parts.\\
$\nu,\nut$ &
Physical molecular kinematic viscosity and outer eddy viscosity.\\
$\ust,z_0,\kappa$ &
Reference friction velocity, aerodynamic roughness length, and the
selected outer logarithmic slope constant.\\
$q_m,C_{\mu,*}$ &
Model-specific neutral root and its equilibrium viscosity coefficient,
$C_{\mu,*}=q_m^{-2}$; $m=93,95,E$.\\
$\kappa_m^{\rm nat}$ &
Compatible $\kappa$ implied by the original outer coefficients.\\
$K=\mathcal K[k_h]$, $K_0$ &
Declared current wall-energy sample; on the fixed reference
$K_0=k_0=q_m\ust^2$. The transported $k_h$ need not equal $K$
away from equilibrium.\\
$u_m(K),t_m(K),A_m(K)$ &
Counterfactual reference scales $\sqrt{K/q_m}$, $K/q_m$,
and $u_m^3/\kappa$ at the same current $K$.\\
$H,Z,h_r,z_m$ &
Matching height, outer top, original equivalent sand-grain roughness
input, and height by which the corrected dissipation tail must start.\\
$b,z_v,z_{\eps,A},z_c$ &
Dimensionless viscosity intercept, its dimensional location, original
dissipation cutoff, and corrected dissipation cutoff. These are distinct.\\
$c_\mu^w,c_\ell^w,\alpha$ &
Original auxiliary wall coefficients, with
$\alpha=c_\mu^wc_\ell^w$.\\
$c_{\ell,m}^{\rm eq}$ &
Effective coefficient $\kappa q_m^{3/2}$ of the corrected equilibrium
dissipation tail; it need not replace $c_\ell^w$.\\
$D_A,J_0,J_1,J_2$ &
Total inner momentum diffusivity and three resistance/forcing integrals;
the integer suffixes label integrands.\\
$\boldsymbol F,\boldsymbol t_A,\boldsymbol T_H$ &
Sampled tangential forcing, signed ground shear per density, and
signed upper-interface shear per density.\\
$p_{A,0},P_A,E_A,E_{c,m}$ &
Counterfactual production, current production, original dissipation
baseline, and corrected dissipation profile.\\
$E_{\rm pre},D_{\rm pre},C_{\rm pre}$ &
Dependent dissipation, effective destruction, and conservative
redistribution before the positivity adaptation.\\
$E_{\rm pos},D_{\rm pos},C_{\rm pos}$ &
Corresponding profiles after that adaptation.\\
$\chi$ &
Product of declared spatial/state/admissibility gates, constant along
each wall-normal ray.\\
$\boldsymbol n,\boldsymbol N_H,J,A_g$ &
Wall-to-fluid unit normal, actual matching-surface normal, normal-tube
Jacobian, and tangent stretch map.\\
\bottomrule
\end{longtable}}



\section{Original formulations and the selected outer equations}
\label{out:models}

The three constitutive models used in this note are the original
stationary-frame A93 and A95 algebraic-stress closures and the common
E94/E95 eddy-viscosity closure. Their coefficients are recorded before
any atmospheric-equilibrium recalibration. A93 and A95 have different
quadratic stress corrections; the 1994 report and 1995 journal paper for
E94/E95 describe the same closure \citep{Shih1993,Shih1995ASM,
Shih1994EVM,Shih1995EVM}. All formulas below specialize the publications
to incompressible motion in a nonrotating frame. Their additional
rotating-frame definitions are not needed for the reference considered
here.

\subsection{Constitutive notation and the full production}
\label{out:notation}

With the coordinate order $(x,y,z)$, define
\begin{equation}
 A_{ij}=\partial_jU_i,\qquad
 \Smat=\tfrac12(\Amat+\Amat^T),\qquad
 \Omat=\tfrac12(\Amat-\Amat^T),\qquad
 \tau=\frac{k}{\eps}.
 \label{out:gradient}
\end{equation}
The Reynolds tensor is the velocity covariance
$R_{ij}=\overline{u_i'u_j'}$, so $k=\tfrac12\tr\R$ and the turbulent
contribution to the mean Cauchy stress is $-\rho\R$. We use
\begin{align}
 s&=(\Smat:\Smat)^{1/2},&
 r&=(\Omat:\Omat)^{1/2},&
 \mathcal S&=\sqrt2s,&
 \mathcal U&=(s^2+r^2)^{1/2}, \label{out:norms}\\
 \dev\boldsymbol B&=\boldsymbol B-\tfrac13\tr(\boldsymbol B)\Id,
 &\boldsymbol B:\boldsymbol C&=B_{ij}C_{ij}. \nonumber
\end{align}
For $s>0$, the strain invariant and its coefficient are
\begin{equation}
 W=\frac{\tr(\Smat^3)}{s^3},\qquad
 A_s=\sqrt6\cos\left[\frac13\arccos(\sqrt6W)\right].
 \label{out:invariants}
\end{equation}
For a symmetric trace-free $\Smat$, $|\sqrt6W|\le1$, and the equivalent
identity
\begin{equation}
 A_s=\frac{3\lambda_{\max}(\Smat)}s
 \label{out:eigenAs}
\end{equation}
uses the largest algebraic strain eigenvalue. The neutral reference has
$s>0$ and distinct strain eigenvalues; no value of the indeterminate
ratio $W$ at zero strain is required in its construction.

Every model has the form
\begin{equation}
 \boxed{\R=\frac23k\Id-2\nut\Smat+\Q,\qquad
 \nut=C_\mu\frac{k^2}{\eps}=C_\mu k\tau.}
 \label{out:stress}
\end{equation}
Its kinetic-energy production is computed from the full covariance:
\begin{equation}
 P_k=-\R:\Amat=-\R:\Smat
     =2\nut s^2-\Q:\Smat.
 \label{out:production}
\end{equation}
The last equality uses incompressibility and the symmetry of $\R$.
It is not permissible to discard $\Q:\Smat$ from a general A93
production calculation. On the simple-shear branch derived below it
vanishes for all three models.

\subsection{A93: the original three-tensor quadratic correction}
\label{out:A93}

In the stationary specialization of A93, the invariants and viscosity
coefficient are
\begin{equation}
 \eta=\sqrt2s\tau,\qquad \xi=\sqrt2r\tau,\qquad
 \boxed{C_\mu^{93}=\frac{2/3}{1.25+\eta+0.9\xi}.}
 \label{out:A93C}
\end{equation}
The original quadratic law, with its published coefficients, is
\begin{equation}
\begin{split}
 \Q^{93}=\frac{k\tau^2}{1000+\eta^3}\big[&
 -4\dev\{\Amat^2+(\Amat^2)^T\}\\
 &+13\dev(\Amat\Amat^T)-2\dev(\Amat^T\Amat)\big].
\end{split}
\label{out:A93QA}
\end{equation}
These are the final model in Eq.~(47), its viscosity relation in
Eq.~(50), and the constants in Eqs.~(52)--(53) of
\citet{Shih1993}; the strain and rotation invariants are specified
in Eqs.~(27)--(28) and (45). The nonlinear denominator remains
$1000+\eta^3$: the rotation invariant $\xi$ enters $C_\mu$, not that
denominator.

Expanding $\Amat=\Smat+\Omat$ gives the exactly equivalent form
\begin{equation}
 \boxed{
 \Q^{93}=\frac{k\tau^2}{1000+\eta^3}
 \left[3\dev(\Smat^2)-19\dev(\Omat^2)
       +15(\Omat\Smat-\Smat\Omat)\right].}
 \label{out:A93Q}
\end{equation}
For example,
$\Amat\Amat^T=\Smat^2-\Omat^2+
\Omat\Smat-\Smat\Omat$, which fixes the commutator sign for the
velocity-gradient convention in \cref{out:gradient}. The tensor basis
is quadratic even though its scalar coefficient contains $\eta^3$.

\subsection{A95: the original strain--rotation correction}
\label{out:A95}

The final A95 constitutive relation is
\begin{equation}
 \boxed{C_\mu^{95}=\frac1{6.5+A_s\mathcal U\tau},\qquad
 \Q^{95}=2C_Qk\tau^2(\Omat\Smat-\Smat\Omat),}
 \label{out:A95law}
\end{equation}
where
\begin{equation}
 C_Q=\frac{\sqrt{1-9(C_\mu^{95})^2(s\tau)^2}}
               {1+6(s\tau)(r\tau)}.
 \label{out:CQ}
\end{equation}
The provenance is Eqs.~(14), (16.1)--(16.2), and (17) of
\citet{Shih1995ASM}; the norms are those of Eqs.~(3.1)--(3.4),
and $A_s$ is given in Eq.~(11.3). The original constants
$A_0=6.5$ and $C_0=1$ have already been inserted.
Here $C_Q$ is the paper's quadratic coefficient $C_2$, renamed to
avoid confusion with a dissipation destruction coefficient.

The square root in \cref{out:CQ} is part of the domain of the printed
law. The reference will satisfy its strict positivity condition.
No clipping or continuation of a negative radicand is included in
\cref{out:A95law,out:CQ}. Also,
$\tr[(\Omat\Smat-\Smat\Omat)\Smat]=0$, so
$\Q^{95}:\Smat=0$ wherever this coefficient is real. This production
identity does not remove the quadratic correction from the normal
stresses or from the constitutive tangent.

\subsection{E94/E95: the original variable eddy-viscosity relation}
\label{out:E}

The two E publications use
\begin{equation}
 \boxed{C_\mu^E=\frac1{4.0+A_s\mathcal U\tau},\qquad
        \Q^E=\boldsymbol0.}
 \label{out:Elaw}
\end{equation}
This is the formulation in Eqs.~(17), (19)--(21) of the journal
presentation, with $A_0=4.0$ in its Table~2 and boundary-layer
calibration \citep{Shih1995EVM}; the report gives the corresponding
model \citep{Shih1994EVM}. The value $4.0$ is retained exactly.
Although the retained tensor is the linear Boussinesq strain tensor,
$C_\mu^E$ depends nonlinearly on the velocity gradient.

\subsection{Published transport equations and the outer reduction}
\label{out:transport}

It is useful to distinguish the complete published transport equations
from the outer approximation used in the rest of this note. For constant
molecular kinematic viscosity $\nu$, the common incompressible momentum
and kinetic-energy equations are
\begin{align}
 \nabla\cdot\boldsymbol U&=0,
 &\frac{D\boldsymbol U}{Dt}
   &=-\nabla\pi+\nu\Delta\boldsymbol U-\nabla\cdot\R,
   \label{out:originalmomentum}\\
 \frac{Dk}{Dt}
   &=\nabla\cdot\left[\left(\nu+\frac{\nut}{\sigma_k}\right)
                  \nabla k\right]+P_k-\eps,
 &&\sigma_k=1.
 \label{out:originalk}
\end{align}
Here $D/Dt=\partial_t+\boldsymbol U\cdot\nabla$ and $\pi$ is pressure
per unit mass with the hydrostatic contribution absorbed. With gravity
$-g_{\rm grav}\boldsymbol e_z$, for example,
$\pi=p/\rho+g_{\rm grav}z$.

A93 and A95 retain the standard dissipation transport structure,
\begin{equation}
 \frac{D\eps}{Dt}=
 \nabla\cdot\left[\left(\nu+\frac{\nut}{\sigma_\eps}\right)
                   \nabla\eps\right]
 +C_{\eps1}\frac{\eps}{k}P_k-C_{\eps2}\frac{\eps^2}{k},
 \label{out:originalepsA}
\end{equation}
with the original values
\begin{equation}
 C_{\eps1}=1.44,\qquad C_{\eps2}=1.92,\qquad
 \sigma_{\eps,\rm org}=1.3.
 \label{out:originalAconstants}
\end{equation}
See Eqs.~(48)--(52) of \citet{Shih1993} and Eqs.~(20)--(23) of
\citet{Shih1995ASM}.

E94/E95 instead uses
\begin{equation}
 \frac{D\eps}{Dt}=
 \nabla\cdot\left[\left(\nu+\frac{\nut}{\sigma_\eps}\right)
                   \nabla\eps\right]
 +C_{\eps S}(\eta)\mathcal S\eps
 -C_{\eps2}\frac{\eps^2}{k+\sqrt{\nu\eps}},
 \label{out:originalepsE}
\end{equation}
where
\begin{equation}
 \eta=\mathcal S\tau,\qquad
 C_{\eps S}(\eta)=\max\left(0.43,\frac{\eta}{\eta+5}\right),
 \qquad C_{\eps2}=1.9,\qquad \sigma_{\eps,\rm org}=1.2.
 \label{out:originalEconstants}
\end{equation}
The source calls $C_{\eps S}$ and $C_{\eps2}$ simply $C_1$ and $C_2$.
The kinetic-energy and dissipation equations are Eqs.~(22)--(23),
the $C_1$ formula is Eq.~(25), and the diffusivity calibration is
Eq.~(28) and Table~2 in \citet{Shih1995EVM}; see also
\citet{Shih1994EVM}. Its dissipation production is
$C_{\eps S}\mathcal S\eps$, not
$C_{\eps1}(\eps/k)P_k$.

The selected \emph{outer} system omits the molecular momentum and
scalar-diffusion terms and uses $k$ in the E94/E95 dissipation
sink. The latter is the large-$k^2/(\nu\eps)$ reduction of the
published denominator. Explicitly, the equations used for the outer
reference and its later outer stability analysis are
\begin{align}
 \nabla\cdot\boldsymbol U&=0,
 &\frac{D\boldsymbol U}{Dt}&=-\nabla\pi-\nabla\cdot\R,
 \label{out:momentum}\\
 \frac{Dk}{Dt}&=\nabla\cdot\left(\frac{\nut}{\sigma_k}\nabla k\right)
                  +P_k-\eps,
 &&\sigma_k=1, \label{out:kouter}\\
 \frac{D\eps}{Dt}&=\nabla\cdot\left(\frac{\nut}{\sigma_\eps}\nabla\eps\right)
       +C_{\eps1}\frac{\eps}{k}P_k-C_{\eps2}\frac{\eps^2}{k},
 &&\text{A93/A95}, \label{out:epsA}\\
 \frac{D\eps}{Dt}&=\nabla\cdot\left(\frac{\nut}{\sigma_\eps}\nabla\eps\right)
       +C_{\eps S}(\eta)\mathcal S\eps-C_{\eps2}\frac{\eps^2}{k},
 &&\text{E94/E95}. \label{out:epsE}
\end{align}
The molecular parameter is not set to zero inside the later inner wall
element. The inner finite-viscosity model and the outer approximation
are different parts of the composite construction. Unless explicitly
identified as an original equation above, all balances in the remaining
three sections concern the outer system \cref{out:momentum,out:kouter,out:epsA,out:epsE}.

\section{The neutral reference and its complete Reynolds covariance}
\label{out:neutral}

\subsection{Exact simple-shear reductions}
\label{out:simpleshear}

Let $\boldsymbol U=U(z)\boldsymbol e_x$ and $G=U'(z)>0$. Then
\begin{equation}
 \Amat=G\begin{pmatrix}0&0&1\\0&0&0\\0&0&0\end{pmatrix},\qquad
 \Smat=\frac G2\begin{pmatrix}0&0&1\\0&0&0\\1&0&0\end{pmatrix},\qquad
 \Omat=\frac G2\begin{pmatrix}0&0&1\\0&0&0\\-1&0&0\end{pmatrix}.
 \label{out:shearmatrices}
\end{equation}
In particular,
\begin{equation}
 s=r=\frac G{\sqrt2},\qquad
 \mathcal S=\mathcal U=G,\qquad
 W=0,\qquad A_s=\frac3{\sqrt2},\qquad
 q=G\tau=\frac{Gk}{\eps}.
 \label{out:shearinvariants}
\end{equation}
Both A93 invariants equal $q$. To distinguish the simple-shear scalar
functions from the full tensor-dependent coefficient, write
\begin{equation}
 C_{93}(q)=\frac{2/3}{1.25+1.9q},\qquad
 C_{95}(q)=\frac1{6.5+(3/\sqrt2)q},\qquad
 C_E(q)=\frac1{4+(3/\sqrt2)q}.
 \label{out:Cq}
\end{equation}
These functions are evaluations of the original laws, not new fits.

The matrix products needed in A93 are
\begin{equation}
 \Amat^2=0,\qquad
 \Amat\Amat^T=G^2\diag(1,0,0),\qquad
 \Amat^T\Amat=G^2\diag(0,0,1).
 \label{out:shearproducts}
\end{equation}
Thus
\begin{equation}
 \boxed{\Q^{93}=
 \frac{kq^2}{1000+q^3}\diag\left(\frac{28}3,-\frac{11}3,-\frac{17}3\right).}
 \label{out:A93shearQ}
\end{equation}
For A95,
\begin{equation}
 \Omat\Smat-\Smat\Omat=\frac{G^2}{2}\diag(1,0,-1),
\end{equation}
and hence
\begin{equation}
 \boxed{\Q^{95}=C_Qkq^2\diag(1,0,-1),\qquad
 C_Q(q)=\frac{\sqrt{1-\frac92 C_{95}(q)^2q^2}}{1+3q^2}.}
 \label{out:A95shearQ}
\end{equation}
The E94/E95 correction remains zero. Therefore all three complete
simple-shear covariances are
\begin{equation}
 \R=
 \begin{pmatrix}
  \tfrac23k+Q_{xx}&0&-\nut G\\
  0&\tfrac23k+Q_{yy}&0\\
  -\nut G&0&\tfrac23k+Q_{zz}
 \end{pmatrix},\qquad
 \boxed{Q_{xz}=0,\quad R_{xz}=-\nut G,\quad P_k=\nut G^2.}
 \label{out:shear}
\end{equation}
In particular, the full shear covariance is nonzero. The normal
corrections in \cref{out:A93shearQ,out:A95shearQ} follow directly from
the original closures; they require no change of their constitutive
coefficients. When comparing with the papers' $U_1(x_2)$ calculations,
the coordinate mapping here is $x_1=x$, $x_2=z$, $x_3=y$.

\subsection{Constant stress, constant kinetic energy, and the three roots}
\label{out:rootssection}

Choose $\ust>0$, $\kappa>0$, and aerodynamic roughness length $z_0>0$.
On the outer interval $H\le z\le Z$, with $Z>H>z_0$, prescribe
\begin{equation}
 U_{o,0}(z)=\frac{\ust}{\kappa}\ln\frac z{z_0},\qquad
 G_0(z)=\frac{\ust}{\kappa z}.
 \label{out:log}
\end{equation}
The suffix $o$ identifies an outer field and the suffix $0$ evaluation
on the selected reference. Neither suffix means evaluation at the
numerical ground $z=0$.

Stationarity, horizontal homogeneity, and zero streamwise pressure
gradient reduce the streamwise momentum equation to
\begin{equation}
 \frac{\dd}{\dd z}(\nu_{t,o,0}G_0)=0.
\end{equation}
Define the friction velocity by this constant turbulent stress,
\begin{equation}
 -R_{xz,0}=\nu_{t,o,0}G_0=\ust^2.
 \label{out:constantstress}
\end{equation}
The explicit canonical ansatz $k_0=K_0$, constant in $z$, makes its
outer diffusion zero. The kinetic-energy equation then requires
$P_{k,0}=\eps_{o,0}$. Combining these balances with \cref{out:log}
gives
\begin{equation}
 \eps_{o,0}(z)=\frac{\ust^3}{\kappa z},\qquad
 \nu_{t,o,0}(z)=\kappa\ust z,\qquad
 q=\frac{K_0}{\ust^2}.
 \label{out:balances}
\end{equation}
Finally, evaluating the original eddy viscosity supplies
\begin{equation}
 \boxed{C_m(q_m)q_m^2=1,\qquad
        K_0=q_m\ust^2,\qquad C_{\mu,*,m}=q_m^{-2},}
 \label{out:rootcondition}
\end{equation}
where $m\in\{93,95,E\}$. There is one positive root for each model:
\begin{align}
 \frac23q_{93}^2&=1.25+1.9q_{93},
 &q_{93}&=\frac{57+\sqrt{6249}}{40},\nonumber\\
 q_{95}^2&=6.5+\frac3{\sqrt2}q_{95},
 &q_{95}&=\frac{3+\sqrt{61}}{2\sqrt2},\nonumber\\
 q_E^2&=4+\frac3{\sqrt2}q_E,
 &q_E&=\frac{3+\sqrt{41}}{2\sqrt2}.
 \label{out:roots}
\end{align}
Thus the common form of the canonical outer reference is
\begin{equation}
 \boxed{
 U_{o,0}=\frac{\ust}{\kappa}\ln\frac z{z_0},\quad
 k_0=q_m\ust^2,\quad
 \eps_{o,0}=\frac{\ust^3}{\kappa z},\quad
 \nu_{t,o,0}=\kappa\ust z,\quad
 R_{xz,0}=-\ust^2.}
 \label{out:reference}
\end{equation}
The dissipation equation has not yet been imposed: its additional
compatibility condition is essential.

\begin{table}[htbp]
\centering
\caption{Constitutive constants on the canonical simple-shear branch.
The exact expressions in \cref{out:roots} define the values; displayed
numbers are rounded. These constants do not depend on the prescribed
$\kappa$.}
\label{out:roottable}
\begin{tabular}{lrrr}
\toprule
Model & $q_m=k_0/\ust^2$ & $C_{\mu,*,m}$ & $C_{\mu,*,m}^{1/4}$\\
\midrule
A93 & 3.4012654174 & 0.0864408350 & 0.5422252508\\
A95 & 3.8220004261 & 0.0684571063 & 0.5115105958\\
E94/E95 & 3.3245064563 & 0.0904785514 & 0.5484492029\\
\bottomrule
\end{tabular}
\end{table}

\subsection{Full normal stresses and reference realizability}
\label{out:covariance}

Define the dimensionless coefficients
\begin{equation}
 F_{93}=\frac{q_{93}^2}{1000+q_{93}^3},\qquad
 C_{Q,*}=\frac{\sqrt{1-\frac{9}{2q_{95}^2}}}{1+3q_{95}^2},\qquad
 B_{95}=C_{Q,*}q_{95}^2.
 \label{out:normalcoeffs}
\end{equation}
The square-root radicand is
$1-9/(2q_{95}^2)=0.6919430214\ldots>0$.
The complete reference covariance is
\begin{equation}
 \boxed{
 \R_0=k_0
 \begin{pmatrix}
 a_x&0&-1/q_m\\
 0&a_y&0\\
 -1/q_m&0&a_z
 \end{pmatrix},}
 \label{out:R0}
\end{equation}
where
\begin{equation}
 (a_x,a_y,a_z)=
 \begin{cases}
 (\frac23+\frac{28}3F_{93},\,
  \frac23-\frac{11}3F_{93},\,
  \frac23-\frac{17}3F_{93}),&m=93,\\[3pt]
 (\frac23+B_{95},\,\frac23,\,\frac23-B_{95}),&m=95,\\[3pt]
 (\frac23,\,\frac23,\,\frac23),&m=E.
 \end{cases}
 \label{out:normalcomponents}
\end{equation}
For all cases $a_x+a_y+a_z=2$, as required by the definition of $k$.
The three eigenvalues of $\R_0/k_0$ are
\begin{equation}
 a_y,\qquad
 \frac{a_x+a_z}{2}\pm
 \frac12\sqrt{(a_x-a_z)^2+\frac4{q_m^2}}.
 \label{out:referenceeigenvalues}
\end{equation}
They are all strictly positive, as quantified in
\cref{out:normalstable}. This establishes positive definiteness of the
\emph{reference covariance}. It is separate from realizability for
arbitrary gradients and from positivity of the incompressible momentum
principal symbol, which require additional analysis.

\begin{table}[htbp]
\centering\small
\caption{Full neutral normal stresses and the smallest covariance
eigenvalue, normalized by $\ust^2$. In every row
$R_{xz,0}/\ust^2=-1$.}
\label{out:normalstable}
\begin{tabular}{lrrrr}
\toprule
Model & $R_{xx,0}/\ust^2$ & $R_{yy,0}/\ust^2$
& $R_{zz,0}/\ust^2$ & $\lambda_{\min}(\R_0)/\ust^2$\\
\midrule
A93 & 2.6208540198 & 2.1286966655 & 2.0529801495 & 1.2973882572\\
A95 & 3.5841102266 & 2.5480002841 & 1.5118903415 & 1.1080267381\\
E94/E95 & 2.2163376375 & 2.2163376375 & 2.2163376375 & 1.2163376375\\
\bottomrule
\end{tabular}
\end{table}

\subsection{Normal momentum and the pressure variable}
\label{out:pressure}

The normal-stress correction must also be included in the vertical
momentum equation. With zero vertical velocity, it reduces to
\begin{equation}
 \frac{\dd}{\dd z}(\pi_0+R_{zz,0})=0,
 \qquad
 \boxed{\pi_0+R_{zz,0}=\Pi_0,}
 \label{out:pressurebalance}
\end{equation}
where $\Pi_0$ is a constant fixed by the pressure datum. All three
canonical normal stresses are constant, so a constant hydrostatically
reduced $\pi_0$ satisfies this condition. If the isotropic turbulent
stress is absorbed into the solved pressure,
\begin{equation}
 \pi^\dagger=\pi+\frac23k,\qquad
 \boxed{\pi_0^\dagger+Q_{zz,0}=\Pi_0,}
 \label{out:modifiedpressure}
\end{equation}
since $S_{zz,0}=0$. The two pressure conventions are equivalent when
the corresponding normal-stress term is retained. One pressure datum
is sufficient; changing the pressure variable does not permit dropping
$Q_{zz}$ from the normal momentum balance.

\section{Dissipation compatibility and the choice of logarithmic slope}
\label{out:compatibility}

\subsection{The additional condition supplied by the dissipation equation}
\label{out:epscompat}

For \cref{out:reference}, direct differentiation gives
\begin{equation}
 \frac{\dd}{\dd z}
  \left(\frac{\nu_{t,o,0}}{\sigma_\eps}\eps_{o,0}'\right)
 =\frac{\ust^4}{\sigma_\eps z^2},\qquad
 \frac{\eps_{o,0}^2}{k_0}
 =\frac{\ust^4}{\kappa^2q_mz^2}.
 \label{out:epsbalances}
\end{equation}
For A93/A95, $P_{k,0}=\eps_{o,0}$ combines the reaction terms into
$-(C_{\eps2}-C_{\eps1})\eps_{o,0}^2/k_0$. The compatibility condition
is therefore
\begin{equation}
 \kappa^2q_m=\sigma_{\eps,m}(C_{\eps2}-C_{\eps1}),
 \qquad m=93,95.
 \label{out:compatA}
\end{equation}
For E94/E95,
\begin{equation}
 \frac{q_E}{q_E+5}=0.39936379097\ldots<0.43,
 \qquad C_{\eps S,*}=0.43.
 \label{out:Ebranch}
\end{equation}
Moreover $\mathcal S_0\eps_{o,0}
=q_E\eps_{o,0}^2/k_0$. Hence its condition is
\begin{equation}
 \kappa^2q_E=\sigma_{\eps,E}(C_{\eps2}-0.43q_E).
 \label{out:compatE}
\end{equation}
With all original reaction coefficients retained, introduce the
positive net dissipation-sink constants
\begin{equation}
 D_m=\begin{cases}
 1.92-1.44=0.48,&m=93,95,\\
 1.9-0.43q_E=0.4704622237849\ldots,&m=E.
 \end{cases}
 \label{out:Dm}
\end{equation}
Here $D_m$ is a dimensionless model constant, not the position-dependent
sum of molecular and turbulent viscosity used in the inner wall
momentum equation. The single outer compatibility identity is
\begin{equation}
 \boxed{\kappa^2q_m=\sigma_{\eps,m}D_m.}
 \label{out:compat}
\end{equation}
Together with the constitutive root and the pressure condition, it
makes \cref{out:reference} a solution of every selected outer equation.

\subsection{The natural logarithmic slopes of the unmodified models}
\label{out:naturalkappa}

When the published reaction and diffusion constants are used, the
canonical calculation predicts
\begin{equation}
 \kappa_m^{\rm nat}
 =\sqrt{\frac{\sigma_{\eps,m,\rm org}D_m}{q_m}}.
 \label{out:kappanatural}
\end{equation}
The resulting values are
\begin{table}[htbp]
\centering
\caption{Canonical logarithmic slopes implied by the original
coefficients in the selected outer high-Re equations.}
\label{out:kappatable}
\begin{tabular}{lrrr}
\toprule
Model & $\sigma_{\eps,m,\rm org}$ & $D_m$ & $\kappa_m^{\rm nat}$\\
\midrule
A93 & 1.3 & 0.4800000000 & 0.4283236287\\
A95 & 1.3 & 0.4800000000 & 0.4040609953\\
E94/E95 & 1.2 & 0.4704622238 & 0.4120874716\\
\bottomrule
\end{tabular}
\end{table}
No coefficient change is needed if a model's own
$\kappa_m^{\rm nat}$ is accepted. The E papers state a log-layer
calibration using $\kappa=0.41$ and the rounded constants
$C_\mu\simeq0.09$, $\sigma_\eps=1.20$
\citep[Eq.~(28) and Table~2]{Shih1995EVM}. Inserting the printed
coefficients and the exact positive root into the canonical equations
instead gives the value in \cref{out:kappatable}. This is the exact
implication of those coefficients for the stated reference, not a claim
that every boundary-layer solution must have this constant-$k$ form.

The distinction between a model's natural $\kappa$ and a prescribed
common $\kappa$ is operational. Choosing a different $z_0$ changes the
velocity intercept, and changing $\ust$ changes the dimensional
amplitudes. Neither alters the dimensionless identity
\cref{out:compat}. A wall boundary condition likewise cannot make a
nonzero dissipation imbalance vanish throughout the outer interval.

\subsection{A common prescribed \texorpdfstring{$\kappa$}{kappa}: one diffusivity change}
\label{out:minimalchange}

If the reference must have one independently prescribed
$\kappa>0$, retain the original stress laws, all reaction coefficients,
and $\sigma_k=1$, and replace only
\begin{equation}
 \boxed{\sigma_{\eps,m}^{\rm eq}(\kappa)
       =\frac{\kappa^2q_m}{D_m}.}
 \label{out:sigma}
\end{equation}
All these values are positive. They are coefficients of the whole
outer dissipation diffusion operator, not values to be changed only
at a boundary. The stress root, $C_{\mu,*,m}$, covariance anisotropy,
and the frozen-$k,\eps$ constitutive tangent remain those of the
original model.

\begin{table}[htbp]
\centering
\caption{Required dissipation diffusivity parameter for two prescribed
common logarithmic slopes. These are modified values; the original
values are shown separately.}
\label{out:sigmatable}
\begin{tabular}{lrrr}
\toprule
Model & $\sigma_{\eps,m,\rm org}$
& $\sigma_{\eps,m}^{\rm eq}(0.40)$
& $\sigma_{\eps,m}^{\rm eq}(0.41)$\\
\midrule
A93 & 1.3 & 1.1337551391 & 1.1911514931\\
A95 & 1.3 & 1.2740001420 & 1.3384963992\\
E94/E95 & 1.2 & 1.1306349503 & 1.1878733447\\
\bottomrule
\end{tabular}
\end{table}

This is minimal in the number of changed scalar constants within the
fixed original stress law, fixed canonical family, and fixed reaction
coefficients: zero changes suffice at $\kappa_m^{\rm nat}$, while one
change suffices at any prescribed positive $\kappa$. It is not a
uniqueness claim under all possible measures of model change. Changing
$\sigma_\eps$ changes inhomogeneous dissipation transport, while leaving
the diffusion-free reaction calibration intact. An alternative change
of one destruction coefficient is given separately in the appendix;
it is not applied simultaneously with \cref{out:sigma}.

The subsequent composite construction uses an explicitly selected
$\kappa$ and the compatible coefficients. Retaining each natural value
is also possible, but then different models have different target
logarithmic slopes at a common prescribed $\ust$.

\subsection{What a failed compatibility check does and does not exclude}
\label{out:compatibilityscope}

Condition \cref{out:compat} is necessary and sufficient for the
\emph{canonical constant-$k$ branch} under the stated outer assumptions.
Its failure does not, by itself, exclude every solution with the same
logarithmic mean velocity and nonconstant $k(z)$ and $\eps(z)$.
For example, suppose the logarithm and constant shear stress are retained
but $q=q(z)$ is allowed to vary. The constitutive shear relation would
then give
\begin{equation}
 k(z)=\frac{\ust^2}{C_m(q)q},\qquad
 \eps(z)=\frac{\ust^3}{\kappa z\,C_m(q)q^2},\qquad
 \nut(z)=\kappa\ust z.
 \label{out:nonconstantcandidate}
\end{equation}
Such fields still have to satisfy both turbulence equations, including
\begin{equation}
 \frac{\dd}{\dd z}\left(\frac{\kappa\ust z}{\sigma_k}k'(z)\right)
 +\frac{\ust^3}{\kappa z}-\eps(z)=0,
 \label{out:nonconstantk}
\end{equation}
as well as the pressure balance and boundary data. This note neither
constructs nor rules out those more general branches. The coefficient
change in \cref{out:sigma} has the specific purpose of making the
standard constant-$k$ neutral reference exact.

\section{Neutral MOST, matching-plane wall data, and integral checks}
\label{out:wallbenchmark}

\subsection{What is matched to the neutral similarity reference}
\label{out:MOST}

The importance of compatible inlet profiles, turbulence balances,
and rough-wall data for homogeneous atmospheric-flow calculations
is discussed by \citet{RichardsHoxey1993,RichardsNorris2011,Blocken2007}.
For the present purpose, the neutral Monin--Obukhov similarity
reference is specified by the mean gradient
\begin{equation}
 \frac{\dd U_M}{\dd z}=\frac{\ust}{\kappa_M z},\qquad
 U_M(z)=\frac{\ust}{\kappa_M}\ln\frac z{z_0}.
 \label{out:MOSTreference}
\end{equation}
Together with local production--dissipation balance, it supplies
$\eps_M=\ust^3/(\kappa_Mz)$. The logarithmic mean relation alone does
not prescribe $k/\ust^2$ or the three normal-stress ratios; the
constitutive model supplies those quantities through
\cref{out:roots,out:normalcomponents}.

For the same $\ust$ and $z_0$, comparing an unmodified model at its
natural slope with this prescribed target gives
\begin{equation}
 \frac{U_{o,0,m}}{U_M}
 =\frac{G_{0,m}}{U_M'}
 =\frac{\eps_{o,0,m}}{\eps_M}
 =\frac{\kappa_M}{\kappa_m^{\rm nat}},\qquad
 \frac{\nu_{t,o,0,m}}{\kappa_M\ust z}
 =\frac{\kappa_m^{\rm nat}}{\kappa_M},
 \label{out:MOSTratios}
\end{equation}
where the velocity ratio is taken at $z>z_0$.
After selecting $\kappa=\kappa_M$ and applying
\cref{out:sigma}, all models have the same $U_{o,0}$,
$\eps_{o,0}$, and $\nu_{t,o,0}$ at the same $\ust,z_0$.
Their $k_0$ and normal stresses still differ, as required by their
original closures.

A different comparison fixes $U_M(z_r)=U_r$ at some $z_r>z_0$ and
infers the friction velocity for each chosen $\kappa$:
\begin{equation}
 \ust=\frac{\kappa U_r}{\ln(z_r/z_0)},\qquad
 U_{o,0}(z)=U_r\frac{\ln(z/z_0)}{\ln(z_r/z_0)}.
 \label{out:velocitynormalization}
\end{equation}
The velocity curves then coincide even when the inferred friction
velocities, shear stresses, kinetic energies, and dissipation rates
differ. A comparison of velocity profiles alone therefore does not
establish equality of the complete neutral reference states.

\subsection{An elementary equilibrium-admitting matching-plane wall law}
\label{out:walllaw}

This subsection records a benchmark for the \emph{resolved outer
boundary} at $z=H$. It is not a finite-viscosity reconstruction between
the ground and $H$; that construction is developed later. In particular,
the logarithm is used only above the matching plane, and $z_0$ is an
aerodynamic roughness length rather than the location of a material
no-slip boundary.

Let $L_H=\ln(H/z_0)>0$, and let
$\boldsymbol U_{\parallel,H}$ be the tangential velocity trace.
An elementary velocity-based friction estimate and corresponding
kinematic traction on the lower outer boundary are
\begin{align}
 u_{\tau,U}&=\frac{\kappa|\boldsymbol U_{\parallel,H}|}{L_H},
 \label{out:utauU}\\
 \boxed{\boldsymbol t_{\parallel,H}^{o}
   =-u_{\tau,U}^2
       \frac{\boldsymbol U_{\parallel,H}}
            {|\boldsymbol U_{\parallel,H}|}
   =-\frac{\kappa^2}{L_H^2}
       |\boldsymbol U_{\parallel,H}|\boldsymbol U_{\parallel,H}.}
 \label{out:walltraction}
\end{align}
The second expression has the continuous zero value at zero tangential
velocity. Multiplication by $\rho$ gives the dynamic traction.
On \cref{out:reference}, $u_{\tau,U}=\ust$ and
$\boldsymbol t_{\parallel,H}^{o}=-\ust^2\boldsymbol e_x$.
This is a traction boundary law using the solution's velocity trace;
one does not additionally prescribe an independent tangential
Dirichlet value at the same boundary.

The corresponding reference turbulence values satisfy
\begin{equation}
 \boxed{k_H=q_m u_\tau^2,\qquad
 \eps_H=\frac{u_\tau^3}{\kappa H}
        =\frac{k_H^{3/2}}{\kappa q_m^{3/2}H}
        =C_{\mu,*,m}^{3/4}\frac{k_H^{3/2}}{\kappa H}.}
 \label{out:wallvalues}
\end{equation}
A kinetic-energy-based friction estimate is equivalently
\begin{equation}
 u_{\tau,k}=\sqrt{\frac{k_H}{q_m}}
           =C_{\mu,*,m}^{1/4}\sqrt{k_H}.
 \label{out:utauk}
\end{equation}
The velocity- and kinetic-energy-based estimates coincide on the
reference; they are not two independently enforceable equalities for
arbitrary off-reference data. For a standalone outer benchmark one may
use a compatible $k_H$ value or the zero normal $k$-gradient that the
constant reference also satisfies. Those are alternative boundary
choices. The transported-$k$ inner construction later replaces this
choice by its own interface coupling.

It is convenient to name the outer reference length coefficient
\begin{equation}
 c_{\ell,m}^{\rm eq}=\kappa q_m^{3/2},\qquad
 \eps_H=\frac{k_H^{3/2}}{c_{\ell,m}^{\rm eq}H}.
 \label{out:cell}
\end{equation}
Its values are
\begin{table}[htbp]
\centering
\caption{Outer reference length coefficient in the equilibrium
matching value for $\eps$. These numbers are not replacements for
coefficients of the later inner momentum kernel.}
\label{out:celltable}
\begin{tabular}{lrr}
\toprule
Model & $c_{\ell,m}^{\rm eq}$ at $\kappa=0.40$
& $c_{\ell,m}^{\rm eq}$ at $\kappa=0.41$\\
\midrule
A93 & 2.5091162113 & 2.5718441166\\
A95 & 2.9887947248 & 3.0635145930\\
E94/E95 & 2.4246595226 & 2.4852760107\\
\bottomrule
\end{tabular}
\end{table}

\subsection{Boundary orientations, normal stress, and scalar fluxes}
\label{out:boundarysigns}

The outward normal of the resolved outer domain is
$\boldsymbol n_H=-\boldsymbol e_z$ at the matching plane and
$\boldsymbol n_Z=+\boldsymbol e_z$ at its top. Its kinematic effective
Cauchy stress under the selected outer approximation is
\begin{equation}
 \boldsymbol\Sigma_o=-\pi\Id-\R.
 \label{out:outerCauchy}
\end{equation}
Thus the canonical tangential tractions are
\begin{equation}
 (\boldsymbol\Sigma_{o,0}\boldsymbol n_H)_\parallel
   =-\ust^2\boldsymbol e_x,\qquad
 (\boldsymbol\Sigma_{o,0}\boldsymbol n_Z)_\parallel
   =+\ust^2\boldsymbol e_x.
 \label{out:boundarytractions}
\end{equation}
At either horizontal plane the normal traction component is
\begin{equation}
 \boldsymbol n\cdot\boldsymbol\Sigma_{o,0}\boldsymbol n
 =-(\pi_0+R_{zz,0})=-\Pi_0.
 \label{out:normaltraction}
\end{equation}
The identity is useful for verifying the pressure convention. When
normal velocity is imposed, it is not an additional independent
pressure boundary condition.

For either scalar $\phi\in\{k,\eps\}$, define the outward
\emph{natural diffusion datum}
\begin{equation}
 \Phi_\phi=\frac{\nut}{\sigma_\phi}\nabla\phi\cdot\boldsymbol n.
 \label{out:fluxdefinition}
\end{equation}
This is the boundary term obtained by integrating the positive
diffusion divergence in the scalar equation. The physical outward
Fickian flux is $-\Phi_\phi$, with the opposite sign. At the neutral
reference,
\begin{equation}
 \Phi_{k,H}=\Phi_{k,Z}=0,\qquad
 \boxed{\Phi_{\eps,H}=+\frac{\ust^4}{\sigma_{\eps,m}H},\qquad
        \Phi_{\eps,Z}=-\frac{\ust^4}{\sigma_{\eps,m}Z}.}
 \label{out:referencefluxes}
\end{equation}
The actual coefficient used in the outer PDE must also be used in these
fluxes; for a prescribed common $\kappa$ this is
$\sigma_{\eps,m}^{\rm eq}$.

Analytic reference values and the corresponding natural diffusion
fluxes are alternative ways of admitting the same scalar reference;
they are not simultaneously prescribed independent data. In the later
wall coupling, the dependent inner dissipation supplies the matching
value, while \cref{out:referencefluxes} remains a reference check.
A homogeneous dissipation Neumann condition at a finite upper height
would exclude this reference. Similarly, a stress-free top or a zero
mean-velocity gradient there is incompatible with nonzero constant
shear stress. One may instead impose analytic top values or compatible
traction and scalar-flux data, subject to a consistent mixed boundary
formulation.

For a streamwise through-flow benchmark, the inlet profiles must include
all of $U_{o,0}$, $k_0$, and $\eps_{o,0}$. Horizontal homogeneity admits
zero streamwise derivatives at an appropriate outflow and periodic
spanwise data, with one compatible pressure datum. The imposed data
must refer to the same $\ust,z_0,\kappa$, model root, and chosen
transport coefficients.

\subsection{Slab-integrated checks and the limit of the continuum result}
\label{out:integrals}

Consider an outer slab $[a,b]$, with $H\le a<b\le Z$, per unit
horizontal area. Production and dissipation coincide pointwise, and
therefore
\begin{equation}
 \int_a^b P_{k,0}\,\dd z
 =\int_a^b\eps_{o,0}\,\dd z
 =\frac{\ust^3}{\kappa}\ln\frac ba,
 \qquad
 \overline{P}_{k,[a,b]}=\overline\eps_{[a,b]}
 =\frac{\ust^3}{\kappa(b-a)}\ln\frac ba.
 \label{out:productionintegral}
\end{equation}
With no mean normal velocity, the tangential stress work on the two
planes is
\begin{equation}
 \ust^2[U_{o,0}(b)-U_{o,0}(a)]
 =\frac{\ust^3}{\kappa}\ln\frac ba
 =\int_a^bP_{k,0}\,\dd z.
 \label{out:outerpower}
\end{equation}
This equality concerns the molecular-viscosity-free outer equations.
In a finite-viscosity inner element, total mechanical viscous work
must additionally account for molecular mean-flow dissipation.

The integrated dissipation diffusion is
\begin{equation}
 \left[\frac{\nu_{t,o,0}}{\sigma_{\eps,m}}
           \eps_{o,0}'\right]_a^b
 =\frac{\ust^4}{\sigma_{\eps,m}}
       \left(\frac1a-\frac1b\right),
 \label{out:epsdiffintegral}
\end{equation}
whereas its integrated reaction term is
\begin{equation}
 -\int_a^bD_m\frac{\eps_{o,0}^2}{k_0}\,\dd z
 =-\frac{D_m\ust^4}{\kappa^2q_m}
       \left(\frac1a-\frac1b\right).
 \label{out:epsreactionintegral}
\end{equation}
They cancel exactly under \cref{out:compat}. The momentum tractions
likewise sum to zero and the $k$ diffusion flux is zero. These checks
verify signs, dimensions, and the use of the same transport coefficients
in the PDE and its boundary terms.

The logarithm and inverse-height dissipation are generally not
represented exactly by a finite polynomial approximation. Point values
also differ from their cell averages in
\cref{out:productionintegral}. Consequently, an exact continuum
reference and compatible continuum boundary data do not alone establish
exact discrete preservation, positivity, or convergence of a numerical
scheme. Those are separate properties of the discretization and wall
coupling.

\section{Original finite-viscosity analytical wall function}
\label{awf:original}

This section specifies the scalar SCI baseline \citep{SCI2006} in
physical coordinates. Its original coefficients and roughness
correlation are retained in the principal formulation. A vector
version applies the same scalar resistance to both physical tangential
velocity components. The later full-strain, curved, and unsteady
completion is separately declared; the scalar column equations alone
are not a complete terrain-flow model.

\subsection{Geometry, sampled energy, and the local reference family}
\label{awf:geometry}

Let $0\leq z\leq H$ be the modeled wall layer and $H<z\leq Z$ the
resolved outer layer. The numerical ground is $z=0$. The matching
height $H$, outer top $Z$, equivalent sand-grain roughness height
$h_r$, and corrected-tail start $z_m$ are physical inputs satisfying
\begin{equation}
 0\leq h_r<H<Z,\qquad 0<z_m\leq H,\qquad 0<z_0<H.
 \label{awf:heights}
\end{equation}
The aerodynamic roughness length $z_0$ is distinct from $h_r$.
The original AWF uses $h_r$ both in its viscosity-intercept correlation
and as the cutoff below which sampled momentum forcing is suppressed.
An independently prescribed $z_0$ must satisfy the resistance
compatibility condition derived below; it is not an additional freely
imposable wall input once all other parameters have been fixed.

Let $K=\mathcal K[k_h]>0$ be a declared, constant-preserving sampling
functional, for example evaluation at a fixed physical height
$z_K\in(0,H]$ or an explicitly specified positive normalized average.
The functional and all physical heights are held fixed under
polynomial-degree refinement. The symbol $K$ is a current wall input;
it is not the transported inner field $k_h(z)$ at every height.
For each model define
\begin{equation}
 \boxed{
 u_m(K)=\sqrt{K/q_m},\qquad
 t_m(K)=u_m^2=K/q_m,\qquad
 A_m(K)=\frac{u_m^3}{\kappa}
       =\frac{K^{3/2}}{\kappa q_m^{3/2}}.
 }
 \label{awf:reference-scales}
\end{equation}
These scales describe a counterfactual zero-forcing reference at the
same current $K$. On the fixed inlet reference,
$K=K_0=q_m\ust^2$ and $u_m(K_0)=\ust$. In a forced or reversed-flow
state, $u_m(K)$ need not equal the friction velocity inferred from the
actual signed wall traction.

\subsection{The retained finite-viscosity SCI kernel}
\label{awf:kernel}

Define
\begin{equation}
 \boxed{
 z^*=\frac{z\sqrt K}{\nu},\qquad
 \nu_{t,A}(z)=\nu\max\{0,\alpha(z^*-b)\},\qquad
 D_A(z)=\nu+\nu_{t,A}(z)>0.
 }
 \label{awf:selected-viscosity}
\end{equation}
The original constants are
\begin{equation}
 c_\mu^w=0.09,\quad c_\ell^w=2.55,\quad
 \alpha=c_\mu^wc_\ell^w=0.2295,\quad z_{vs}^*=10.7.
 \label{awf:original-constants}
\end{equation}
These are auxiliary wall-model constants, not the variable outer
$C_\mu$. In particular the velocity scale in $z^*$ is $\sqrt K$,
not $u_m$.
Writing $h^*=h_r\sqrt K/\nu$ and $x=h^*/70$, the rough-wall
correlation is
\begin{equation}
 b=z_{vs}^*(1-x^{m_r}),\qquad
 m_r=\max\{0.5-0.4x^{0.7},\,1-0.79x^{-0.28}\},\qquad h_r>0.
 \label{awf:original-intercept}
\end{equation}
For $h_r=0$, set $b=10.7$ directly without evaluating the negative
power at $x=0$. Equivalently, set
\begin{equation}
 a_A=\alpha\sqrt K,\qquad z_v=\frac{\nu b}{\sqrt K},\qquad
 \nu_{t,A}=a_A(z-z_v)_+,
 \label{awf:dimensional-kernel}
\end{equation}
where $(q)_+=\max(q,0)$. If $z_v>0$, the interval below $z_v$
is molecular-only in this auxiliary momentum kernel. If $z_v<0$,
the eddy viscosity is positive at the numerical ground; $z_v$ is
then a virtual intercept, not a displaced ground or a negative
physical sublayer thickness.

\subsection{No slip, signed shear, and the correct interface traction}
\label{awf:momentum}

Let $\boldsymbol U_H$ be the tangential outer velocity trace and
$\boldsymbol F$ the sampled tangential pressure/convection forcing,
with the sign convention
\begin{equation}
 \begin{aligned}
 \frac{\dd}{\dd z}\left(D_A\boldsymbol U'_{A,\parallel}\right)
 &=\boldsymbol F\,\mathbf1_{z>h_r},\\
 \boldsymbol U_{A,\parallel}(0)&=0,\qquad
 \boldsymbol U_{A,\parallel}(H)=\boldsymbol U_H.
 \end{aligned}
 \label{awf:momentum-ode}
\end{equation}
Hydrostatic pressure and gravity must cancel in the forcing sampler.
For example, a two-dimensional flat reduction may use the declared
one-sided outer sampler
$F_x=\mathcal E_H^+[\partial_x(U_o^2)+\rho^{-1}\partial_xp_o-b_x]$,
where $b_x$ is the streamwise body acceleration. This retains the
original AWF's simplified conservative convection approximation;
it is not the complete unsteady inner momentum equation. The
curved patch formulation below retains storage and transport.
Define $L_r(z)=(z-h_r)_+$ and
\begin{equation}
 \begin{aligned}
 J_0(z)&=\int_0^z\frac{\dd\xi}{D_A(\xi)},&
 J_1(z)&=\int_0^z\frac{L_r(\xi)}{D_A(\xi)}\,\dd\xi,\\
 J_2(z)&=\int_0^z\frac{\dd\xi}{D_A(\xi)^2},&
 B&=\frac{J_1(H)}{J_0(H)}.
 \end{aligned}
 \label{awf:integrals}
\end{equation}
All three integrals are finite because $D_A\geq\nu>0$.
Closed forms, including all molecular-only and rough-active cases,
are derived in \cref{app:integrals}.

Integration of the momentum equation gives
\begin{align}
 \boldsymbol t_A&=\frac{\boldsymbol U_H}{J_0(H)}
                  -B\boldsymbol F, \label{awf:wall-shear}\\
 \boldsymbol U_{A,\parallel}(z)
 &=\frac{\boldsymbol U_H}{J_0(H)}J_0(z)
   +\boldsymbol F[J_1(z)-BJ_0(z)], \label{awf:velocity}\\
 D_A\boldsymbol U'_{A,\parallel}
 &=\boldsymbol t_A+\boldsymbol F L_r(z). \label{awf:shear-profile}
\end{align}
Reversed shear remains signed. The ground force on the fluid is
$-\rho\boldsymbol t_A$, whereas the upper interface carries
\begin{equation}
 \boxed{\boldsymbol T_H
 =\boldsymbol t_A+\boldsymbol F(H-h_r)
 =\frac{\boldsymbol U_H}{J_0(H)}
  +[(H-h_r)-B]\boldsymbol F.}
 \label{awf:interface-shear}
\end{equation}
The outer domain's lower outward normal is $-\boldsymbol e_z$;
its tangential traction is therefore $-\rho\boldsymbol T_H$.
Using the ground shear instead omits the integrated inner forcing.
Do not add that forcing again to an outer volume element.

For frozen $K,\boldsymbol F,h_r,H,\nu$,
$\partial\boldsymbol T_H/\partial\boldsymbol U_H
=J_0(H)^{-1}I_\parallel$ is positive. These frozen derivatives do
not constitute the coupled wall/pressure Jacobian or a stability proof.
The equilibrium dissipation/source correction leaves the original
SCI momentum and production response unchanged at fixed inputs.
The reconstructed field, rather than an unrelated polynomial
interpolant, defines the numerical inner velocity and its derivatives.
Adding an independent wall traction to an unchanged strongly imposed
no-slip Galerkin boundary would give zero work against its vanishing
velocity tests and would not implement this wall element.

\subsection{Original finite-viscosity dissipation baseline}
\label{awf:baseline}

Retain the original finite-viscosity dissipation profile
\begin{equation}
 \boxed{
 z_{\eps,A}=\frac{2c_\ell^w\nu}{\sqrt K},\qquad
 E_A(z;K)=\frac{K^{3/2}}{c_\ell^w\max(z,z_{\eps,A})}.
 }
 \label{awf:baseline-epsilon}
\end{equation}
Its constant inner branch is equivalently $2\nu K/z_{\eps,A}^2$.
This is the original continuous plateau/tail choice; no new
fully rough cutoff is introduced. The dissipation cutoff
$z_{\eps,A}$, viscosity intercept $z_v$, roughness height $h_r$,
and corrected cutoff $z_c$ have different definitions. The baseline
is finite even if $z_{\eps,A}>H$.


\section{Equilibrium-preserving construction and wall-coefficient consistency}
\label{awf:construction}

The outer model index enters through $q_m$ and the chosen $\kappa$.
The construction first matches momentum, then matches a positive
dissipation continuation to the actual finite-viscosity turbulent
production budget, and finally supplies the source compatibility
needed for all retained kinetic-energy modes.

\subsection{Reference anchoring and prescribed aerodynamic roughness}
\label{awf:anchor}

For every current $K>0$, use the original-kernel anchor
\begin{equation}
 \boxed{U_{H,0}(K)=t_m(K)J_0(H;K,h_r,\nu).}
 \label{awf:anchor-identity}
\end{equation}
At zero forcing, the continuous composite reference is
\begin{equation}
 U_{\mathrm{comp},0}(z;K)=
 \begin{cases}
 t_mJ_0(z),&0\leq z\leq H,\\[1mm]
 t_mJ_0(H)+\dfrac{u_m}{\kappa}\ln(z/H),&H\leq z\leq Z .
 \end{cases}
 \label{awf:composite-velocity}
\end{equation}
It satisfies numerical no slip at $z=0$. The outer branch can be
written $u_m\ln(z/z_{0,m}(K))/\kappa$, with
\begin{equation}
 \boxed{z_{0,m}(K)=H\exp[-\kappa u_m(K)J_0(H;K,h_r,\nu)].}
 \label{awf:generated-roughness}
\end{equation}
The counterfactual family therefore generally has a $K$-dependent
aerodynamic intercept even though $h_r$ is fixed. The fixed inlet
reference uses only $K_0=q_m\ust^2$.

To retain an independently prescribed $z_0$ at that reference,
require
\begin{equation}
 \boxed{
 J_0(H;K_0,h_r,\nu)
 =J_{\rm tar}:=\frac{\ln(H/z_0)}{\kappa\ust}.
 }
 \label{awf:roughness-compatibility}
\end{equation}
This condition can be imposed by a one-time calibration of the
original input $h_r$, retaining the SCI correlation itself. Keep
the resulting $h_r$ fixed as $K$ subsequently evolves. If $h_r$ is
already fixed independently, evaluate this condition; the original
wall law and an arbitrary additional $z_0$ cannot both be imposed
without checking it. A dissipation correction cannot repair a
momentum-resistance mismatch.
The matching height $H$ is part of this calibration. Changing it
requires renewed resistance and cutoff checks; a fixed $h_r$ does
not automatically produce the same $z_0$ at a different interface
height.

The calibration has a precise range. At fixed $K_0,H,\nu$, put
$\mathcal J(h)=J_0(H;K_0,h,\nu)$,
$J_s=\mathcal J(0)$, and
$J_{\lim}=\lim_{h\uparrow H}\mathcal J(h)$.
The SCI intercept $b(h\sqrt{K_0}/\nu)$ is continuous and strictly
decreasing. To see this, the two exponent branches cross once at
$x_\times\in(3/4,1)$, since their difference is strictly decreasing.
The logarithmic derivatives of $x^{m_r}$ on the respective active
branches are
\begin{equation}
 \begin{cases}
 \frac12-0.4x^{0.7}-0.28x^{0.7}\ln x,&x\leq x_\times,\\
 1+x^{-0.28}(0.2212\ln x-0.79),&x\geq x_\times.
 \end{cases}
 \label{awf:roughness-monotonicity}
\end{equation}
Both are positive on their stated intervals. Hence $\mathcal J$
is nonincreasing, and strictly decreasing wherever a nonzero
turbulent branch lies in the layer. A unique calibration with
$0\leq h_r<H$ and positive integrated turbulent production exists
if and only if
\begin{equation}
 \boxed{J_{\lim}<J_{\rm tar}\leq J_s,
 \qquad J_{\rm tar}<H/\nu.}
 \label{awf:roughness-calibration-gate}
\end{equation}
The dissipation gate derived below is a separate additional condition.
If these conditions fail, preservation of that prescribed target by
this unchanged wall kernel is not claimed. Recalibrating $h_r$ or $b$
at every current $K$ to hold the same $z_0$ throughout the entire
counterfactual family would be a further model change, and is not
part of this construction.

For an implementation whose wall input is an interpolation functional
rather than the actual interface trace, apply the anchor to that
same functional and select the additive outer velocity offset
accordingly. A continuum point-value anchor cannot simply be reused
for a different discrete sampler.

The derivatives generally jump at the model interface:
\begin{equation}
 \frac{U'_{A,0}(H)}{U'_{o,0}(H)}
 =\frac{\kappa u_mH}{D_A(H)},\qquad
 \boxed{D_A(H)U'_{A,0}(H)
 =\nu_{t,o,0}(H)U'_{o,0}(H)=t_m.}
 \label{awf:derivative-flux}
\end{equation}
This is continuity of \emph{total} tangential traction. Inner Reynolds
shear alone is
$R_{A,xz}=-t_m\nu_{t,A}/D_A$, whereas outer Reynolds shear is
$R_{o,xz}=-t_m$. Their difference is exactly the inner molecular
shear contribution; Reynolds shear is not separately required to
be continuous at this model interface.

\subsection{The selected \texorpdfstring{$\kappa$}{kappa} and the original wall constants}
\label{kap:selection}

Every occurrence of $\kappa$ in the wall reference denotes the value
selected for the \emph{actual outer equations}. There are two consistent
outer choices: retain the original coefficients and use
$\kappa_m^{\rm nat}$, or prescribe a common $\kappa$ and use the
corresponding $\sigma_{\eps,m}^{\rm eq}$. A wall coefficient must not
use $\kappa_m^{\rm nat}$ after the outer equations have been recalibrated
to a different slope.

The corrected dissipation amplitude gives the identity
\begin{equation}
 \eps_{o,0}(z)=\frac{K^{3/2}}{c_{\ell,m}^{\rm eq}z},
 \qquad
 \boxed{c_{\ell,m}^{\rm eq}=\kappa q_m^{3/2}}.
 \label{kap:ell-equilibrium}
\end{equation}
It depends on both the selected slope and the model's energy level.
The original baseline $E_A$ still uses $c_\ell^w=2.55$ unless its
coefficient is separately changed. The two coefficients are therefore
given different symbols. A corrected profile with coefficient $2.55$
for every model would not, in general, match the outer dissipation.

\begin{table}[htbp]
\centering\small
\caption{Equilibrium-derived dissipation and viscosity-slope coefficients.
The exact formulas, rather than these rounded entries, define the model.}
\label{kap:table}
\begin{tabular}{@{}lrrrr@{}}
\toprule
 & \multicolumn{2}{c}{Natural outer $\kappa_m^{\rm nat}$}
 & \multicolumn{2}{c}{Prescribed outer $\kappa=0.41$}\\
Model & $c_{\ell,m}^{\rm eq}$ & $\kappa/\sqrt{q_m}$
      & $c_{\ell,m}^{\rm eq}$ & $\kappa/\sqrt{q_m}$\\
\midrule
A93 &2.6867844012&0.2322478870&2.5718441166&0.2223123528\\
A95 &3.0191384281&0.2066814804&3.0635145930&0.2097193443\\
E94/E95&2.4979295305&0.2260090453&2.4852760107&0.2248641732\\
\bottomrule
\end{tabular}
\end{table}

\subsection{The logarithmic slope implied by the original auxiliary viscosity}
\label{kap:implicit}

On the active turbulent branch, the finite-viscosity kernel is
\begin{equation}
 D_A(z)=\alpha\sqrt K\,z+\nu(1-\alpha b).
 \label{kap:active-D}
\end{equation}
At zero forcing and reference shear $t_m=u_m^2$,
\begin{equation}
 U'_{A,0}(z)=\frac{u_m^2}{D_A(z)}
 \sim\frac{u_m}{\alpha\sqrt{q_m}\,z}
\end{equation}
when the term proportional to $z$ dominates the molecular/roughness
intercept. The leading auxiliary logarithmic slope is thus
\begin{equation}
 \boxed{\kappa_{A,m}=\alpha\sqrt{q_m}}.
 \label{kap:auxiliary-kappa}
\end{equation}
For the original $\alpha=0.2295$, its values are
$0.4232558326$, $0.4486710576$, and $0.4184526092$
for A93, A95, and E94/E95. They generally differ from both a prescribed
$0.41$ and the respective natural outer slopes. In particular, the
A95 difference is appreciable.
Even if one used the conventional relation
$q=(0.09)^{-1/2}$, the original constants would imply
$2.55(0.09)^{3/4}=0.4190077565$, rather than exactly $0.41$.
These statements follow from the wall momentum equation; they are
not additional constants asserted by the source papers.

This mismatch does not invalidate the continuous composite matching
in \cref{awf:composite-velocity,awf:derivative-flux}. That construction
matches velocity and total traction at a selected $H$, allowing a
velocity-gradient jump across the model interface.
It does mean that a common inner/outer logarithmic overlap is an
additional requirement.

\subsection{Optional smallest changes for a common leading overlap}
\label{kap:slope-option}

Velocity-overlap consistency alone requires
\begin{equation}
 \boxed{\alpha_m^{\rm match}=\frac{\kappa}{\sqrt{q_m}}
 =\kappa C_{\mu,*}^{1/4}.}
 \label{kap:slope-match}
\end{equation}
This is one effective wall-coefficient change. Keeping the original
baseline $c_\ell^w=2.55$ and the product parameterization
$\alpha=c_\mu^wc_\ell^w$ gives the equivalent choice
\begin{equation}
 c_{\mu,m}^{w,\rm slope}=\frac{\kappa}{2.55\sqrt{q_m}},
 \qquad c_\ell^w=2.55.
 \label{kap:one-coefficient}
\end{equation}
The corrected EP tail already uses \cref{kap:ell-equilibrium}; this
one-coefficient option does not require changing the uncorrected
dissipation baseline.

If one additionally requires the \emph{uncorrected} baseline tail
$K^{3/2}/(c_\ell^wz)$ to equal the outer tail and keeps the shared
product parameterization, the two conditions determine
\begin{equation}
 \boxed{
 c_{\ell,m}^{w,\rm match}=\kappa q_m^{3/2},\qquad
 c_{\mu,m}^{w,\rm match}=q_m^{-2},\qquad
 \alpha_m^{\rm match}=\kappa/\sqrt{q_m}.}
 \label{kap:shared-pair}
\end{equation}
These are constant reference values for the auxiliary wall law.
The outer variable-coefficient constitutive relation is unchanged.
Changing $c_\ell^w$ alone while retaining $c_\mu^w=0.09$ matches both
tails only if $0.09q_m^2=1$, which is not exactly true here.

Either viscosity-slope modification changes $D_A$, all $J_i$,
the forced momentum response, the reference production, and the
roughness calibration. All wall budgets and cutoff gates must be
recomputed. If the baseline $c_\ell^w$ is changed, its own cutoff
$z_{\eps,A}=2c_\ell^w\nu/\sqrt K$ changes as well.
The original roughness-correlation constants and $z_{vs}^*=10.7$
have no algebraic rule requiring multiplication by a ratio of
von Karman constants. Their further recalibration would be another
empirical wall-model modification.

The principal formulation in this note retains the original SCI
coefficients and claims the exact composite reference at the selected
$H$. Equations \eqref{kap:slope-match} and \eqref{kap:shared-pair}
are explicitly named optional variants. Their nonequilibrium accuracy
requires separate validation.

\subsection{Finite viscosity, roughness shift, and matching-height dependence}
\label{kap:height}

Even with \cref{kap:slope-match}, the active-branch derivative is
\begin{equation}
 U'_{A,0}(z)=\frac{u_m}{\kappa(z+\delta_A)},\qquad
 \delta_A=\frac{\nu(1-\alpha b)}{\kappa u_m}.
 \label{kap:shifted-gradient}
\end{equation}
An exact unshifted derivative throughout that branch additionally
requires $\delta_A=0$. Slope matching alone gives the leading overlap.
The intercept is a molecular/roughness contribution: at fixed nonzero
$h_r$, the original fully rough correlation has
$z_v\to-10.7h_r/70$ as $\nu\to0$, so a geometric shift can remain.

Holding $K,h_r,\nu$ and the wall coefficients fixed while varying the
matching height, \cref{awf:generated-roughness} implies
\begin{equation}
 \boxed{
 \frac{\dd\ln z_{0,m}(H)}{\dd\ln H}
 =1-\frac{\kappa u_m H}{D_A(H)}
 \longrightarrow
 1-\frac{\kappa}{\alpha\sqrt{q_m}}.}
 \label{kap:H-derivative}
\end{equation}
Here the limit means heights large relative to the active-branch
intercept scale. The slope adjustment removes the leading mismatch,
while the finite intercept generally leaves residual $H$ dependence.
Exact preservation at one calibrated interface therefore does not
establish matching-height independence.

At fixed $K,h_r,H,\nu$ and unchanged auxiliary coefficients, changing
$\kappa$ leaves $D_A,J_i,E_A$, and the inner zero-forcing velocity
unchanged. It changes the attached outer slope, generated $z_{0,m}$,
corrected dissipation amplitude, normalized budget $\Lambda_m$, and cutoff.
The dimensional budget $\int_0^H E_{c,m}\,\dd z=\int_0^H p_{A,0}\,\dd z$ stays fixed at these fixed wall inputs.
If a fixed physical $z_0$ is required, the resistance calibration must
be reconsidered. These dependencies are kept explicit in the
verification programme.

\subsection{Turbulent production and the corrected dissipation continuation}
\label{awf:dissipation}

Turbulence production contains only the eddy viscosity. Thus the
counterfactual zero-forcing reference production and current
flat-kernel production are
\begin{align}
 \boxed{p_{A,0}(z;K)
 =\nu_{t,A}(z)\left[\frac{t_m}{D_A(z)}\right]^2
 =t_m^2\left[\frac1{D_A(z)}-\frac{\nu}{D_A(z)^2}\right],}
 \label{awf:reference-production}\\
 \boxed{P_A(z)=\nu_{t,A}(z)
 \frac{|\boldsymbol t_A+\boldsymbol F L_r(z)|^2}{D_A(z)^2}\geq0.}
 \label{awf:current-production}
\end{align}
Using $D_A|U_A'|^2$ as turbulent production would incorrectly
transfer molecular mean-flow dissipation into $k$.
Define the layer mean $P_0$ and normalized budget $\Lambda_m$ by
\begin{equation}
 \begin{aligned}
 HP_0&=\int_0^Hp_{A,0}\,\dd z
     =t_m^2[J_0(H)-\nu J_2(H)],\\
 \Lambda_m&:=\frac{HP_0}{A_m}
     =\kappa u_m[J_0(H)-\nu J_2(H)].
 \end{aligned}
 \label{awf:production-budget}
\end{equation}
Evaluate the nonnegative integral directly near a molecular-only
state, rather than subtracting nearly equal $J_0$ and $\nu J_2$.
This finite reference production is not an integral of the singular
outer profile $A_m/z$ down to the ground.

Choose the positive, continuous, nonincreasing corrected profile
\begin{equation}
 E_{c,m}(z)=\frac{A_m}{\max(z,z_c)}.
 \label{awf:corrected-epsilon}
\end{equation}
For $z_c\leq H$, its integral is $A_m[1+\ln(H/z_c)]$.
Matching that integral to $HP_0$ gives
\begin{equation}
 \boxed{z_c=H\exp(1-\Lambda_m).}
 \label{awf:cutoff}
\end{equation}
If the tail must agree with $A_m/z$ by $z_m$, the admissibility gate is
\begin{equation}
 \boxed{z_c\leq z_m\leq H
 \quad\Longleftrightarrow\quad
 \Lambda_m\geq1+\ln(H/z_m).}
 \label{awf:gate}
\end{equation}
The gate is necessary for any nonnegative nonincreasing continuation
that joins $A_m/z_m$ and equals $A_m/z$ above $z_m$: its integral
below $z_m$ is at least $A_m$. The displayed plateau/tail profile
proves sufficiency. A strict margin is preferable on a switching
plateau. If $\nu_{t,A}=0$ throughout the layer, $\Lambda_m=0$ and
the gate fails. Clipping an inadmissible cutoff changes the budget
and removes the claimed equilibrium.

Let
\begin{equation}
 M_{\nu,m}=\kappa u_m\nu J_2(H)>0.
 \label{awf:molecular-budget}
\end{equation}
Combining the original resistance anchor with the production budget
gives the exact finite-viscosity identity
\begin{equation}
 \boxed{z_c=e\,z_{0,m}(K)\exp(M_{\nu,m}).}
 \label{awf:cutoff-roughness}
\end{equation}
At the fixed reference satisfying
\cref{awf:roughness-compatibility}, $z_{0,m}(K_0)=z_0$ and
\begin{equation}
 \Lambda_m=\ln(H/z_0)-M_{\nu,m},\qquad
 \boxed{e z_0\exp(M_{\nu,m})\leq z_m.}
 \label{awf:finite-nu-gate}
\end{equation}
Thus $z_c=e z_0$ is not exact with finite $\nu$; it is recovered
only when the molecular contribution $M_{\nu,m}$ is negligible.
Keeping the old cutoff would give
$\int_0^H(E_{c,m}-p_{A,0})\,\dd z=\nu t_m^2J_2(H)>0$,
so the redistribution would cease to be conservative.


\subsection{Pointwise and all-mode kinetic-energy compatibility}
\label{awf:moments}

Let $0\leq\chi\leq1$ be constant along each wall-normal ray.
First define the dependent dissipation and the effective destruction
\begin{equation}
 \begin{aligned}
 E_{\mathrm{pre}}&=(1-\chi)E_A+\chi E_{c,m},\\
 D_{\mathrm{pre}}&=(1-\chi)E_A+\chi p_{A,0},\\
 C_{\mathrm{pre}}&=E_{\mathrm{pre}}-D_{\mathrm{pre}}
                 =\chi(E_{c,m}-p_{A,0}).
 \end{aligned}
 \label{awf:pre-profiles}
\end{equation}
Equal layer means alone leave nonconstant Galerkin source moments:
$\int_0^H\phi_i(p_{A,0}-E_{c,m})\,\dd z$ is generally nonzero.
The required signed inner redistribution is $C_{\mathrm{pre}}$.
Its integral is zero, and the resulting inner source is
\begin{equation}
 \boxed{
 S_{k,w}=P_A-E_{\mathrm{pre}}+C_{\mathrm{pre}}
 =P_A-D_{\mathrm{pre}}.
 }
 \label{awf:net-source}
\end{equation}
At the flat reference, $\chi=1$, $k=K_0$, $F=0$, and
$U_H=t_mJ_0(H)$. Then $P_A=p_{A,0}$ and $S_{k,w}=0$
pointwise. Therefore every active source load satisfies
\begin{equation}
 \boxed{\int_{\Omega_w}\rho\phi_iS_{k,w}\,\dd V=0.}
 \label{awf:all-moments}
\end{equation}
This is an explicit new inner transport term, not a positive
dissipation profile or a quadrature adjustment. Assemble the simplified
source in \cref{awf:net-source}. Use the counterfactual $p_{A,0}$
at the same current $K$, never current $P_A$ in its place.

Equivalently, define
\begin{equation}
 j_c(z)=\int_0^z[E_{c,m}(\xi)-p_{A,0}(\xi)]\,\dd\xi,\qquad
 j_c(0)=j_c(H)=0,\qquad C_{\mathrm{pre}}=\partial_z(\chi j_c).
 \label{awf:redistribution-flux}
\end{equation}
Its weak action is $-\int\partial_z\phi_i\,\chi j_c\,\dd V$.
Use the source form or this flux action, never both.

\section{Wall-element balances, field constraints, and interface coupling}
\label{awf:properties}

\subsection{Mean-flow and turbulent-energy accounting}

The reference mechanical-power identity is
\begin{equation}
 \begin{aligned}
 t_mU_{H,0}
 &=\int_0^HD_A(U'_{A,0})^2\,\dd z\\
 &=\int_0^Hp_{A,0}\,\dd z+\nu t_m^2J_2(H)\\
 &=\boxed{\int_0^HE_{c,m}\,\dd z+\nu t_m^2J_2(H)} .
 \end{aligned}
 \label{awf:energy}
\end{equation}
The final term is molecular dissipation of mean kinetic energy; it
is not part of the modeled dissipation of turbulent kinetic energy.
The stationary ground does no work because its velocity is zero.
For a forced column, direct integration of the momentum equation gives
\begin{equation}
 \boldsymbol T_H\cdot\boldsymbol U_H
 =\int_0^HD_A|\boldsymbol U_A'|^2\,\dd z
  +\int_{h_r}^H\boldsymbol F\cdot\boldsymbol U_A\,\dd z.
 \label{awf:forced-energy}
\end{equation}
The forcing work must be retained and need not have one sign.
Zero integrated redistribution does not mean that the full blend
retains the baseline loss: changing that loss is the intended
equilibrium correction.

\subsection{Transported and dependent fields, traction, and pressure}
\label{awf:interface}

Retain $k_h$ as a transported field throughout the wall and outer
regions. In the wall layer use
$D_{k,w}=\nu+\nu_{t,A}/\sigma_k>0$ with $\sigma_k=1$, so that
$D_{k,w}=D_A$ for the selected value. Use reconstructed velocity in advection and
\cref{awf:net-source}. Use zero modeled ground $k$ flux, a common
$k$ trace at $H$, and one shared diffusive and advective interface flux.
Constant reference $k=K_0$ has zero diffusion on both sides.
Retaining molecular viscosity does not impose $k=0$ at the numerical
ground: that would change the wall-function model and exclude the
constant-$K_0$ reference. The selected condition remains zero modeled
ground $k$ flux.

Inner dissipation is dependent: before the positivity adaptation below,
$\eps_w=E_{\mathrm{pre}}$. Remove independent inner dissipation
transport degrees of freedom and impose only
\begin{equation}
 \boxed{\eps_{o,H}=E_{\mathrm{pre}}(H).}
 \label{awf:epsilon-trace}
\end{equation}
At equilibrium it is $A_m/H$. Do not also impose an independent
dissipation flux at $H$. The outer equation, using its corrected
$\sigma_{\eps,m}$, determines that flux. This closure does not claim
a solved inner dissipation transport conservation law.

For pressure and curved-flow completion, take
\begin{equation}
 \begin{aligned}
 R_A&=\tfrac23kI-2\nu_{t,A}S_w,\qquad
 S_w=\tfrac12[\nabla\boldsymbol U_w+(\nabla\boldsymbol U_w)^T],\\
 T_w&=-(\pi_w+\tfrac23k)I+2D_AS_w,\qquad
 T_o=-\pi_oI-R_m .
 \end{aligned}
 \label{awf:inner-outer-stress}
\end{equation}
The kinematic pressures use the same hydrostatic convention.
This full-strain inner completion is an additional constitutive
choice introduced here. The scalar SCI AWF does not itself specify
this three-dimensional tensor closure. Its incompressible production is $P_A=2\nu_{t,A}S_w:S_w$,
reducing to \cref{awf:current-production} in parallel flat shear.
For only weakly divergence-free fields, the full contraction includes
$-(2k/3)\nabla\cdot\boldsymbol U_w$ and must be treated consistently
with the isotropic-stress pressure convention. The molecular mean-flow
loss is separately $2\nu S_w:S_w$. In particular, molecular stress
belongs in $T_w$ but not in the Reynolds covariance $R_A$.

Use a shared velocity trace and full traction
$T_w\boldsymbol N_H=T_o\boldsymbol N_H$.
Its normal component gives
\begin{equation}
 \pi_w-\pi_o
 =R_{NN,o}-\tfrac23k+2D_A\boldsymbol N_H\cdot S_w\boldsymbol N_H,
 \qquad R_{NN,o}=\boldsymbol N_H\cdot R_m\boldsymbol N_H.
 \label{awf:normal-pressure}
\end{equation}
At the flat reference $S_{zz,w}=0$, so this still requires
\begin{equation}
 \boxed{\pi_w-\pi_o=R_{zz,o,0}-\tfrac23K.}
 \label{awf:pressure-offset}
\end{equation}

\begin{table}[htbp]
\centering
\caption{Required flat-reference pressure offset. The rows use the
corresponding $q_m$; the offset multiplies $K$.}
\label{awf:pressure-table}
\begin{tabular}{lll}
\toprule
Outer model & $(\pi_w-\pi_o)/K$ & Approximate value\\
\midrule
A93 & $-17q_m^2/[3(1000+q_m^3)]$ & $-0.06307362186$\\[1mm]
A95 & $-\dfrac{q_m^2\sqrt{1-9/(2q_m^2)}}{1+3q_m^2}$
    & $-0.27109100656$\\[2mm]
E94/E95 & $0$ & $0$\\
\bottomrule
\end{tabular}
\end{table}

It is total traction that is continuous. Independent mechanical-pressure
continuity would add an unintended surface force for A93/A95.
The offset is a consequence of the model interface, not a pressure
jump attributed to a material interface. Retain the relative pressure
modes until traction matching determines them, with one compatible
global gauge.

At the outer top $Z$, the reference requires tangential traction
$+\rho\ust^2\boldsymbol e_x$, zero $k$ flux, and natural dissipation
flux $(\nut/\sigma_{\eps,m})\partial_n\eps
=-\ust^4/(\sigma_{\eps,m}Z)$ for $n=+\boldsymbol e_z$.
Compatible analytic values are alternatives to these conditions.
A stress-free top or zero dissipation gradient does not admit the
constant-stress reference.


\subsection{Explicit forcing response and its scope}
\label{awf:response}

For fixed $K,H,h_r,\nu$ and coefficient variant,
\begin{equation}
 \frac{\partial\boldsymbol t_A}{\partial\boldsymbol F}=-BI_\parallel,
 \qquad
 \frac{\partial\boldsymbol T_H}{\partial\boldsymbol F}
 =[(H-h_r)-B]I_\parallel,
 \qquad
 \frac{\partial\boldsymbol T_H}{\partial\boldsymbol U_H}
 =J_0(H)^{-1}I_\parallel.
 \label{awf:frozen-response}
\end{equation}
Since $0\le L_r\le H-h_r$ under the positive resistance weight,
$0\le B\le H-h_r$. The ground and upper-interface forcing responses
are different. The EP dissipation/source construction preserves these
explicit derivatives when the original viscosity coefficients are
retained. It does not preserve the full coupled solution derivative:
$K$, the forcing sampler, pressure, and the transported fields can
respond to the solution. The optional slope recalibration changes
$J_i$ and therefore changes the fixed-input wall response as well.

\subsection{A conditional exact composite equilibrium}
\label{awf:equilibrium-theorem}

\begin{proposition}
Choose an outer model with its compatible $\kappa$ and dissipation
coefficient. Let $K_0=q_m\ust^2>0$, $0\le h_r<H$, and $\nu>0$.
Assume the resistance condition \eqref{awf:roughness-compatibility},
the cutoff gate \eqref{awf:gate}, and compatible outer boundary data.
On a flat unit plateau of the correction switch, the composite
velocity \eqref{awf:composite-velocity}, constant $k=K_0$, dependent
inner dissipation $E_{c,m}$, outer dissipation $A_m/z$, and pressure
offset \eqref{awf:pressure-offset} satisfy the specified stationary
inner/outer system and its full interface conditions.
\end{proposition}
\begin{proof}
The original inner momentum equation has $F=0$ and constant total
shear $t_m=\ust^2$, with no slip and the resistance-matched upper
velocity. The compatible outer reference satisfies its momentum,
constitutive, and turbulence equations. Full tangential and normal
traction match by \cref{awf:derivative-flux,awf:pressure-offset}.
Constant $k$ has zero diffusion and advection. Its inner source is
$P_A-p_{A,0}=0$ pointwise, so all weak source moments vanish.
The corrected dissipation trace is $A_m/H$ because $z_c\le H$.
The source redistribution has zero endpoint flux and zero integrated
budget. Remaining boundary conditions are those of the compatible
outer reference. The positivity modification introduced later is
inactive on an open neighborhood of this reference.
\end{proof}

This proposition is conditional on the specified flat reference and
admissibility conditions. It does not assert a curved-wall equilibrium,
an arbitrary prescribed pair $(h_r,z_0)$, a common logarithmic gradient
at $H$, matching-height independence, or exact stationarity of a
polynomial interpolation under arbitrary discrete operators.


\section{Finite-volume interpretation of the analytical wall element}
\label{geo:fv}

\subsection{Which interval is being eliminated?}

The analytical interval is the physical layer $0<z<H$ beneath a
specified interface with the outer equations.  Its upper trace is
$\boldsymbol U_H$, its ground velocity is zero, and its upper shear
is $\boldsymbol T_H$.  These are distinct from a cell-centre velocity,
a ground-face stress, and a first-cell dissipation value.  A
finite-volume implementation must declare how those numerical quantities
are reconstructed from the model traces.  In particular, identifying
$H$ with a first-cell centre while integrating a different control
volume requires an explicit transfer; it is not an identity of the
continuous wall model.

The closed column is an elimination of its momentum profile under the
stated flat, quasisteady approximation.  With frozen $K$ and forcing,
it gives a resistance relation between the trace velocity and the
\emph{upper} traction.  The corresponding implicit contribution has
positive velocity derivative $J_0(H)^{-1}I_\parallel$.  The difference
$\boldsymbol T_H-\boldsymbol t_A=\boldsymbol F(H-h_r)$ accounts for
forcing within the eliminated interval.  Using $\boldsymbol t_A$ at
its upper face would omit that integrated force.  Retaining transient
storage or tangential transport instead requires the patch equations
in \cref{geo:curved}; their elimination contains history and lateral
coupling as well as instantaneous resistance.

For subsequent sections, denote the profiles \emph{actually used} by
\begin{equation}
 \begin{aligned}
 (E_w,D_w,C_w)&=
 \begin{cases}
 (E_{\mathrm{pre}},D_{\mathrm{pre}},C_{\mathrm{pre}}),
       &\text{preliminary profiles},\\
 (E_{\mathrm{pos}},D_{\mathrm{pos}},C_{\mathrm{pos}}),
       &\text{positivity-adapted profiles},
 \end{cases}\\
 C_w&=E_w-D_w .
 \end{aligned}
 \label{geo:used-profiles}
\end{equation}
The positive profiles are defined in \cref{st:positive-profiles}.
Thus $\eps_w=E_w$ and the assembled inner reaction is $P_A-D_w$.
Here $D_w$ is a destruction profile, whereas $D_A=\nu+\nu_{t,A}$
is a momentum diffusivity.  Their dimensions and roles differ.
All selected wall coefficients, including an optional overlap-slope
adjustment, are fixed parts of the model variant being compared.

\subsection{Cell averages and the source/flux choice}

For a flat cell spanning the complete layer, the turbulent-energy
balance per unit wall area, suppressing lateral terms only for this
illustration, is
\begin{equation}
 \frac{\dd}{\dd t}\int_0^H k\,\dd z
 +[wk-D_{k,w}\partial_z k]_0^H
 =\int_0^H(P_A-D_w)\,\dd z.
 \label{geo:fv-balance}
\end{equation}
At the reference, advection and diffusion vanish and the right-hand
side is zero.  Using $E_{c,m}$ without redistribution also gives a
zero \emph{whole-layer} mean at that reference, by the cutoff budget.
That observation explains why a mean balance can suffice for a
single wall-cell source test.  It does not determine the source in
individual subcells or against nonconstant test functions.

If the layer is resolved by several control volumes, retain the local
source $P_A-D_w$ in each cell, or use $P_A-E_w$ together with the
conservative redistribution flux.  For the unadapted profiles this
flux is $-\chi j_c\boldsymbol e_z$ on the left-hand side of the
conservation equation.  Fluxes at an internal subcell face enter
its two neighbours with opposite signs; the redistribution is zero
at the ground and at $H$.  Adding this flux to a cell source already
formed as $P_A-D_w$ counts the redistribution twice.

The $k$ field remains transported.  A prescribed inner dissipation
is instead a dependent model field.  In a cell formulation its
independent transport equation must be replaced by that constraint,
or the corresponding unknown eliminated.  Solving an independent
inner $\eps$ equation and subsequently overwriting its solution is
a different algorithm whose remaining residual and conservation
properties require separate analysis.  At the outer interface use
the single value $E_w(H)$; the outer dissipation equation determines
its flux.

\section{Spectral/\texorpdfstring{$hp$}{hp} fields, no slip, and all-mode compatibility}
\label{geo:hp}

\subsection{A field formulation rather than a pointwise wall overwrite}

The unknowns of the composite problem are velocity and pressure in
the wall patch and outer region, transported $k$ in both regions,
and transported $\eps_o$ only in the outer region.  The inner
$\eps_w=E_w$ is dependent.  Inner velocity and pressure may be
numerically condensed onto the interface, provided their storage,
continuity, pressure modes, and lateral fluxes are retained before
condensation.  The closed flat column is a useful exact reduction
and reference, not an additional independent momentum equation to
be imposed on a general patch solution.
\label{geo:fields}

The sample $K=\mathcal K[k_h]$ is a specified physical functional.
A fixed physical sampling height or a fixed positive normalized
average is admissible; moving to the first interior interpolation
point whenever polynomial order changes would change the model.
The same sample is used in $D_A$, $E_A$, the counterfactual $p_{A,0}$,
the cutoff, and the state sensors.  Constant preservation,
$\mathcal K[K_0]=K_0$, is necessary for the reference.  Positivity of
this sample does not prove pointwise positivity of the transported
field.

Strong no slip means that the selected inner trial velocity has
$\boldsymbol U_w=0$ at the numerical ground.  It can be enforced by
a constrained patch space or by a reconstruction/enrichment whose
ground trace is identically zero.  The outer problem sees the
traction and velocity at $H$.  A natural load placed on the ground
of an otherwise unchanged conforming velocity space cannot replace
this construction: with strongly imposed ground velocity, its free
test functions vanish there.  The interface traction must act on
the actual retained interface degrees of freedom.

The wall layer does not set $k=0$ at the ground.  Its specified
modeled $k$ flux is zero, and constant $K_0$ remains admissible.
Molecular viscosity in the momentum and $k$ diffusivities does not
turn this wall-function layer into a resolved viscous-sublayer
transport model for all turbulence variables.

\subsection{Continuity supplies transpiration at the matching surface}

Even on a flat wall, a tangentially varying wall field generally
requires a nonzero vertical velocity.  With
$\boldsymbol u_\parallel=(U_x,U_y,0)$, physical coordinates $(x,y,z)$,
and $w=U_z$, continuity and ground impermeability give
\begin{equation}
 w(x,y,z)=-\int_0^z
 \nabla_\parallel\cdot\boldsymbol u_\parallel(x,y,\zeta)
 \,\dd\zeta .
 \label{geo:flat-transpiration}
\end{equation}
For constant $H$, the outer trace has $w_H$ from this relation.
For variable $H(x,y)$, the flux through the graph includes its slope,
as derived in \cref{geo:massflux}.  Imposing an independent zero
normal velocity at that graph generally violates the integrated
continuity equation.  Conversely, computing a continuity-compatible
normal velocity does not eliminate the normal momentum equation or
its pressure balance.

\subsection{Every source mode and the physical quadrature}

For a basis/test function $\phi_i$, the wall source load is
\begin{equation}
 b_i^{k,w}=\rho\int_{\Omega_w}\phi_i(P_A-D_w)\,\dd V.
 \label{geo:wall-load}
\end{equation}
On the admissible flat reference the positivity adaptation is
inactive, $\chi=1$, and the same analytical reconstruction gives
$P_A=p_{A,0}$.  Therefore \cref{awf:all-moments} holds for every
active test function.  The weaker identity
$\int(P_A-E_{c,m})\,\dd z=0$ checks only the constant test; it
cannot be promoted to an all-mode statement by increasing the
quadrature order.  The signed redistribution in
\cref{awf:pre-profiles,awf:net-source} is the specified inner transport
needed for the latter statement.

Use the actual basis functions, physical volume Jacobian, and
common interface transfer.  Split normal integration at every
interior positive breakpoint among $h_r,z_v,z_{\eps,A},z_c$ and
any active positivity-profile breakpoint.  A negative $z_v$ remains
in the viscosity formula; it is not replaced by zero.  Reciprocal,
logarithmic, maximum, and state-dependent coefficient functions are
not generally polynomial, so integration convergence must be
checked independently of polynomial order.  For cancellation,
evaluate current and counterfactual reference production through a
consistent arithmetic path.  For verification, also integrate their
budgets independently, so an identical coding error is not mistaken
for conservation.

These source identities do not alone make an interpolated logarithm
an exact stationary state of the assembled momentum and outer
$\eps$ operators.  Representation, differentiation, pressure,
interface transfer, and time stages have additional requirements;
\cref{num:implementation} separates them.

\section{Curved terrain: conservative Cartesian completion}
\label{geo:curved}

The following completion defines a model on a smooth stationary
wall patch.  Physical velocity and stress components always refer
to $(x,y,z)$.  Surface parameters describe the geometry only; they
do not replace those components by a moving streamwise/normal/binormal
component system.  The geometric identities are exact.  The
full-strain inner stress and the roughness forcing extension are
constitutive choices whose terrain predictions require validation.

\subsection{Normal map, metric, and geometric admissibility}

Let $\boldsymbol X(\xi^1,\xi^2)$ describe the ground surface
$\Gamma_w$, and let $\boldsymbol n$ point from the ground into the
fluid.  Write
\begin{equation}
 \begin{gathered}
 \boldsymbol x=(x,y,z)=\boldsymbol X(\boldsymbol\xi)
                    +r\boldsymbol n(\boldsymbol\xi),
 \qquad 0\leq r\leq H(\boldsymbol\xi),\qquad
 P=I-\boldsymbol n\otimes\boldsymbol n,\\
 B_g=\nabla_\Gamma\boldsymbol n,\qquad
 A_g=I_T+rB_g,\qquad J=\det_T A_g,\qquad
 B_r=B_gA_g^{-1},\qquad \dd V=J\,\dd S_w\,\dd r.
 \end{gathered}
 \label{geo:map}
\end{equation}
Here $r$ is physical normal distance, not generally the Cartesian
height $z$.  Within this geometric section, $r$ denotes that distance;
it is unrelated to the rotation norm used in the constitutive
notation.  The convention makes the principal curvatures
$\kappa_i$ of an outward-oriented sphere positive.
The inverse of $A_g$ acts on the tangent plane and is extended by
zero on the normal direction when used in tensor fluxes.  At fixed
$r$ define
\begin{equation}
 \nabla_r f=A_g^{-T}\nabla_\Gamma f,\qquad
 \nabla f=\nabla_r f+\boldsymbol n\partial_r f,
 \qquad \nabla\boldsymbol n=B_r.
 \label{geo:metric-operators}
\end{equation}
Apply the gradient componentwise to Cartesian vector components.
In particular, differentiation of a decomposition involving
$\boldsymbol n$ or $P$ includes their spatial derivatives.

The required geometric and model domain is
\begin{equation}
 \begin{gathered}
 1+r\kappa_i\geq a_{\min}>0\quad (0\leq r\leq H),
 \qquad \text{the normal map is globally injective},\\
 H>0,\qquad 0\leq h_r<H,\qquad 0<r_K,r_m\leq H,
 \qquad K>0,\qquad \nu>0 .
 \end{gathered}
 \label{geo:admissibility}
\end{equation}
The first condition gives local nonsingularity and $J>0$; it does
not rule out intersections of distant normal rays.  A layer lying
within the one-sided normal reach supplies a sufficient global
prescription.  The numerical element map and the inner/outer
interface maps also require positive Jacobians, compatible normals,
and compatible measures.  Between-node geometry needs checking;
validity at a finite set of solution nodes is insufficient.  Corners
and nonsmooth terrain require separate patches or a declared smooth
geometric approximation.  The sampling distance $r_K$, tail-start
$r_m$, and matching height are physical modeling choices.  A thin
layer relative to curvature radii and tangential flow-variation scales
concerns model adequacy in addition to coordinate invertibility.

For a stationary, spatially variable $H$, define
\begin{equation}
 \boldsymbol q_H=A_{g,H}^{-T}\nabla_\Gamma H,
 \qquad \beta_H=\sqrt{1+|\boldsymbol q_H|^2},\qquad
 \boxed{\boldsymbol N_H=\frac{\boldsymbol n-\boldsymbol q_H}{\beta_H},
 \qquad \dd S_H=J_H\beta_H\,\dd S_w.}
 \label{geo:interface-geometry}
\end{equation}
$\boldsymbol N_H$ points into the outer region, whose lower outward
normal is $\boldsymbol n_o=-\boldsymbol N_H$.  A time-dependent
moving interface would additionally need moving-boundary fluxes;
it is outside this stationary-geometry specification.

\subsection{Mass conservation and the interface normal velocity}

Decompose $\boldsymbol U=\boldsymbol u_t+w\boldsymbol n$, with
$\boldsymbol u_t=P\boldsymbol U$, and define
$\boldsymbol v=A_g^{-1}\boldsymbol u_t$.  The exact transformed
continuity equation is
\begin{equation}
 \boxed{\partial_r(Jw)+\operatorname{div}_\Gamma(J\boldsymbol v)=0.}
 \label{geo:continuity}
\end{equation}
Ground no slip gives $w(0)=0$ and $\boldsymbol u_t(0)=0$, hence
$J(r)w(r)=-\int_0^r\operatorname{div}_\Gamma(J\boldsymbol v)\,\dd\zeta$.
The variable-limit rule then yields the interface mass flux per unit
ground area,
\begin{equation}
 \boxed{\begin{aligned}
 m_H&=J_H\beta_H\widehat{\boldsymbol U}\cdot\boldsymbol N_H
     =J_H\bigl(w_H-\boldsymbol v_H\cdot\nabla_\Gamma H\bigr)\\
    &=-\operatorname{div}_\Gamma\int_0^H J\boldsymbol v\,\dd r,
 \qquad
 g_n=\boldsymbol U_o\cdot\boldsymbol n_o
     =-\frac{m_H}{J_H\beta_H} .
 \end{aligned}}
 \label{geo:massflux}
\end{equation}
Here $\widehat{\boldsymbol U}$ is the full three-component shared
trace.  The height-slope term is essential even on a flat wall.
Share this weighted flux between inner continuity, outer continuity,
scalar advection, and pressure coupling.  Separate nodal normal
velocities whose projected integrated fluxes differ do not meet
this requirement.

\subsection{Momentum, full strain, and the roughness force}

For a kinematic total stress $T$, including pressure, Cartesian
momentum conservation becomes
\begin{equation}
 \boxed{\begin{aligned}
 J\partial_t\boldsymbol U
 &+\partial_r\{J(w\boldsymbol U-T\boldsymbol n)\}\\
 &+\operatorname{Div}_\Gamma
 \{J(\boldsymbol U\otimes\boldsymbol v-TA_g^{-T})\}
 =J\boldsymbol f .
 \end{aligned}}
 \label{geo:momentum}
\end{equation}
$\operatorname{Div}_\Gamma$ takes the surface divergence of each
Cartesian row.  These metric fluxes already include the geometric
effects of pressure and stress; additional copies of curvature
forces must not be inserted into this conservative equation.

Use $T_w,T_o$ from \cref{awf:inner-outer-stress}.  In the wall
layer, $D_A$ and $\nu_{t,A}$ are evaluated at distance $r$ and sample
$K$, while the isotropic covariance uses the local transported $k$.
For incompressible flow the production is
$P_A=2\nu_{t,A}S_w:S_w$.  The molecular mean-flow loss is
$2\nu S_w:S_w$ and is excluded from $P_A$.  If only weak
incompressibility is enforced, the contraction
$-R_A:\nabla\boldsymbol U_h$ contains the additional term
$-(2k_h/3)\nabla\cdot\boldsymbol U_h$; its treatment must agree with
the chosen isotropic-stress pressure convention and discrete energy
accounting.

Let $\boldsymbol f_\pi$ be the body force remaining after the same
hydrostatic pressure convention used in the outer equations.  One
explicit conservative extension of the original forcing suppression
below roughness height is
\begin{equation}
 \begin{aligned}
 a_r(r)&=\mathbf1_{r>h_r},\qquad
 \boldsymbol f_w=\boldsymbol f_\pi+(1-a_r)\boldsymbol F_A,\\
 \boldsymbol F_A&=\mathcal E_H^+P\left[
 \operatorname{Div}_r(\boldsymbol u_{o,t}\otimes\boldsymbol u_{o,t})
       +\nabla_r\pi_o-\boldsymbol f_\pi\right],
 \qquad \boldsymbol u_{o,t}=P\boldsymbol U_o .
 \end{aligned}
 \label{geo:rough-force}
\end{equation}
For a tensor $Q$ with $Q\boldsymbol n=0$, define
$\operatorname{Div}_rQ=J^{-1}\operatorname{Div}_\Gamma(JQA_g^{-T})$.
The sampler $\mathcal E_H^+$ takes the one-sided outer values after
computing derivatives at fixed $r$, before following the graph
$r=H(\boldsymbol\xi)$.  Extend $\boldsymbol F_A$ constantly along
each normal ray.  Hydrostatic gravity and pressure must cancel
consistently in this sampler.

Recovering the scalar column requires a declared parallel-flow
approximation in addition to flatness: neglect inner storage,
normal advection, lateral viscous transport, and the terms involving
$\nabla_\parallel w$ in the tangential viscous shear. Replace the
remaining inner convection, pressure-gradient, and body-force balance
by the same height-constant sample $\boldsymbol F_A$. In particular,
this replaces the inner effective-pressure gradient
$\nabla_\parallel(\pi_w+2k/3)$ by the sampled outer mechanical-pressure
gradient. The discarded gradient
$\nabla_\parallel(\pi_w+2k/3-\pi_o)$ is not zero in a general
inhomogeneous flow. With these approximations the tangential equation
reduces to $(D_A\boldsymbol U_\parallel')'=a_r\boldsymbol F_A$.
All discarded terms and gradients vanish on the flat homogeneous
reference, so its exactness is unaffected.
In the conservative curved or unsteady patch the full acceleration,
pressure--normal-stress balance, and lateral transport remain in
\cref{geo:momentum}; the scalar reduction is not imposed as a second
equation on that patch.  The additional force in
\cref{geo:rough-force} has no guaranteed drag sign; include its work
in energy budgets.  Multiplying the entire residual or its mass
term by $a_r$ would suppress different physics and would introduce
a degenerate storage operator.

\subsection{Normal pressure and curved shear}

Put $\pi_e=\pi_w+2k/3$, $V=2D_AS_w$,
$\boldsymbol t=PV\boldsymbol n$, $V_{nn}=\boldsymbol n\cdot V\boldsymbol n$,
and $V_{tt}=PVP$.  The tangential traction on a base-normal surface is
\begin{equation}
 \boxed{\boldsymbol t=D_A\left[
 \partial_r\boldsymbol u_t+\nabla_r w-B_r\boldsymbol u_t\right].}
 \label{geo:curved-shear}
\end{equation}
It includes tangential gradients of normal velocity and curvature.
The normal acceleration and pressure equation following from
\cref{geo:momentum} are
\begin{align}
 a_n&=(\partial_t+w\partial_r+
          \boldsymbol v\cdot\nabla_\Gamma)w
       -\boldsymbol u_t\cdot B_r\boldsymbol u_t,
 \label{geo:normal-acceleration}\\
 \partial_r\pi_e
 &=\boldsymbol n\cdot\boldsymbol f_\pi-a_n
   +J^{-1}\partial_r(JV_{nn})
   +\operatorname{div}_r\boldsymbol t-V_{tt}:B_r,
 \qquad V_{nn}=2D_A\partial_rw,
 \label{geo:normal-pressure}
\end{align}
where $\operatorname{div}_r\boldsymbol t
=J^{-1}\operatorname{div}_\Gamma(JA_g^{-1}\boldsymbol t)$.
The roughness force is base-tangential, so contributes no normal
component here.  For an outward circular offset of radius $R+r$
and tangential speed $u$, the curvature acceleration is
$-u^2/(R+r)$.  With no other normal forces or stress gradients the
corresponding pressure gradient is $+u^2/(R+r)$, fixing the sign
convention.

More generally, write $T_o=-\pi_oI+T'_o$, with
$T'_o=-R_m$ for the high-Re outer equations.  Then
\begin{equation}
 \partial_r(\pi_o-T'_{o,nn})
 =\boldsymbol n\cdot\boldsymbol f_\pi-a_{o,n}
  +\operatorname{div}_r(P T'_o\boldsymbol n)
  +(\operatorname{tr}B_r)T'_{o,nn}-B_r:(PT'_oP).
 \label{geo:outer-normal-pressure}
\end{equation}
Thus anisotropic normal stresses as well as curvature enter the
outer pressure balance.  A globally hydrostatic mechanical pressure
is generally incompatible with the hill or transition solution.

\subsection{Conservative coupling of pressure, velocity, and momentum}

Integrating \cref{geo:momentum} through the variable-height layer
gives
\begin{equation}
 \begin{aligned}
 \partial_t\int_0^H J\boldsymbol U_w\,\dd r
 &+\operatorname{Div}_\Gamma\int_0^H
   J(\boldsymbol U_w\otimes\boldsymbol v-T_wA_g^{-T})\,\dd r\\
 &+m_H\widehat{\boldsymbol U}
  -J_H\beta_HT_w\boldsymbol N_H+T_w(0)\boldsymbol n
 =\int_0^H J\boldsymbol f_w\,\dd r .
 \end{aligned}
 \label{geo:integrated-momentum}
\end{equation}
The ground stress term has the displayed sign because the layer's
lower outward normal is $-\boldsymbol n$.  On $\Gamma_H$ impose
\begin{equation}
 \boldsymbol U_w=\boldsymbol U_o=\widehat{\boldsymbol U},\qquad
 T_w\boldsymbol N_H=T_o\boldsymbol N_H,\qquad
 \widehat{\boldsymbol f}_M=
 (\widehat{\boldsymbol U}\otimes\widehat{\boldsymbol U}
                         -\widehat T)\boldsymbol N_H .
 \label{geo:interface-momentum}
\end{equation}
Use opposite orientations of this same flux, and the same physical
interface measure, in the two regions.  The corresponding conservative
weak momentum form in either subdomain $\Omega_\ell$ is
\begin{equation}
 (\boldsymbol v_h,\partial_t\boldsymbol U)
 -(\nabla\boldsymbol v_h,\boldsymbol U\otimes\boldsymbol U-T)
 +\langle\boldsymbol v_h,
        \widehat{(\boldsymbol U\otimes\boldsymbol U-T)}
        \boldsymbol n_\ell\rangle_{\partial\Omega_\ell}
 =(\boldsymbol v_h,\boldsymbol f).
 \label{geo:momentum-weak}
\end{equation}
Internal loads cancel for common traces or compatible mortar tests.
For an outer right-hand-side stress load specifically, the lower
normal gives $-\rho\int_{\Gamma_H}\boldsymbol v_h\cdot
T_w\boldsymbol N_H\,\dd S_H$.

The normal component of \cref{geo:interface-momentum} is precisely
\cref{awf:normal-pressure}.  Relative pressure modes must be kept
until traction matching determines them; a single global gauge
fixes the remaining nullspace.  An isolated incompressible patch
cannot be assigned all interface velocity components and an
additional independent pressure trace.  A pressure space or lifting
must permit the model's stress-dependent pressure offset.  If the
solver absorbs $2k/3$ into pressure, it must still retain the
anisotropic part of the normal traction.

On the fixed interface, normal projection of the outer momentum
also gives the compatibility relation
\begin{equation}
 \partial_{n_o}\pi_o=-\partial_tg_n
 -\boldsymbol n_o\cdot[(\boldsymbol U_o\cdot\nabla)\boldsymbol U_o]
 +\boldsymbol n_o\cdot(\nabla\cdot T'_o+\boldsymbol f_\pi).
 \label{geo:pressure-compatibility}
\end{equation}
This determines what a pressure stage must be consistent with;
it is not an extra freely assignable boundary condition.

\subsection{Retained \texorpdfstring{$k$}{k}, metric redistribution, and energy}

For any transported scalar $\phi$ with diffusivity $D_\phi$, its
Cartesian conservative equation becomes
\begin{equation}
 \begin{aligned}
 J\partial_t\phi
 &+\partial_r\{J(w\phi-D_\phi\partial_r\phi)\}\\
 &+\operatorname{div}_\Gamma\{J[
       \boldsymbol v\phi-D_\phi A_g^{-1}A_g^{-T}
                                  \nabla_\Gamma\phi]\}
 =JS_\phi .
 \end{aligned}
 \label{geo:scalar}
\end{equation}
For the inner $k$ equation use $D_{k,w}=\nu+\nu_{t,A}/\sigma_k$
and $S_{k,w}=P_A-D_w$.  For the outer $k$ and $\eps_o$ equations,
use the chosen model's high-Re coefficients and sources.  The metric
transformation does not add molecular viscosity to those outer
operators or change their constitutive laws.

At a common $k$ trace, the scalar flux per unit ground area is
\begin{equation}
 \boxed{J_H\beta_H\boldsymbol f_k\cdot\boldsymbol N_H
 =m_H\widehat k-J_HD_k\left[
 \partial_rk-\boldsymbol q_H\cdot A_{g,H}^{-T}
                                   \nabla_\Gamma k\right],
 \qquad \boldsymbol f_k=\boldsymbol U k-D_k\nabla k.}
 \label{geo:k-flux}
\end{equation}
Use one shared advective and diffusive flux and zero modeled ground
$k$ flux.  Lateral interfaces between neighbouring wall patches
require the same conservation rule.  At variable $H$, using only
$D_k\partial_rk$ would miss the tangential contribution to the
actual interface-normal diffusion.

If an equilibrium correction is evaluated on a curved ray, its
budget must be metric weighted.  Set
\begin{align}
 \mathcal P_J&=\int_0^HJp_{A,0}\,\dd r,
 &E_{c,m}(r)&=\frac{A_m}{\max(r,r_c)},\nonumber\\
 \int_0^HJE_{c,m}\,\dd r&=\mathcal P_J,
 &\mathcal P_J&\geq
 A_m\left[\frac1{r_m}\int_0^{r_m}J\,\dd r
                    +\int_{r_m}^H\frac Jr\,\dd r\right]
 =:\mathcal B_J .
 \label{geo:metric-cutoff}
\end{align}
The derivative of the integral with respect to $r_c$ is
$-A_mr_c^{-2}\int_0^{r_c}J\,\dd r<0$, and the integral diverges
as $r_c\downarrow0$.  For positive $J$, the displayed inequality
is therefore necessary and sufficient for a unique
$0<r_c\leq r_m$.  The $p_{A,0}$ in this construction remains the
\emph{flat counterfactual} kernel at current $K$.  It is not the
production of an assumed curved neutral equilibrium.

Before positivity adaptation the conservative flux is
\begin{equation}
 j_c(r)=\frac1{J(r)}\int_0^rJ(\zeta)
 [E_{c,m}(\zeta)-p_{A,0}(\zeta)]\,\dd\zeta,
 \qquad
 \nabla\cdot(\chi j_c\boldsymbol n)
       =\chi(E_{c,m}-p_{A,0}),\quad
 j_c(0)=j_c(H)=0 .
 \label{geo:metric-redistribution}
\end{equation}
This uses $\partial_r\chi=0$.  With the positivity adaptation,
replace $\chi j_c$ by
$j_w=J^{-1}\int_0^rJC_w\,\dd\zeta$; its endpoint values also
vanish when the specified weighted budget holds.  This flux adds
no interface term even when $\boldsymbol N_H\ne\boldsymbol n$.
One may assemble $P_A-D_w$, or assemble $P_A-E_w$ and put
$-j_w\boldsymbol n$ in the conservative $k$ flux, but not both.
The flat finite-$\nu$ cutoff in \cref{awf:cutoff} is recovered only
when $J=1$.

The single outer dissipation trace is $\eps_{o,H}=E_w(H)$, including
any attenuation induced by positivity.  There is no separately
solved inner dissipation conservation equation and no independently
prescribed dissipation flux.  Conservation claims for this model
concern mass, momentum, and transported $k$.

For exactly incompressible inner flow, the common momentum and
$k$ equations imply the total-energy identity
\begin{equation}
 \begin{aligned}
 &\partial_t\left(\tfrac12|\boldsymbol U_w|^2+k\right)\\
 &\quad+\nabla\cdot\left[
   (\tfrac12|\boldsymbol U_w|^2+k)\boldsymbol U_w
   -T_w\boldsymbol U_w-D_{k,w}\nabla k-j_w\boldsymbol n
                  \right]\\
 &\qquad=\boldsymbol f_w\cdot\boldsymbol U_w-2\nu S_w:S_w-E_w .
 \end{aligned}
 \label{geo:total-energy}
\end{equation}
The turbulent production cancels between the two equations.  The
molecular loss remains separate, pressure work is contained in
$-T_w\boldsymbol U_w$, and the roughness force contributes its
actual signed work.  Common velocity, traction, and $k$ flux ensure
cancellation of interface energy transfers between the two regions.
Ground no slip removes ground mechanical work.  The same statement
at the discrete level additionally requires compatible contractions,
quadrature, and numerical fluxes.

\section{Spatial transition and exact recovery of the adapted AWF baseline}
\label{geo:transition}

\subsection{Fixed physical support and independent state gates}

For an incoming mean wind in the Cartesian $x$ direction, choose
$x_a<x_b$ on an exactly flat upstream fetch, with $x_b$ before the
terrain variation.  Use the fixed physical mask
\begin{equation}
 g(\zeta)=1-10\zeta^3+15\zeta^4-6\zeta^5,\qquad
 \chi_{\rm sp}(x)=
 \begin{cases}
 1,&x\leq x_a,\\
 g((x-x_a)/(x_b-x_a)),&x_a<x<x_b,\\
 0,&x\geq x_b.
 \end{cases}
 \label{geo:spatial-mask}
\end{equation}
It is $C^2$ and has flat unit and zero plateaus.  In three-dimensional
terrain, prescribe an equally regular function of physical position
on the wall patches.  Keep its physical support and transition length
fixed under refinement.  The policy is one-way: the correction stays
off over the hill, wake, and outlet unless a separate reactivation
rule is deliberately introduced.  If no flat upstream fetch exists,
the flat-reference preservation claim does not apply there.

For a nonnegative sensor $d$, define the compact gate
\begin{equation}
 \mathcal G(d;a,b)=
 \begin{cases}
 1,&d\leq a,\\
 g((d-a)/(b-a)),&a<d<b,\\
 0,&d\geq b,
 \end{cases}
 \qquad a<b .
 \label{geo:gate}
\end{equation}
Extend its first constant branch to negative arguments when it is
used on a signed margin.  At fixed current $K$, the counterfactual
matching velocity is $\boldsymbol U_{H,0}=t_mJ_0(H)\boldsymbol e_x$.
Its outer reference has slope $u_m/\kappa$ and generated roughness
$z_{0,m}(K)$ from \cref{awf:generated-roughness}.  Calibration to a
specified $z_0$ occurs once at $K_0$; a state sensor must not silently
retune $h_r$ or impose that same $z_0$ on every current-$K$ family
member.

As one explicit state-sensing design, let $\mathcal E_a,\mathcal E_b$
evaluate the outer fields at fixed physical heights $z_a,z_b>H$ on
the flat support.  Compare fields with the corresponding reference
using the same evaluation and representation maps; write
$\delta_0\mathcal E f=\mathcal E f-\mathcal E f_0$.
Set $u_s=\sqrt{u_m^2+(\nu/H)^2}$ and use
\begin{equation}
 \begin{aligned}
 \boldsymbol r_U&=P(\widehat{\boldsymbol U}-\boldsymbol U_{H,0})/u_s,
 &\boldsymbol r_{\Delta U}&=
   P(\delta_0\mathcal E_b\boldsymbol U-
     \delta_0\mathcal E_a\boldsymbol U)/u_s,\\
 r_n&=\boldsymbol n\cdot\delta_0\mathcal E_a\boldsymbol U/u_s,
 &r_k&=(\mathcal E_a k-K)/K,\\
 r_\eps&=\delta_0\mathcal E_a\eps/(\mathcal E_a\eps_0),
 &\boldsymbol r_p&=H(\boldsymbol F_p-\boldsymbol F_{p,0})/u_s^2,\\
 &&\boldsymbol r_a&=H(\boldsymbol F_a-\boldsymbol F_{a,0})/u_s^2 .
 \end{aligned}
 \label{geo:state-sensors}
\end{equation}
Here $\boldsymbol F_p$ and $\boldsymbol F_a$ are the pressure/body-force
and convection contributions to \cref{geo:rough-force}, evaluated
separately.  This avoids masking a large departure by cancellation
between them.  The dissipation sample is an independent outer
unknown, away from its constrained trace.  The denominator
$\mathcal E_a\eps_0$ must be positive.  With positive weights,
let $d^2$ be the weighted sum of the seven squared sensor magnitudes
and set $\chi_{\rm st}=\mathcal G(d^2;d_{\rm on}^2,d_{\rm off}^2)$,
where $0<d_{\rm on}<d_{\rm off}$.  All sensors must vanish for the
actual represented target.  This is a comparison of physical
states, not a change to the governing residual.

Support the correction strictly inside its budget domain.  For
example, use the margin
$\mu_J=(\mathcal P_J-\mathcal B_J)/\mathcal B_J$ and define
\begin{equation}
 \chi_{\rm adm}=
 1-\mathcal G(\mu_J;\mu_{\rm off},\mu_{\rm on}),
 \quad 0<\mu_{\rm off}<\mu_{\rm on},\qquad
 \boxed{\chi=\chi_{\rm sp}\chi_{\rm st}\chi_{\rm adm},
 \quad\partial_r\chi=0.}
 \label{geo:correction-gate}
\end{equation}
A separately calibrated wall-regime gate, for example on
$H\sqrt K/\nu$, can be included in the product if the intended
application requires it.  Gate thresholds are model parameters
requiring reporting and sensitivity checks; they are not universal
constants supplied by the outer closures.  The selected reference
must lie on every unit plateau.  The condition $K>0$, geometric
validity, and the real coefficient domain of the outer model are
prior admissibility requirements.  Setting $\chi=0$ cannot repair
an undefined coefficient, intersecting normal rays, or an unstable
outer principal symbol.

If $\chi=0$, evaluate the finite baseline branch directly.  An
inadmissible $r_c$ need not and must not be evaluated merely to
multiply its result by zero.  On active branches, all dependencies
through $K$, $H$, the forcing, the cutoff, and the gates remain.
For any variation $\delta$, in particular,
\begin{equation}
 \begin{aligned}
 \delta E_{\mathrm{pre}}&=(1-\chi)\delta E_A+
 \chi\delta E_{c,m}+(E_{c,m}-E_A)\delta\chi,\\
 \delta S_{k,w}^{\mathrm{pre}}&=
 \delta P_A-(1-\chi)\delta E_A-\chi\delta p_{A,0}
                  +(E_A-p_{A,0})\delta\chi .
 \end{aligned}
 \label{geo:transition-derivatives}
\end{equation}
Gate derivatives follow the ordinary product rule and
$g'(\zeta)=-30\zeta^2(1-\zeta)^2$.  In a state linearization the
fixed spatial mask has zero state derivative, while its spatial
gradient remains.  Differentiating the dependent dissipation or
its trace includes $(E_{c,m}-E_A)\nabla\chi$ and the derivatives
of the positivity adaptation when active.  The momentum equation
acquires no separately invented force proportional to $\nabla\chi$;
its response follows the coupled fields and the stated stress law.

\subsection{What the zero-switch identity includes}

Define the \emph{adapted original-AWF baseline} to use exactly the
same chosen auxiliary coefficient variant, geometry, sample $K$,
spaces, quadrature, no-slip patch, roughness force, pressure coupling,
transported $k$, dependent dissipation, and positivity treatment,
with $\chi=0$ everywhere.  Before positivity adaptation this gives
\begin{equation}
 E_{w,A}=E_A,\qquad D_{w,A}=E_A,\qquad C_{w,A}=0,
 \qquad S_{k,w,A}=P_A-E_A,
 \qquad\eps_{o,H}=E_A(H).
 \label{geo:adapted-baseline}
\end{equation}
With positivity adaptation, use \cref{st:positive-profiles} with
$E_{\mathrm{pre}}=D_{\mathrm{pre}}=E_A$.  More explicitly, if its
weight is $w_{\rm p}$ and cap is $\delta_{\rm p}$, then
\begin{equation}
 \begin{aligned}
 \mathcal B_A&=\int_0^HJE_A\,\dd r,&
 b_{{\rm p},A}&=\mathcal B_A^{-1}\int_0^HJE_Aw_{\rm p}\,\dd r,\\
 m_{{\rm p},A}&=\max(b_{{\rm p},A},\delta_{\rm p}),&
 \lambda_{{\rm p},A}&=b_{{\rm p},A}/m_{{\rm p},A},\\
 D_{w,A}&=E_Aw_{\rm p}/m_{{\rm p},A},&
 E_{w,A}&=\lambda_{{\rm p},A}E_A,\\
 C_{w,A}&=E_{w,A}-D_{w,A},&
 \eps_{o,H}&=E_{w,A}(H).
 \end{aligned}
 \label{geo:positive-baseline}
\end{equation}
Consequently, zero equilibrium correction need not mean zero
positivity redistribution or an unattenuated dissipation trace.
The original baseline profiles are recovered when that additional
adaptation is inactive.

For identical admissible inputs and an identical numerical contract,
the complete wall/outer operator therefore satisfies
\begin{equation}
 \boxed{\mathcal A_{\rm composite}(\boldsymbol Q;\chi=0)
        =\mathcal A_{A}^{\rm adapted}(\boldsymbol Q).}
 \label{geo:baseline-identity}
\end{equation}
This includes the scalar trace and pressure/interface coupling,
not just the $k$ source.  The full-strain curved completion remains
present as part of both operators.  The identity specifies recovery
of this explicit baseline; it does not identify that baseline with
an unspecified historical finite-volume implementation.  Separate
all-baseline and transition calculations can have different fields
because of upstream history and elliptic pressure coupling even
where their local zero-switch operators agree.

Finally, trace transfer must preserve the stated spatial support.
An unconstrained global projection can spread a dissipation
correction onto nominal baseline faces.  A suitable design uses a
baseline trace projection plus a correction projection constrained
to zero on those faces, with a physical transition buffer matching
its support.  Apply the same policy to the actual positivity-adapted
trace, rather than forcing its attenuation to vanish with $\chi$.
The reference trace, positivity requirements, and common flux
moments must all be checked under that chosen transfer.  An
unconstrained polynomial projection of $\chi$ itself need not preserve
its range or its flat plateaus.

\section{Constitutive stability, realizability and regularity}
\label{st:constitutive}

The outer closures admit the neutral reference derived above, but
stationarity does not establish stability. This section assesses the
original outer stress laws of
\cref{out:stress,out:A93Q,out:A95law,out:Elaw} at fixed positive
$k,\eps$. Molecular
viscosity remains absent from the outer momentum operator. The finite
physical $\nu$ retained inside the wall element supplies no additional
outer diffusion.

\subsection{The principal momentum form and the meaning of the bounds}
\label{st:principal}

Write $A_{ij}=\partial_jU_i$, $S=(A+A^T)/2$,
$\Omega=(A-A^T)/2$, and $\tau=k/\eps$. For a gradient variation
$\mathsf H$, let $\delta R[\mathsf H]$ denote the derivative of the
complete covariance with respect to $A$, including every invariant
coefficient. For a solenoidal plane-wave polarization define
\begin{equation}
 \boxed{
 \mathcal D_m(\boldsymbol a,\boldsymbol n)
 =-\mathsf H:\delta R_m[\mathsf H],\qquad
 \mathsf H=\boldsymbol a\otimes\boldsymbol n,\qquad
 \boldsymbol a\cdot\boldsymbol n=0.}
 \label{st:principal-form}
\end{equation}
Here $\boldsymbol n$ is a wave direction. It is unrelated to the wall
normal used in the wall-element construction. For $|\boldsymbol n|=1$,
the pressure-projected acoustic operator is
\[
 \mathsf A_m(\boldsymbol n)\boldsymbol a
 =(I-\boldsymbol n\boldsymbol n^T)
   \{\delta(-R_m)[\boldsymbol a\otimes\boldsymbol n]\boldsymbol n\}.
\]
Its quadratic form is \cref{st:principal-form}. The locally frozen
momentum principal part has Fourier symbol
$-|\boldsymbol\xi|^2\mathsf A_m(\boldsymbol n)$, with
$\boldsymbol n=\boldsymbol\xi/|\boldsymbol\xi|$.

The test is positivity on incompressible rank-one variations. It is
neither positivity of an unrestricted nine-dimensional gradient
Jacobian nor covariance realizability. The latter concerns $R_m$
itself, whereas \cref{st:principal-form} concerns its derivative.
Positive production and positive scalar $\nut$ are different
conditions again.

The three reference roots are
\begin{equation}
 q_{93}=\frac{57+\sqrt{6249}}{40},\qquad
 q_{95}=\frac{3+\sqrt{61}}{2\sqrt2},\qquad
 q_E=\frac{3+\sqrt{41}}{2\sqrt2},
 \qquad C_\mu(q_m)q_m^2=1.
 \label{st:roots}
\end{equation}
The calculations below establish
\begin{equation}
 \boxed{
 \mathcal D_m(\boldsymbol a,\boldsymbol n)
 \ge c_m\nu_{t,o,0}|\boldsymbol a|^2|\boldsymbol n|^2,\qquad
 c_{93}=\frac9{250},\quad c_{95}=\frac1{50},\quad
 c_E=\frac4{q_E^2}=0.36191420548\ldots .}
 \label{st:reference-bounds}
\end{equation}
The A93 and A95 constants are certified conservative bounds, not sharp
minima. The E bound is sharp. These are local constitutive statements:
they do not prove nonlinear, wall-coupled, discrete, or long-time
stability of the full turbulence system.

Throughout the simple-shear calculations let
\[
 B=\boldsymbol e_x\otimes\boldsymbol e_z,\qquad
 A^{(0)}=GB,\qquad G>0,\qquad q=G\tau .
\]
The parameter $q$ labels the base state. It is not held fixed as $A$
varies: $k,\eps$ are fixed, while all gradient invariants vary.
In particular, differentiating only the reduced scalar expression
$C_\mu(q)$ would miss some three-dimensional perturbations.

\subsection{A93: the complete tangent and an all-direction estimate}
\label{st:a93}

Introduce the homogeneous quadratic tensor
\begin{equation}
 \begin{aligned}
 T(A)&=-4\dev\{A^2+(A^2)^T\}
      +13\dev(AA^T)-2\dev(A^TA),\\
 T(B)&=T_0=\diag(28/3,-11/3,-17/3).
 \end{aligned}
 \label{st:a93-polynomial}
\end{equation}
The A93 nonlinear stress is $k\tau^2T(A)/(1000+\eta^3)$, with
$\eta=\sqrt2\,\tau\|S\|_F$ and
$\xi=\sqrt2\,\tau\|\Omega\|_F$. On a regular branch,
\[
 \delta\eta=\sqrt2\,\tau\,\frac{S:\delta S}{\|S\|_F},
 \qquad
 \delta\xi=\sqrt2\,\tau\,\frac{\Omega:\delta\Omega}{\|\Omega\|_F},
 \qquad
 \frac{\delta\nut}{\nut}
 =-\frac{\delta\eta+0.9\delta\xi}{1.25+\eta+0.9\xi}.
\]
At simple shear these become
$\delta\eta=\tau(\mathsf H_{xz}+\mathsf H_{zx})$ and
$\delta\xi=\tau(\mathsf H_{xz}-\mathsf H_{zx})$.

For arbitrary positive simple-shear $q$, define
\begin{equation}
 \begin{aligned}
 v_1&=\frac{1.9q}{1.25+1.9q},&
 v_2&=\frac{0.1q}{1.25+1.9q},\\
 \beta_{93}&=\frac{q}{C_\mu(1000+q^3)},&
 \gamma_{93}&=\frac{3q^3}{1000+q^3},
 \qquad C_\mu=\frac{2/3}{1.25+1.9q}.
 \end{aligned}
 \label{st:a93-coefficients}
\end{equation}
The full derivative is
\begin{equation}
 \begin{aligned}
 \frac{\delta(-R_{93})[\mathsf H]}{\nut}
 ={}&\mathsf H+\mathsf H^T
 -(v_1\mathsf H_{xz}+v_2\mathsf H_{zx})(B+B^T)\\
 &-\beta_{93}\delta T_B[\mathsf H]
 +\beta_{93}\gamma_{93}
       (\mathsf H_{xz}+\mathsf H_{zx})T_0 ,
 \end{aligned}
 \label{st:a93-tangent}
\end{equation}
where
\begin{equation}
 \begin{aligned}
 \delta T_B[\mathsf H]={}&
 -4\dev\{B\mathsf H+\mathsf H B+
                   (B\mathsf H+\mathsf H B)^T\}\\
 &+13\dev(\mathsf H B^T+B\mathsf H^T)
 -2\dev(\mathsf H^TB+B^T\mathsf H).
 \end{aligned}
 \label{st:a93-polynomial-derivative}
\end{equation}
The last term in \cref{st:a93-tangent} differentiates the nonlinear
denominator. Omitting it would change the principal form.

For unit orthogonal $\boldsymbol a,\boldsymbol n$, set
$X=a_xn_z$ and $Y=a_zn_x$. Direct contraction gives
\[
 \mathsf H:\delta T_B=-6a_xa_z+9n_xn_z,\qquad
 \mathsf H:T_0=13a_xn_x-2a_zn_z.
\]
Consequently,
\begin{equation}
 \begin{aligned}
 \frac{\mathcal D_{93}}{\nut}
 ={}&1-(v_1X+v_2Y)(X+Y)
       +6\beta_{93}a_xa_z-9\beta_{93}n_xn_z\\
 &+\beta_{93}\gamma_{93}(X+Y)
                      (13a_xn_x-2a_zn_z).
 \end{aligned}
 \label{st:a93-rank-one}
\end{equation}

Here is an explicit bound covering every such pair. Put
$r_0=|a_xn_z|$ and $t=(1-r_0)^{1/2}$. Cauchy--Schwarz yields
$|X|+|Y|\le1$, whence
\[
 (v_1X+v_2Y)(X+Y)\le v_2+(v_1-v_2)r_0.
\]
Orthonormality also gives
$a_z^2,n_x^2\le1-\max(a_x^2,n_z^2)\le1-r_0$.
Thus $|a_xa_z|,|n_xn_z|,|a_xn_x|,|a_zn_z|\le t$, and
\begin{equation}
 \begin{aligned}
 \frac{\mathcal D_{93}}{\nut}
 &\ge1-v_1+(v_1-v_2)t^2
              -15\beta_{93}(1+\gamma_{93})t\\
 &\ge L_{93}(q):=
 1-v_1-
 \frac{\{15\beta_{93}(1+\gamma_{93})\}^2}{4(v_1-v_2)}.
 \end{aligned}
 \label{st:a93-q-bound}
\end{equation}
This sufficient, possibly conservative, bound applies at any positive
simple-shear $q$ for which $L_{93}(q)>0$.

On the canonical root,
\[
 v_1=\frac{2.85}{q_{93}},\quad v_2=\frac{0.15}{q_{93}},\quad
 \beta_{93}=\frac{q_{93}^3}{1000+q_{93}^3},\quad
 \gamma_{93}=3\beta_{93}.
\]
Hence
\begin{equation}
 L_{93}(q_{93})
 =1-v_1-\frac{(15\beta_{93}+45\beta_{93}^2)^2}
                    {4(v_1-v_2)}
 =0.03613783593\ldots>\frac9{250}.
 \label{st:a93-reference-bound}
\end{equation}
The exact algebraic root and this strict inequality are enclosed by
the outward-rounded arithmetic in \cref{st:certificate}. Homogeneity
then proves the A93 part of \cref{st:reference-bounds} for arbitrary
$\boldsymbol a,\boldsymbol n$.

\subsection{Invariant derivatives for A95 and E94/E95}
\label{st:invariant-derivatives}

For these models write
\[
 s=\|S\|_F,\quad r=\|\Omega\|_F,\quad
 \mathcal U=(s^2+r^2)^{1/2},\quad
 A_s=\frac{3\lambda_{\max}(S)}s,\quad
 C=\frac1{a_0+\tau A_s\mathcal U},
\]
where $a_0=6.5$ for A95 and $a_0=4$ for E94/E95.
The eigenvalue identity is valid for traceless symmetric $S\ne0$;
$\lambda_{\max}$ is the largest signed eigenvalue, not the one of
largest absolute magnitude.

At simple shear the largest eigenvalue is $G/2$, with unit
eigenvector $(\boldsymbol e_x+\boldsymbol e_z)/\sqrt2$. For a
trace-free $\mathsf H$,
\begin{equation}
 \begin{aligned}
 \delta s&=(\mathsf H_{xz}+\mathsf H_{zx})/\sqrt2,&
 \delta r&=(\mathsf H_{xz}-\mathsf H_{zx})/\sqrt2,\\
 \delta\mathcal U&=\mathsf H_{xz},&
 \delta\lambda_{\max}
  &=\tfrac12(-\mathsf H_{yy}+\mathsf H_{xz}+\mathsf H_{zx}),\\
 \delta(A_s\mathcal U)&=\frac3{\sqrt2}
                    (\mathsf H_{xz}-\mathsf H_{yy}).
 \end{aligned}
 \label{st:invariant-variations}
\end{equation}
In particular, a normal-strain perturbation changes the invariant
coefficient even when $\mathsf H_{xz}=0$.

Define the derivative coefficient
\begin{equation}
 \alpha_c=\frac3{\sqrt2}Cq,\qquad
 C=\frac1{a_0+(3/\sqrt2)q},\qquad
 \frac{\delta\nut}{\nut}
 =-\frac{\alpha_c}{G}(\mathsf H_{xz}-\mathsf H_{yy}).
 \label{st:viscosity-derivative}
\end{equation}
The symbol $\alpha_c$ here is unrelated to the auxiliary wall slope
$\alpha$.

For A95 also put
\begin{equation}
 \begin{gathered}
 f=1-\frac92C^2q^2,\qquad
 C_Q=\frac{\sqrt f}{1+3q^2},\qquad
 \beta_{95}=\frac{C_Qq}{C},\\
 c=\frac{1-f}{f},\qquad e=\frac{6q^2}{1+3q^2},\qquad
 \ell_p=c(1-\alpha_c)+e .
 \end{gathered}
 \label{st:a95-coefficients}
\end{equation}
The general differential
\[
 \delta\log C_Q
 =\frac{\delta f}{2f}
 -\frac{6\tau^2(r\,\delta s+s\,\delta r)}
             {1+6\tau^2sr},\qquad
 \delta f=-18\tau^2(Cs^2\,\delta C+C^2s\,\delta s)
\]
then gives
\begin{equation}
 \boxed{\frac{\delta C_Q}{C_Q}
 =-\frac1G\{\ell_p\mathsf H_{xz}
              +c\mathsf H_{zx}+c\alpha_c\mathsf H_{yy}\}.}
 \label{st:a95-cq-derivative}
\end{equation}
All square-root, strain, rotation, and $C_\mu$ derivatives are
included. At simple shear $f>0$ for every finite $q>0$; that
observation does not extend the A95 real domain to arbitrary
gradients.

\subsection{E94/E95: a sharp simple-shear bound}
\label{st:e-bound}

Since $Q=0$, \cref{st:viscosity-derivative} gives, for incompressible
rank-one $\mathsf H$,
\begin{equation}
 \boxed{\frac{\mathcal D_E(\mathsf H)}{\nut}
 =\|\mathsf H\|_F^2
 -\alpha_c(\mathsf H_{xz}-\mathsf H_{yy})
              (\mathsf H_{xz}+\mathsf H_{zx}).}
 \label{st:e-rank-one}
\end{equation}
Each bilinear form
$a_xn_z-a_yn_y$ and $a_xn_z+a_zn_x$ has absolute value at most
$|\boldsymbol a||\boldsymbol n|$: its defining matrix has spectral
norm one. It follows that
\begin{equation}
 \mathcal D_E\ge
 \frac{4}{4+(3/\sqrt2)q}\,
 \nut|\boldsymbol a|^2|\boldsymbol n|^2
 \qquad(q>0,\ \hbox{simple shear}).
 \label{st:e-q-bound}
\end{equation}
Equality is achieved by
$\mathsf H=\boldsymbol e_x\otimes\boldsymbol e_z$.
On the canonical root,
$q_E^2=4+(3/\sqrt2)q_E$, which proves the sharp E constant in
\cref{st:reference-bounds}. This all-$q$ simple-shear result does
not cover other classes of velocity gradient.

\subsection{A95: the full tangent and a certified rank-one bound}
\label{st:a95-bound}

Let $S_B=(B+B^T)/2$, $\Omega_B=(B-B^T)/2$,
$\delta S=(\mathsf H+\mathsf H^T)/2$,
$\delta\Omega=(\mathsf H-\mathsf H^T)/2$, and
$[X,Y]=XY-YX$. With $D_0=\diag(1,0,-1)$, the full A95 tangent is
\begin{equation}
 \begin{aligned}
 \frac{\delta(-R_{95})[\mathsf H]}{\nut}
 ={}&\mathsf H+\mathsf H^T
 -\alpha_c(\mathsf H_{xz}-\mathsf H_{yy})(B+B^T)\\
 &-2\beta_{95}\{[\delta\Omega,S_B]+[\Omega_B,\delta S]\}\\
 &+\beta_{95}\{\ell_p\mathsf H_{xz}
              +c\mathsf H_{zx}+c\alpha_c\mathsf H_{yy}\}D_0 .
 \end{aligned}
 \label{st:a95-tangent}
\end{equation}
The last line is the coefficient derivative of the nonlinear stress;
it is not present in a frozen-coefficient calculation.

Contracting with incompressible rank-one $\mathsf H$ gives
\begin{equation}
 \begin{aligned}
 F(\mathsf H):=\frac{\mathcal D_{95}}{\nut}
 ={}&\|\mathsf H\|_F^2
 -\alpha_c(\mathsf H_{xz}-\mathsf H_{yy})
                       (\mathsf H_{xz}+\mathsf H_{zx})\\
 &+\beta_{95}\{(\mathsf H\mathsf H^T)_{xz}
                      -(\mathsf H^T\mathsf H)_{xz}\}\\
 &+\beta_{95}\{\ell_p\mathsf H_{xz}
              +c\mathsf H_{zx}+c\alpha_c\mathsf H_{yy}\}
                    (\mathsf H_{xx}-\mathsf H_{zz}).
 \end{aligned}
 \label{st:a95-rank-one}
\end{equation}
At $q=q_{95}$ the coefficients are
\[
 \alpha_c=0.555028808758460\ldots,\quad
 \beta_{95}=1.036109942594599\ldots,\quad
 f=0.691943021448165\ldots,
\]
\[
 c=0.445205701919075\ldots,\qquad
 e=1.955380111911424\ldots,\qquad
 \ell_p=2.153483823441880\ldots .
\]

Extend the polynomial on the right of \cref{st:a95-rank-one} to
arbitrary trace-free matrices, solely for the following algebraic
certificate. This extension need not equal the physical tangent
contraction on a non-rank-one gradient. Form
\begin{equation}
 \begin{aligned}
 \mathcal P(\mathsf H)={}&F(\mathsf H)
 +\frac15(\mathsf H_{xy}\mathsf H_{yz}
                  -\mathsf H_{xz}\mathsf H_{yy})\\
 &-\frac85(\mathsf H_{xx}\mathsf H_{zz}
                  -\mathsf H_{xz}\mathsf H_{zx})
 -\frac1{50}\|\mathsf H\|_F^2.
 \end{aligned}
 \label{st:a95-certificate-form}
\end{equation}
The two added minors vanish for every rank-one matrix. Eliminate
$\mathsf H_{zz}=-\mathsf H_{xx}-\mathsf H_{yy}$ and use the eight
coordinates
\[
 (\mathsf H_{xx},\mathsf H_{xy},\mathsf H_{xz},\mathsf H_{yx},
   \mathsf H_{yy},\mathsf H_{yz},\mathsf H_{zx},\mathsf H_{zy}).
\]
The symmetric matrix of $\mathcal P$ has a positive $LDL^T$
factorization. Its pivots, rounded only for display, are
\begin{equation}
 \begin{aligned}
 (&3.56,\ 0.98,\ 0.023749225694546,\ 0.98,\\
 &0.698569139712622,\ 0.695937802769545,\
   0.321159138225151,\ 0.702126481651832).
 \end{aligned}
 \label{st:a95-pivots}
\end{equation}
The executable 60-digit interval calculation in
\cref{st:certificate} constructs this matrix directly from
\cref{st:a95-certificate-form} and verifies that every pivot
interval has a strictly positive lower endpoint. It follows that
$\mathcal P\ge0$ on all trace-free matrices. On rank-one matrices
the minors vanish, proving
$\mathcal D_{95}\ge\nut\|\mathsf H\|_F^2/50$.
The proof covers all incompressible wave directions and polarizations;
a directional search is not substituted for it.

The neutral coefficients are smooth and the unit orthogonal-direction
set is compact. The strict bounds for all three closures therefore
persist in some open neighborhood of each reference state.
No certified multidimensional neighborhood of a prescribed size is
claimed.

\subsection{A transverse sector and its limitations}
\label{st:transverse}

For $\delta U_y=v(x,z,t)$ with $\partial_y=0$, the first variations
of the invariant coefficients and production vanish at simple shear.
The complete transverse responses are
\begin{equation}
 \begin{aligned}
 \begin{pmatrix}-\delta R_{yx}\\-\delta R_{yz}\end{pmatrix}_{93}
 &=\nut\begin{pmatrix}1&-13\beta_{93}\\4\beta_{93}&1\end{pmatrix}
          \begin{pmatrix}v_x\\v_z\end{pmatrix},\\
 \begin{pmatrix}-\delta R_{yx}\\-\delta R_{yz}\end{pmatrix}_{95}
 &=\nut\begin{pmatrix}1&-\beta_{95}\\0&1\end{pmatrix}
          \begin{pmatrix}v_x\\v_z\end{pmatrix},\\
 \begin{pmatrix}-\delta R_{yx}\\-\delta R_{yz}\end{pmatrix}_{E}
 &=\nut\begin{pmatrix}1&0\\0&1\end{pmatrix}
          \begin{pmatrix}v_x\\v_z\end{pmatrix}.
 \end{aligned}
 \label{st:transverse-blocks}
\end{equation}
Their minimum transverse principal coefficients are
$1-\tfrac92\beta_{93}$, $1-\tfrac12\beta_{95}$, and $1$,
respectively, in units of $\nut$. At the neutral roots they equal
$0.82963783890\ldots$, $0.48194502870\ldots$, and $1$.
These checks are useful but do not prove the stronger all-direction
bounds above.

\subsection{Explicit negative principal eigenvalues away from reference}
\label{st:failures}

The following algebraic states use positive $k,\eps$ and the original
stress laws. They are not additional stationary neutral references.
Each has positive covariance and positive production. Thus neither
condition detects the constitutive failure.

\subsubsection{A93: a realizable simple-shear state at \texorpdfstring{$q=8$}{q=8}}

Take $k=\eps=1$ and $A=8\boldsymbol e_x\otimes\boldsymbol e_z$.
Then
\[
 C_\mu=\frac{40}{987},\quad
 v_1=\frac{304}{329},\quad v_2=\frac{16}{329},\quad
 \beta_{93}=\frac{47}{360},\quad\gamma_{93}=\frac{64}{63}.
\]
For
$\boldsymbol n=(2,0,1)/\sqrt5$ and
$\boldsymbol a=(1,0,-2)/\sqrt5$,
\begin{equation}
 \frac{\mathcal D_{93}}{\nut}
 =1+\frac3{25}(v_1-4v_2)-6\beta_{93}
             -\frac{18}{5}\beta_{93}\gamma_{93}
 =-\frac{51301}{296100}<0.
 \label{st:a93-negative}
\end{equation}
Reflection in $y$ decouples the two solenoidal polarizations.
In the basis $\{\boldsymbol a,\boldsymbol e_y\}$ the projected
acoustic matrix is
\[
 \nut\,\diag\left(-\frac{51301}{296100},\frac{53}{100}\right).
\]
The negative value is an actual eigenvalue, not merely a negative
quadratic-form value of a nonnormal matrix.

At the same state,
\begin{equation}
 \frac Rk=
 \begin{pmatrix}
 86/81&0&-320/987\\
 0&290/567&0\\
 -320/987&0&242/567
 \end{pmatrix}\succ0,\qquad P_k=\nut G^2>0.
 \label{st:a93-failure-covariance}
\end{equation}
The $xz$ principal determinant is
$247166356/710169201>0$; the covariance eigenvalues are approximately
$0.29051012,0.51146384,1.19802604$.

\subsubsection{A95: a decoupled transverse failure at \texorpdfstring{$q=1$}{q=1}}

For positive simple-shear $q$,
\begin{equation}
 \beta_{95}(q)
 =\frac{q\sqrt{6.5^2+2(6.5)(3/\sqrt2)q}}{1+3q^2}.
 \label{st:a95-beta-q}
\end{equation}
At $G=k=\eps=1$,
\[
 C_\mu=0.115991514077890\ldots,\quad
 f=0.939456858978634\ldots>0,\quad
 \beta_{95}=2.089066245752444\ldots .
\]
Choose $\boldsymbol a=\boldsymbol e_y$ and
$\boldsymbol n=(\boldsymbol e_x+\boldsymbol e_z)/\sqrt2$.
The transverse subsystem is decoupled, and
\begin{equation}
 \boxed{\frac{\mathcal D_{95}}{\nut}
       =1-\frac{\beta_{95}}2
       =-0.044533122876222\ldots<0.}
 \label{st:a95-negative}
\end{equation}
Pressure projection does not remove this polarization.
Moreover $\delta k=\delta\eps=0$ is consistent with this linearized
subsystem because the first variations of $C_\mu,C_Q,P_k$ vanish.
The covariance eigenvalues are approximately
$0.39802176,0.66666667,0.93531157$, and $P_k=C_\mu>0$.
The failure is inside the real coefficient domain, away from its
square-root boundary.

\subsubsection{E94/E95: a pure-strain failure}

Take
\[
 k=\eps=1,\qquad A=S=\diag(10,9,-19),\qquad\Omega=0.
\]
The largest strain eigenvalue is the distinct value $10$.
Since $\mathcal U=s$,
$A_s\mathcal U=30$ and $\nut=C_\mu=1/34$.
For
$\boldsymbol n=(1,0,1)/\sqrt2$ and
$\boldsymbol a=(1,0,-1)/\sqrt2$,
\[
 \delta\lambda_{\max}=\tfrac12,\qquad
 S:\mathsf H=\tfrac{29}{2},\qquad
 \delta(A_s\mathcal U)=\tfrac32,\qquad
 2\delta S:\mathsf H=1.
\]
It follows that
\begin{equation}
 \boxed{\frac{\mathcal D_E}{\nut}
 =1-2C_\mu\,\delta(A_s\mathcal U)(S:\mathsf H)
 =1-\frac{87}{68}=-\frac{19}{68}<0.}
 \label{st:e-negative}
\end{equation}
The projected acoustic matrix in the basis
$\{\boldsymbol a,\boldsymbol e_y\}$ is
$\nut\diag(-19/68,1)$. Yet
\begin{equation}
 \boxed{R/k=\diag(4/51,7/51,91/51)\succ0,\qquad
        P_k=2\nut(S:S)=1084/34>0.}
 \label{st:e-failure-covariance}
\end{equation}
A realizable scalar eddy-viscosity value therefore does not ensure
a monotone constitutive mapping.

\subsection{Covariance realizability and production}
\label{st:realizability}

\subsubsection{A93: unrestricted gradients can violate covariance positivity}

For pure rotation in the $xz$ plane,
\[
 S=0,\qquad
 \Omega=\begin{pmatrix}0&0&\omega\\0&0&0\\-\omega&0&0\end{pmatrix},
\]
the printed A93 law yields
\begin{equation}
 Q=-\frac{19k\tau^2}{1000}\dev(\Omega^2),\qquad
 \frac{R_{yy}}k=\frac23-\frac{38(\tau\omega)^2}{3000}.
 \label{st:a93-rotation}
\end{equation}
Thus $R_{yy}<0$ if $\tau|\omega|>\sqrt{1000/19}$; at
$\tau|\omega|=8$ it is $-0.144k$.

A93 production can also be negative at positive covariance. Set
$k=\eps=1$, $S=\diag(1,1,-2)$, and
$\Omega_{xy}=-\Omega_{yx}=5$, with the other entries zero.
Writing
\[
 C=\frac{2/3}{41/4+2\sqrt3},\qquad D=1000+24\sqrt3,
\]
gives
\begin{equation}
 \begin{aligned}
 R&=\diag\left(\frac23-2C+\frac{466}{3D},
              \frac23-2C+\frac{466}{3D},
              \frac23+4C-\frac{932}{3D}\right)\\
  &\simeq\diag(0.7185770907,0.7185770907,0.5628458186)\succ0,\\
 P_k&=12C-\frac{932}{D}=-0.3114625443\ldots<0.
 \end{aligned}
 \label{st:a93-negative-production}
\end{equation}
Negative production represents modeled transfer from turbulence to
the mean; it is not alone an impossible covariance state. It rules
out arguments that assume the full A93 production is everywhere
nonnegative. Clipping this contraction changes the model and its
mean/turbulent energy-transfer identity.

\subsubsection{A95: positive covariance on its real coefficient domain}

Let $T=\tau S$, $V=\tau\Omega$,
$s_T=\|T\|_F$, $r_V=\|V\|_F$, and $x=C_\mu^2s_T^2$.
The radicand requires $0\le x\le1/9$. Write
\[
 R/k=\tfrac23I+B_R,\qquad
 B_R=-2C_\mu T+2C_Q[V,T].
\]
Both summands are symmetric and traceless, and
$T:[V,T]=0$. Diagonalizing $T$ proves
$\|[V,T]\|_F^2\le2s_T^2r_V^2$, since
$(t_i-t_j)^2\le2\sum_\ell t_\ell^2$.
Consequently,
\begin{equation}
 \begin{aligned}
 \|B_R\|_F^2
 &\le4x+\frac{8(1-9x)(s_Tr_V)^2}{(1+6s_Tr_V)^2}\\
 &\le\frac29+2x\le\frac49.
 \end{aligned}
 \label{st:a95-covariance-norm}
\end{equation}
For a symmetric traceless three-dimensional tensor,
$\lambda_{\min}(B_R)\ge-\sqrt{2/3}\|B_R\|_F$. Therefore
\begin{equation}
 \boxed{\lambda_{\min}(R/k)
 \ge\frac23-\sqrt{\frac8{27}}
 =0.122335612714849\ldots>0.}
 \label{st:a95-covariance-bound}
\end{equation}
This controls the full covariance, including its determinant.
It does not establish ellipticity on that domain.

The real domain excludes some incompressible gradients. For
$S=\diag(L,0.9L,-1.9L)$, $\Omega=0$, and $L>0$, the square root
is real only if
\begin{equation}
 \tau L\le\frac{6.5}{3(\sqrt{5.42}-1)}
          \simeq1.631.
 \label{st:a95-real-domain}
\end{equation}
A zero radicand generally produces a singular derivative; beyond it
the printed coefficient is nonreal. An implementation must declare
a restricted state domain or an explicitly assessed continuation.

\subsubsection{E94/E95: covariance realizability does not imply ellipticity}

For $s>0$, $\mathcal U\ge s$ and the largest-eigenvalue identity imply
\[
 C_\mu\tau\lambda_{\max}(S)
 =\frac{\tau\lambda_{\max}(S)}
        {4+3\tau\lambda_{\max}(S)\mathcal U/s}<\frac13.
\]
Hence
\begin{equation}
 \boxed{\lambda_{\min}(R)
 =\frac23k-2C_\mu k\tau\lambda_{\max}(S)>0}
 \label{st:e-covariance}
\end{equation}
at finite positive $k,\eps$ wherever the coefficient is defined.
The stress has the isotropic extension $R=2kI/3$ at $S=0$.
The explicit failure in \cref{st:e-negative} lies on a smooth branch
of this realizable mapping.

For A95, $S:(\Omega S-S\Omega)=0$. Thus
\begin{equation}
 P_k^{95}=2\nut S:S\ge0
 \quad\hbox{on its real domain},\qquad
 P_k^E=2\nut S:S\ge0 .
 \label{st:positive-production}
\end{equation}
These production identities do not test the stress derivative.

\subsection{Regularity, parameter changes and the limits of the claims}
\label{st:scope}

The neutral states have nonzero strain and rotation, distinct strain
eigenvalues, and a strictly positive A95 radicand. None of the
reference proofs selects an arbitrary derivative at a degeneracy.

Away from these states, the invariant formula for A95/E is undefined
at $s=0$, and the largest-eigenvalue coefficient can be nonsmooth
when the largest eigenvalue is repeated. An isotropic stress
extension at $s=0$ does not itself assign a unique scalar viscosity
for turbulence diffusion when rotation is nonzero. The chosen
coefficient continuation must also be stated. A93 norm coefficients
require care at vanishing strain or rotation. Some complete stress
products can remain differentiable when a factor is not; the
required object is the complete constitutive mapping.

Choosing another $\kappa$ and recalibrating $\sigma_\eps$ leaves
$q_m$ and the original stress function unchanged. The dimensionless
constants in \cref{st:reference-bounds} therefore remain the same,
while the dimensional reference scale
$\nu_{t,o,0}=\kappa\ust z$ changes. An optional recalibration of the
inner wall slope $\alpha$, or of the auxiliary wall constants,
also leaves these outer certificates unchanged. It changes the
wall operator and must satisfy its own resistance, cutoff,
boundary, and discrete checks. A constitutive modification of
the outer stress would require a new derivative assessment.

A negative projected principal eigenvalue gives a locally frozen
growth rate proportional to $+|\boldsymbol\xi|^2$. The examples
above supply actual eigenvalues and admissible polarizations;
they are stronger than negative numerical-range observations.
A wall-only redistribution cannot remove a failure supported in
the outer region. An exactly compensated implicit Laplacian,
under-relaxation, a smaller time step, or a converged steady
iteration does not change the principal part of the specified
continuum stress law.

One may restrict the admissible state envelope around a certified
reference or introduce a constitutive regularization. An active
regularization changes the model and must be used consistently in
momentum, production and traction. No universally stable extension
over all separated-flow gradients is established here.

\section{Positivity of transported turbulence quantities and its remedy}
\label{st:positivity}

A stationary reference, positive covariance, nonnegative production,
and positive scalar diffusivity do not individually ensure that
transported $k$ and $\eps$ remain positive. The inner sampled-energy
source creates an additional issue, common to all three outer
closures.

\subsection{The sampled-energy obstruction}
\label{st:sampled-obstruction}

Recall that $K=\mathcal K[k_h]$ is a representative positive sample,
not the local transported field. At $\chi=0$ the unmodified wall
source is $P_A-E_A(K)$, where $E_A>0$. Take a locally flat zero
minimum of $k_h$ away from its positive sample $K$, with zero
current velocity and production. Advection and diffusion vanish
at this minimum, but
\begin{equation}
 \partial_tk=-E_A(K)<0.
 \label{st:zero-minimum}
\end{equation}
The finite molecular diffusion also vanishes at a locally flat
minimum and does not remove this counterexample. A wholly
molecular sublayer with $\nu_{t,A}=P_A=0$ makes the same source
obstruction explicit.

The proposed remedy below modifies the destruction and dependent
dissipation. It leaves an open neighborhood of the neutral
reference unchanged. It is a separate modeling choice; positivity
is not attributed to the original sampled-$K$ source.

\subsection{Metric budgets and a capped destruction rate}
\label{st:positive-construction}

On one admissible wall ray let $r\in[0,H]$ and let $J(r)>0$ be the
geometric weight; $J=1$ on a flat ray. For clarity, the profiles
from \cref{awf:pre-profiles,awf:net-source} are
\begin{equation}
 \begin{aligned}
 E_{\mathrm{pre}}&=(1-\chi)E_A+\chi E_{c,m},\\
 D_{\mathrm{pre}}&=(1-\chi)E_A+\chi p_{A,0},\\
 C_{\mathrm{pre}}&=E_{\mathrm{pre}}-D_{\mathrm{pre}},\qquad
 S_{k,w}^{\mathrm{pre}}=P_A-D_{\mathrm{pre}}.
 \end{aligned}
 \label{st:pre-profiles}
\end{equation}
The weight $\chi\in[0,1]$ is constant along each normal ray.
The corrected profiles satisfy their metric budget when active,
so
\begin{equation}
 \mathcal B:=\int_0^HJ E_{\mathrm{pre}}\,\dd r
            =\int_0^HJ D_{\mathrm{pre}}\,\dd r>0.
 \label{st:pre-budget}
\end{equation}
At $\chi=0$ evaluate the baseline directly; an inadmissible
corrected profile must not be evaluated and then multiplied by zero.

Choose fixed $0<\theta_{\mathrm p}<1$ and
$0<\delta_{\mathrm p}<1$. For nonnegative local $k$, set
\begin{equation}
 \begin{aligned}
 w_{\mathrm p}(r)&=\min\left(1,\frac{k(r)}{\theta_{\mathrm p}K}\right),\\
 b_{\mathrm p}&=\mathcal B^{-1}
        \int_0^HJ D_{\mathrm{pre}}w_{\mathrm p}\,\dd r,\qquad
 m_{\mathrm p}=\max(b_{\mathrm p},\delta_{\mathrm p}),\qquad
 \lambda_{\mathrm p}=\frac{b_{\mathrm p}}{m_{\mathrm p}} .
 \end{aligned}
 \label{st:normalization}
\end{equation}
Then $0\le b_{\mathrm p}\le1$,
$\delta_{\mathrm p}\le m_{\mathrm p}\le1$, and
$0\le\lambda_{\mathrm p}\le1$. Define
\begin{equation}
 \boxed{
 \begin{aligned}
 D_{\mathrm{pos}}&=\frac{D_{\mathrm{pre}}w_{\mathrm p}}{m_{\mathrm p}},
 &E_{\mathrm{pos}}&=\lambda_{\mathrm p}E_{\mathrm{pre}},\\
 C_{\mathrm{pos}}&=E_{\mathrm{pos}}-D_{\mathrm{pos}},
 &S_{k,w}^{\mathrm{pos}}&=P_A-D_{\mathrm{pos}} .
 \end{aligned}}
 \label{st:positive-profiles}
\end{equation}
This total redistribution replaces $C_{\mathrm{pre}}$.
It is not an additional correction to the already modified source.

The exact identities are
\begin{equation}
 \begin{aligned}
 \int_0^HJ D_{\mathrm{pos}}\,\dd r
 &=\int_0^HJ E_{\mathrm{pos}}\,\dd r
   =\lambda_{\mathrm p}\mathcal B,\\
 \int_0^HJ C_{\mathrm{pos}}\,\dd r&=0,\\
 0\le\frac{D_{\mathrm{pos}}}{k}
 &\le\frac{D_{\mathrm{pre}}}
            {\delta_{\mathrm p}\theta_{\mathrm p}K}
       \qquad(k>0).
 \end{aligned}
 \label{st:positive-budget}
\end{equation}
The last inequality follows from
$w_{\mathrm p}\le k/(\theta_{\mathrm p}K)$ and
$m_{\mathrm p}\ge\delta_{\mathrm p}$.
At a local zero, destruction vanishes exactly.

Three regimes must be distinguished. On the all-ray plateau
$k/K>\theta_{\mathrm p}$, one has
$w_{\mathrm p}=b_{\mathrm p}=m_{\mathrm p}
=\lambda_{\mathrm p}=1$; the model and all existing local
derivatives are unchanged. This includes $k=K$ on the reference.
For $\delta_{\mathrm p}\le b_{\mathrm p}<1$, the loss is redistributed
while retaining the total budget $\mathcal B$.
For $b_{\mathrm p}<\delta_{\mathrm p}$, the total loss is attenuated
to $\lambda_{\mathrm p}\mathcal B$ in order to keep the destruction
rate bounded.

That last concession is necessary. If a nonnegative destruction
satisfies $D\le Lk$ with bounded $L$, then
$\int JD\le\|L\|_\infty\int Jk\to0$ as the available integrated energy
vanishes. It cannot retain a fixed positive total loss in that
limit. The molecular mean-flow dissipation is not part of
$\mathcal B$ and is not rescaled.

\subsection{The conditional continuum positivity statement}
\label{st:positive-assumptions}

The following assumptions delimit what the remedy proves:
\begin{enumerate}
 \item The normal map is admissible, $J>0$, and all wall inputs and
       profiles are regular and bounded on the time interval considered.
       The sampler is constant-preserving and positive on nonnegative
       fields, with $K\ge K_{\min}>0$; the normalization
       $\mathcal B>0$ remains defined. For uniform comparison and
       regularity estimates, $J$ and $\mathcal B$ also have positive
       lower bounds and the normalized inputs remain bounded.
 \item The inner scalar diffusivity is
       $D_{k,w}=\nu+\nu_{t,A}/\sigma_k>0$, with the finite physical
       $\nu>0$. The advecting velocity and coefficients are sufficiently
       regular for the scalar comparison argument.
 \item Production is the consistent incompressible contraction
       $P_A=2\nu_{t,A}S_w:S_w\ge0$. A formulation retaining discrete
       divergence errors must also retain the isotropic contribution
       $-(2k/3)\nabla\cdot\boldsymbol U_w$ and account for it consistently.
 \item Initial and inflow data and the supplied interface trace of
       $k$ are nonnegative. The wall has zero modeled $k$ flux, and
       the remaining boundary/transmission conditions do not inject
       negative values.
\end{enumerate}

Under these assumptions the inner reaction is quasi-positive:
\[
 k=0\quad\Longrightarrow\quad
 P_A-D_{\mathrm{pos}}=P_A\ge0.
\]
More explicitly, for $k>0$ write
$D_{\mathrm{pos}}=a(r,t)k$ with
$0\le a\le D_{\mathrm{pre}}/
(\delta_{\mathrm p}\theta_{\mathrm p}K)$.
A regular scalar transport--diffusion equation with this bounded
nonnegative destruction coefficient, nonnegative production and
compatible data preserves nonnegativity by the comparison principle.
The nonlocal sampler and normalization do not change the sign at
a first zero while their admissibility assumptions hold.

This is a conditional inner result, not a proof for the coupled
outer/inner system. In particular, nonnegative interface data are
an assumption on the coupled solution, not permission to impose
an extra trace independently of the common-flux coupling.
Vanishing $K$, an undefined normalization, or a negative trial field
lies outside the stated construction and requires an explicitly
qualified continuation or an invalid-state outcome. A hidden floor
is not part of this proof.
The numerical no-slip condition $\boldsymbol U_w(0)=0$ remains in
force; it does not introduce a wall-resolved condition $k(0)=0$.

\subsection{Conservative redistribution and interface dissipation}
\label{st:positive-interface}

Define the normal redistribution flux
\begin{equation}
 j_{\mathrm{pos}}(r)=\frac1{J(r)}
          \int_0^rJ(\zeta)C_{\mathrm{pos}}(\zeta)\,\dd\zeta,
 \qquad
 j_{\mathrm{pos}}(0)=j_{\mathrm{pos}}(H)=0.
 \label{st:positive-flux}
\end{equation}
Then $J^{-1}\partial_r(Jj_{\mathrm{pos}})=C_{\mathrm{pos}}$.
The retained $k$ equation may be assembled with source
$P_A-D_{\mathrm{pos}}$, or with source $P_A-E_{\mathrm{pos}}$
and conservative flux
$\boldsymbol U_wk-D_{k,w}\nabla k-j_{\mathrm{pos}}\boldsymbol n$.
These are equivalent formulations, not two contributions to add.

For every wall test function $\phi$,
\begin{equation}
 \int_0^H J\phi C_{\mathrm{pos}}\,\dd r
 =-\int_0^H J(\partial_r\phi)j_{\mathrm{pos}}\,\dd r .
 \label{st:positive-moments}
\end{equation}
Thus the conservative representation determines every source
moment. The zero endpoint flux contributes no new interface term,
including when $H$ varies tangentially. The physical common mass,
momentum and $k$ fluxes remain those of the selected coupled wall
element; the remedy adds no independently assigned interface flux.

The dependent inner dissipation becomes $E_{\mathrm{pos}}$.
Its outer boundary value must be
\begin{equation}
 \boxed{\eps_{o,H}=E_{\mathrm{pos}}(H)
                 =\lambda_{\mathrm p}E_{\mathrm{pre}}(H).}
 \label{st:positive-epsilon-trace}
\end{equation}
There is no solved inner dissipation transport equation and no
additional prescribed outer dissipation flux. Leaving the trace
at $E_{\mathrm{pre}}(H)$ when $\lambda_{\mathrm p}<1$ would violate
the selected common dissipation distribution.
Any trace projection or lifting must respect the attenuated value.

If $b_{\mathrm p}=0$, then $m_{\mathrm p}=\delta_{\mathrm p}$,
$\lambda_{\mathrm p}=0$, and the outer dissipation trace is zero.
This can occur even with $K>0$: for a ray with $\chi=1$, a positive
point sample can lie in a molecular-only region where
$D_{\mathrm{pre}}=p_{A,0}=0$, while $k$ vanishes wherever
$p_{A,0}>0$. The weighted numerator defining $b_{\mathrm p}$ then
vanishes despite that positive sample.
The strict outer constitutive domain is $k>0,\eps>0$, so such a
state requires a declared outer continuation or an invalid-state
outcome; the inner reaction remedy does not supply a full coupled
admissibility theorem. The same qualification applies to near-zero
traces and loss of regular coefficient bounds. An enforced positive
lower bound on $b_{\mathrm p}$ is an additional admissibility policy.
A positive dissipation floor is a model change whose boundary trace
and modified budget must be specified together.

The limiter can remain active at $\chi=0$. Therefore exact recovery
of the unmodified finite-viscosity SCI baseline holds only when
the positivity limiter is inactive. With it active, the zero-$\chi$
operator is the explicitly positivity-adapted baseline.
This distinction applies to dissipation traces as well as source
loads.

\subsection{Inner covariance and what the remedy leaves unresolved}
\label{st:inner-covariance}

For the inner full-strain completion,
\[
 R_A=\tfrac23kI-2\nu_{t,A}S_w,\qquad
 T_w=-(\pi_w+\tfrac23k)I+2(\nu+\nu_{t,A})S_w.
\]
Molecular stress belongs in $T_w$, not in the Reynolds covariance.
At incompressibility, the exact covariance condition is
\begin{equation}
 \boxed{\nu_{t,A}\lambda_{\max}(S_w)\le k/3.}
 \label{st:inner-realizability}
\end{equation}
On the flat reference, $G_A=t_m/D_A$ and $K=q_m u_m^2$, so
\begin{equation}
 \lambda_{\min}(R_A)
 =K\left[\frac23-\frac{\nu_{t,A}/D_A}{q_m}\right]
 \ge K\left(\frac23-\frac1{q_m}\right)>0.
 \label{st:inner-reference-covariance}
\end{equation}
All three roots exceed $3/2$. This conclusion remains valid for any
positive auxiliary slope $\alpha$, including an optional
overlap-consistent recalibration, because
$0\le\nu_{t,A}/D_A<1$ at finite $\nu$.

Away from reference the sampled viscosity can remain positive while
local $k$ depletes. The source remedy restores a reaction sign but
does not enforce \cref{st:inner-realizability}. An active stress
limiter would be a further change and must be used consistently
in momentum, production and traction. It generally changes the
analytic momentum kernel and invalidates an unmodified integral
formula. Positivity at a subset of interpolation nodes also does
not imply positivity at quadrature, trace and sampling locations.

The exact budget identities require common geometric weights and
quadrature in $\mathcal B$, $b_{\mathrm p}$, the redistribution,
and the dependent boundary trace. Derivatives and stage updates
must include the dependence through $K,\chi,b_{\mathrm p}$ and
$\lambda_{\mathrm p}$. The min/max transitions are nonsmooth away
from their inactive plateaus; smoothing them is another specified
variant, not an automatic consequence of the reference proof.

\subsection{Outer transport positivity and the final separation of claims}
\label{st:outer-positivity}

The outer transport equations
\cref{out:kouter,out:epsA,out:epsE} are defined on $k>0,\eps>0$,
with the additional
A95 real-domain constraint. On a regular continuation with $k$
bounded away from zero, their dissipation reactions contain a
factor of $\eps$ and vanish as $\eps\to0$. With positive regular
diffusion and compatible data this is the necessary reaction
sign for a scalar comparison argument. It is not a global theorem
for the original coupled equations: invariant coefficients,
ratios and diffusion can lose regularity at the domain boundary.

Likewise, $P_k\ge0$ for A95/E does not establish invariance of a
closed nonnegative $(k,\eps)$ quadrant. The $k$ sink is $-\eps$,
and parts of that quadrant's boundary are outside the constitutive
domain. For a useful restricted check, homogeneous zero-production
decay satisfies
\[
 \dot k=-\eps,\qquad
 \dot\eps=-C_{\eps2}\frac{\eps^2}{k},\qquad
 \eps=c\,k^{C_{\eps2}}.
\]
For $C_{\eps2}>1$ and positive initial data,
\[
 k(t)=\left[k(0)^{1-C_{\eps2}}
          +(C_{\eps2}-1)c\,t\right]^{-1/(C_{\eps2}-1)}>0.
\]
This special solution remains positive at every finite time; it
does not prove positivity for an inhomogeneous wall-coupled
spectral discretization.

\begin{table}[htbp]
\centering\small
\caption{Distinct properties of the unchanged outer stress laws.
The wall reaction remedy supplies only its conditional inner result.}
\label{st:comparison}
\begin{tabularx}{\linewidth}{@{}>{\raggedright\arraybackslash}p{.28\linewidth}YYY@{}}
\toprule
Property & A93 & A95 & E94/E95\\
\midrule
Certified neutral principal bound
& $0.036\nu_{t,o,0}$ & $0.02\nu_{t,o,0}$
& $4\nu_{t,o,0}/q_E^2$\\
Negative principal eigenvalue away from reference
& Simple shear, $q=8$ & Simple shear, $q=1$ & Pure strain\\
Full covariance
& Can fail for unrestricted gradients
& Positive on the real coefficient domain
& Positive where defined; isotropic stress extension at $S=0$\\
Production everywhere nonnegative
& No & Yes on the real domain & Yes where defined\\
Sampled-$K$ source obstruction
& Present & Present & Present\\
Global coupled or discrete stability proved
& No & No & No\\
\bottomrule
\end{tabularx}
\end{table}

The continuum reference, local momentum diffusion, covariance,
transported-field positivity, exact discrete stationarity and
nonlinear robustness therefore require separate evidence.
The constitutive calculations and wall-source construction provide
the stated algebraic and conditional continuum results. The
implementation and validation
programme must additionally test source moments, full traction,
common fluxes, admissibility, time stages, perturbation growth and
physical terrain response.


\section{Proposed Nektar++ realization and discrete stationary states}
\label{num:implementation}

\subsection{A concrete implementation contract}

Nektar++ provides a spectral/$hp$ framework with continuous and
discontinuous Galerkin infrastructure
\citep{Nektar2015,Nektar2020}.  Its documented incompressible
velocity-correction formulation separates advection, pressure, and
velocity solves and requires compatible pressure boundary conditions
for the intended temporal accuracy \citep{NektarTheory}.  These
capabilities motivate the following implementation proposal.  They
do not supply an existing implementation or verification of this
composite turbulence/wall model.

The numerical specification should record the following objects
before assembling a reference test.
\begin{table}[htbp]
\centering\small
\caption{Proposed field and operator contract.  Each choice is part
of the actual discrete model to be verified.}
\label{num:assembly-contract}
\begin{tabularx}{\linewidth}{@{}p{0.23\linewidth}X@{}}
\toprule
Object & Required specification\\
\midrule
Cartesian velocity & Inner patch/reconstruction, outer trial space,
strong ground no slip, common interface trace, physical gradients,
and treatment of storage before condensation.\\
Pressure & Isotropic-stress convention, admissible relative
inner/outer pressure modes or lifting, full traction coupling,
one global gauge, and pressure-stage compatibility.\\
Transported $k$ & Inner and outer spaces, positive sample
$\mathcal K$, advection/diffusion, actual local production,
positivity-adapted destruction when enabled, and shared fluxes.\\
Dissipation & Outer transport space, dependent inner profile $E_w$,
eliminated independent inner transport rows, and the single
actual trace $E_w(H)$.\\
Geometry and transfer & Physical $H,r_K,r_m,h_r$, valid maps,
Jacobian and normal orientation, mortar/common trace spaces,
and constant-preserving interface projections.\\
Coefficients and integration & A93/A95/E choice, selected $\kappa$
and compatible outer coefficients, wall coefficient variant,
cutoff/state gates, split quadrature, and all limiter parameters.\\
\bottomrule
\end{tabularx}
\end{table}

A wall-patch module should expose its common mass flux, vector
momentum flux, $k$ flux, dissipation trace, and their dependence on
the retained variables.  It should also expose diagnostics for
ground reaction, forcing work, molecular loss, turbulent production,
dissipation, and redistribution.  For the flat quasisteady kernel,
closed derivatives of $J_i$ can be used.  A curved transient patch
instead needs its own coupled solve and consistent condensation;
freezing away its stored velocity history changes the model.

The outer momentum operator must evaluate the complete selected
stress tensor, including the algebraic terms in A93/A95 and the
state dependence of $C_\mu$ in every model.  It cannot be replaced
by a constant positive Laplacian while retaining the same
constitutive claim.  A positive implicit Laplacian plus an explicit
compensating stress is an algebraic time split; it does not remove
the demonstrated negative principal directions of the original
continuous operator.  Filtering, stress limiting, or added diffusion
must be stated as additional numerical or modeling choices and
included in both the operator contract and its tests.

The implementation order should begin with a flat coupled patch,
then add the conservative curved map and common fluxes, and then
the spatial transition.  This makes failures in the wall law,
pressure coupling, metric identities, and switch transfer separately
observable.  A scalar boundary-condition entry alone cannot encode
the independent field constraints and conservative patch system
specified here.

\subsection{Reference representation and spatial equilibrium}
\label{num:reference-representation}

A finite polynomial space generally does not contain the
logarithmic velocity, reciprocal outer dissipation, or the complete
piecewise inner reconstruction.  Nodal interpolation of the analytic
reference therefore defines a new discrete state whose derivatives
and nonlinear stress need not satisfy the analytic identities.
Raising polynomial order can reduce this defect, but does not make
it identically zero at a fixed order.

An equilibrium-aware representation can keep the fixed admissible
analytic reference explicitly and approximate perturbations:
\begin{equation}
 \boldsymbol Q_h(\boldsymbol x)
 =\boldsymbol Q_0(\boldsymbol x)
  +\sum_j\boldsymbol q_j\psi_j(\boldsymbol x).
 \label{num:reference-lifting}
\end{equation}
Here $\boldsymbol Q_0$ denotes the appropriate piecewise velocity,
transported scalars, and pressure lifting; dependent fields still
obey their stated constraints.  The shape functions and trace
constraints must respect the ground and interface conditions.
For the designated target, zero perturbation represents that target
exactly.  This construction changes the representation while
leaving the selected physical equations and sources unchanged.
It does not by itself make their numerical integration exact.

A claim of exact discrete spatial preservation requires all of the
following: analytical or compatible differentiation of the retained
reference; the same constitutive invariants in stress and production;
constant preservation by the $K$ sampler; the same coefficient
variant and $\kappa$ throughout; all-mode source compatibility;
correct actual interface normals, fluxes, and pressure offsets;
and quadrature satisfying the remaining weak identities.  Closed
weak primitives or algebraically compatible evaluations may enforce
particular identities.  Otherwise establish integration convergence
to a stated tolerance.  Polynomial-order and quadrature refinements
are separate checks.

Let $\mathcal R_h$ be the \emph{native assembled} spatial residual
of the specified equations, including boundary and constraint rows,
with the global pressure gauge applied.  The qualification test is
\begin{equation}
 \mathcal R_h(\boldsymbol Q_{0,h})=0
 \quad\text{to the stated arithmetic/assembly tolerance},
 \label{num:spatial-stationarity}
\end{equation}
where $\boldsymbol Q_{0,h}$ is the exact represented target of the
chosen numerical contract.  Inspect the separate momentum,
continuity, $k$, outer-$\eps$, constraint, and interface residuals,
not only a global norm that can hide cancellations between equations.
If \cref{num:spatial-stationarity} fails for an ordinary interpolated
reference, report that defect and its refinement trend.  The
continuum equilibrium proposition remains a different statement.

\subsection{Boundary and pressure compatibility}
\label{num:pressure-compatibility}

The lower outer normal is $-\boldsymbol N_H$.  Apply the common
mass and momentum fluxes from \cref{geo:massflux,geo:interface-momentum}
with that orientation and the same surface measure in both regions.
A continuous mechanical-pressure representation that forces the
A93/A95 reference offset to zero cannot represent the prescribed
normal traction.  Possible designs include independent patch
pressure modes or an appropriate pressure lifting, always with a
single consistent gauge and traction condition.

For a pressure-correction stage, the normal momentum relation
\cref{geo:pressure-compatibility} couples the pressure condition to
the time-dependent transpiration trace, outer convective term,
full stress divergence, and body force.  These terms must use the
stage levels of the selected time discretization.  Homogeneous
pressure Neumann data are not an automatic replacement.  A pure
Neumann pressure problem must also satisfy its integrated
compatibility condition; the wall flux and other boundary fluxes
must sum consistently with the divergence constraint.

Use reference-compatible inlet, top, and outlet data in a flat
preservation test.  The top must supply the reference shear and
outer-$\eps$ flux, or compatible analytic values, as in
\cref{out:referencefluxes}.  A stress-free top or a zero dissipation
gradient is a different problem.  At the model interface, the
single outer dissipation value is the actual $E_w(H)$, including
positivity attenuation.  Prescribing its flux as well would add
an unwanted constraint.  Common $k$ and its shared flux must remain
consistent with constant reference $k$ and zero ground flux.

Trace transfer is part of these boundary equations.  Preserve its
constant tests, its designated reference, and its prescribed
zero-correction support.  A conservation-preserving mortar alone
does not guarantee a positive dissipation trace; a positivity
procedure alone does not guarantee the required flux moments.
Both properties must be qualified for the selected spaces.

\subsection{Stationary time stepping is a separate property}
\label{num:stationary-map}

Let $\mathcal T_{\Delta t,h}$ denote a complete time step, including
all velocity/pressure substeps, turbulence updates, wall solves,
trace projections, gates, and limiting.  For a multistep method,
initialize every required history level consistently with the
same target.  The discrete stationary-state test is
\begin{equation}
 \boxed{\mathcal T_{\Delta t,h}
 (\boldsymbol Q_{0,h},\boldsymbol Q_{0,h},\ldots)
 =\boldsymbol Q_{0,h}}
 \label{num:temporal-stationarity}
\end{equation}
to its declared solver and arithmetic tolerance.  Also initialize
pressure and inner patch histories consistently.  All relevant
stages must preserve the reference values of $K$, the roughness
calibration, cutoff, switch, dependent dissipation, and pressure
traction.  Repeated coefficient lagging, extrapolation, or resetting
of a dependent trace can otherwise introduce a defect even when
individual continuum terms cancel.

A zero sum of spatial residual terms is not enough to prove this
identity for an arbitrary split method: pieces that cancel in the
coupled equations can be evaluated at different states or enter
incompatible boundary stages.  Check a complete single step first,
then repeated steps over several flow-through times, and refine
$\Delta t$ separately from mesh size, polynomial order, quadrature,
and iterative tolerances.  Record both the maximum step change and
cumulative drift.  No specialized steady nonlinear algorithm is
required merely to state or test this property of a URANS
integrator.

Exact stationary stepping, perturbation stability, turbulence
positivity, and constitutive realizability are distinct properties.
The neutral frozen-$k,\eps$ momentum certificates do not certify
all coupled turbulence or pressure-splitting modes.  Conversely,
a smaller time step cannot correct an admissible continuous state
whose momentum principal operator has a demonstrated negative
diffusion eigenvalue.  A solver may pass the stationary-reference
test and still require a restricted constitutive state domain or
a separately declared continuation for off-reference computation.

\section{Reproducible verification and the evidence obtained here}
\label{num:verification}

\subsection{What is already checked}

The analytic derivations establish the conditional continuum
reference, exact finite-viscosity wall identities, conservative
metric transformations, and the constitutive statements in the
preceding stability assessment.  The executable quadrature test in
\cref{app:wall-verification} supplies a separate numerical check of
the wall identities.  Its recorded run passed \textbf{873 assertions}
in \textbf{22 cases}: eighteen rough-reference combinations from
three outer models, two choices of outer $\kappa$, and three wall
coefficient variants; one smooth molecular-sublayer case; one
all-molecular rejection case; and two forced/reversed vector cases.
The two $\kappa$ choices are each model's natural value and the
common value $0.41$ with its compatible outer coefficient.  The
wall variants retain the original coefficients, match the viscosity
slope, or match both tail coefficients as specified earlier.

The script compares closed $J_0,J_1,J_2$ with independent adaptive
quadrature, calibrates $h_r$ by the original roughness correlation,
and checks the production/molecular-work split, the corrected
cutoff and its gate, the dissipation budget, full reference traction,
and pressure offsets.  Its modal check evaluates the
redistribution through integration by parts against Legendre tests
of degrees zero through six, rather than merely reusing the
pointwise simplified zero.  For the Legendre polynomial $L_i$,
the evaluated quantity is
\begin{equation}
 M_i=\int_0^H L_i(2z/H-1)(p_{A,0}-E_{c,m})\,\dd z
     -\int_0^H\frac2H L_i' (2z/H-1)j_c\,\dd z,
 \qquad 0\leq i\leq6 .
 \label{num:modal-check}
\end{equation}
These moments have units $\mathrm{m^3\,s^{-3}}$.  It also compares the matching-height
derivative with a finite difference and checks the overlap-slope
relations.  These are local analytical and quadrature tests; they
do not invoke a PDE solver.

For illustration, take $H=2\,\mathrm m$, $z_m=1\,\mathrm m$,
$z_0=0.03\,\mathrm m$, $\ust=0.4\,\mathrm{m\,s^{-1}}$,
$\nu=1.5\times10^{-5}\,\mathrm{m^2\,s^{-1}}$, common
$\kappa=0.41$, and the original wall coefficients.  The target
resistance is $J_0(H)=25.607957792\,\mathrm{s\,m^{-1}}$.
\Cref{num:finite-wall-table} reports values from the reproduced
calculation; full precision and the input definitions are retained
in \cref{app:wall-verification}.
\begin{table}[htbp]
\centering\small
\caption{Finite-$\nu$ rough-wall reference check.  $I_0=\int_0^H
p_{A,0}\,\dd z$ and $M_0=\nu t_m^2J_2(H)$ are kinematic work
rates per unit wall area, in $\mathrm{m^3\,s^{-3}}$.}
\label{num:finite-wall-table}
\begin{tabular}{@{}lrrrrr@{}}
\toprule
Model & $h_r\ (\mathrm m)$ & $I_0$ & $M_0$ & $\Lambda_m$ & $z_c\ (\mathrm m)$\\
\midrule
A93 & 0.4338100064 & 0.6550655195 & 0.0004982000 & 4.1965134844 & 0.08180914015\\
A95 & 0.3396336383 & 0.6549850211 & 0.0005786984 & 4.1959977913 & 0.08185133944\\
E94/E95 & 0.4543819865 & 0.6550791365 & 0.0004845830 & 4.1966007180 & 0.08180200395\\
\bottomrule
\end{tabular}
\end{table}

All three pass the independent cutoff gate.  Each cutoff exceeds
$e z_0=0.08154845485\ldots\,\mathrm m$ by the finite molecular
correction.  Across the eighteen rough-reference cases, the largest
absolute corrected weak moment is $9.86\times10^{-16}$ $\mathrm{m^3\,s^{-3}}$, whereas the corresponding largest
uncorrected nonconstant moment in each case ranges from $0.04627$
to $0.07786$.  These values show the difference between a matched
mean and the necessary nonconstant modes for this test family.
They are not a universal floating-point error bound.

The smooth test has $H=0.05\,\mathrm m$, $h_r=0$, and a
molecular-only interval of thickness
$z_v=2.20065243\times10^{-4}\,\mathrm m$.  Its integrated turbulent
production is $0.8141890702$ and its molecular mean-flow loss is
$0.5282476121\,\mathrm{m^3\,s^{-3}}$; the corrected weak moments
are below $1.45\times10^{-14}$ $\mathrm{m^3\,s^{-3}}$.  The
all-molecular interval instead has $P_A=0$, $\Lambda_m=0$, and
fails the EP cutoff gate.  A finite total shear work in that interval
is molecular loss; it cannot finance a positive corrected turbulent
dissipation.

In the rough forced/reversed case the ground shear is
$(-0.080,-0.056)\,\mathrm{m^2\,s^{-2}}$, while the upper shear is
approximately $(1.78410,0.378956)\,\mathrm{m^2\,s^{-2}}$.
Upper mechanical work $7.30993684$ equals turbulent production
$6.27552505$, molecular loss $0.00053297$, and forcing work
$1.03387882$, in $\mathrm{m^3\,s^{-3}}$.
Substituting the ground shear at the upper interface would miss
$7.63771870\,\mathrm{m^3\,s^{-3}}$ in that work comparison.
The signs and magnitudes are deliberate test inputs, not an
experimental prediction.  The smooth forced/reversed case supplies
an independent check with a substantial molecular contribution.

\subsection{Verification levels and acceptance measures}
\label{num:evidence-levels}

\begin{table}[htbp]
\centering\small
\caption{Evidence levels.  ``Required'' denotes work for an actual
implementation or application, not a test claimed as completed.}
\begin{tabularx}{\linewidth}{@{}p{0.07\linewidth}p{0.28\linewidth}X@{}}
\toprule
Level & Question & Evidence or acceptance requirement\\
\midrule
V0 & Is the model fully specified? & Record the complete contract in
\cref{num:assembly-contract}, including coefficient variants,
state domain, projections, gates, and physical sampling scales.\\
V1 & What is known about the outer principal operator? & Use the
analytic neutral certificates and explicit off-reference
counterexamples; a finite directional scan is not a proof over
all states or directions.\\
V2 & Are the scalar wall identities correct? & The 873-assertion
quadrature run and its full executable source are supplied in
\cref{app:wall-verification}.\\
V3 & Are curved metric and shared fluxes implemented correctly? &
Required: manufactured patch checks of geometry, continuity,
Cartesian momentum, normal pressure, $k$ flux, and energy.\\
V4 & Does the native spatial operator preserve its target? &
Required: separately report every residual group in
\cref{num:spatial-stationarity}, constant-test budgets, and
interface flux cancellation.\\
V5 & Does the actual time integrator keep that state fixed? &
Required: consistent-history one-step and long-time checks of
\cref{num:temporal-stationarity}, with independent temporal and
iterative-tolerance refinement.\\
V6 & Are positivity and constitutive gates respected? &
Required: depleted turbulence, active source adaptation,
trace transfer, sampled versus local energy, and all declared
invalid-state outcomes.\\
V7 & Is off-reference computation convergent and robust? &
Required: perturbations, forced/reversed shear, transition,
curvature, and $h$, $p$, quadrature, and time refinement.\\
V8 & Are physical predictions adequate? & Required:
independent rough-wall and terrain data with uncertainty and
sensitivity to the wall/outer model choices.\\
\bottomrule
\end{tabularx}
\end{table}

A useful native residual report includes nondimensionalized norms
for each equation, its largest individual row, and its constant-test
balance.  Also record global mass imbalance; the jump in common
traction and $k$ flux after projection; the dependent dissipation
constraint; the pressure gauge; minimum $k$ and $\eps_o$ over the
claimed positivity set; and the minimum covariance eigenvalue and
constitutive-domain margin over the sampled state set.  A sampled
minimum is a diagnostic for those samples, not an all-space
certificate.  Local wall tests should log
\begin{equation}
 \begin{aligned}
 \Delta_P&=\int_0^HJ(P_A-D_w)\,\dd r,\qquad
 \Delta_C=\int_0^HJC_w\,\dd r,\\
 \Delta_{\phi_i}&=\int_{\Omega_w}\phi_i(P_A-D_w)\,\dd V,
 \qquad
 \Delta_{\eps,H}=\eps_{o,H}-E_w(H).
 \end{aligned}
 \label{num:diagnostic-balances}
\end{equation}
At the reference the source moments vanish; away from it,
$\Delta_P$ is generally a physical modeled source and should not
be forced to zero.  The redistribution budget $\Delta_C$ remains
zero for the stated conservative construction.

For V3, use a smooth admissible curved patch with nonzero tangential
variation and variable $H$.  Construct a divergence-free Cartesian
velocity satisfying ground no slip, evaluate the exact metric
operators, and compare independent boundary and volume integrals.
Include the height-slope term in mass and scalar flux, ground
traction orientation, lateral patch fluxes, pressure curvature,
and the metric-weighted cutoff.  A flat-only geometry test cannot
exercise these terms.  Invalid normal maps must be rejected before
evaluating wall coefficients.  The published geometric identities
and a positive numerical Jacobian at selected nodes do not establish
global ray injectivity for an unspecified terrain mesh.

For V4--V7, keep the physical model fixed during numerical
refinement: $H$, $h_r$, the sampling maps, gate support, and calibrated
roughness must not drift with polynomial order.  Separate
matching-height sensitivity from discretization convergence using
\cref{kap:H-derivative}.  Refinement of the native residual and
one-step map resolves numerical compatibility; perturbation tests
resolve dynamic behavior.  Measure growth or decay with its spatial
scale, distinguish physically admissible flow response from
constitutive or numerical instability, and include the already
identified negative-principal-symbol states as regression cases
for whatever state restriction or continuation is adopted.

\section{Implementation and physical validation still required}
\label{num:remaining}

The supplied result is a self-contained continuous model, its
conditional flat reference, constitutive analysis, and reproducible
local numerical checks.  It does not include a modified Nektar++
repository, a built coupled solver, a native residual evaluation,
a pressure-split stationary run, or a terrain-flow simulation.
Accordingly, no code-inspection, solver-acceptance, or experimental
validation claim is made for this proposed implementation.

The first remaining engineering task is the native coupled wall
patch and outer turbulence assembly with the exact field contract,
followed by the V3--V6 checks.  A pressure formulation capable of
the required normal traction and a conservative common-flux transfer
are essential parts of that work.  The transported/dependent scalar
split, positivity-induced trace change, and inactive-branch
coefficient evaluation must remain visible in the implementation,
including at $\chi=0$.

Before an unrestricted off-reference computation is accepted, the
outer state policy must be resolved.  The original closures have
explicit limitations documented in the stability assessment.
Restricting the computational state to a justified admissible
neighbourhood, or introducing a separately analyzed constitutive
continuation, is a modeling decision.  A positive wall resistance,
a positive covariance, positive mean production, and neutral
frozen-variable ellipticity each provide useful information, but
together they still do not supply a global coupled existence,
positivity, or stability theorem.  Any active stress modification
also changes production and must be applied consistently.

Physical validation should first assess flat rough-wall mean
velocity, friction, $k$, and dissipation over independent roughness
and Reynolds-number ranges.  The one-time resistance calibration
of $h_r$ to a target $z_0$ is an internal model relation, not a
universal empirical conversion between aerodynamic and sand-grain
roughness.  Compare original and overlap-adjusted wall coefficients
using their corresponding outer $\kappa$, and quantify the effects
of $H$, sampling height, tail start, and positivity parameters.
Hold out independent data when fitting those choices.

Next test the conservative curved completion on terrain with
measured or otherwise independently established pressure,
acceleration, separation/reattachment, shear, and turbulent-energy
profiles.  Include crossflow, adverse pressure gradients, roughness
changes, and the transition between the equilibrium correction and
its explicit baseline.  Vary the transition support and state-gate
thresholds; report their physical length scales and the influence
of upstream pressure coupling.  Exact preservation is asserted for
the designated flat reference on its unit gate plateau, not for a
nonuniform hill solution merely because it is near that plateau.

Finally, separate numerical, model, and data uncertainties.  Fix
$H$ and all physical sampling scales for $h/p/\Delta t$ convergence;
then vary those model scales independently.  If $H$ is changed,
report whether $h_r$ is held fixed or a new one-time calibration
is performed, because \cref{kap:H-derivative} shows that these are
different physical model comparisons.  A change of coefficient
variant, limiter, filter, or trace projection creates a new
verification contract.  The equilibrium and conservation identities
provide precise tests of that contract, while prediction quality
must be established by the remaining implementation and physical
validation programme.

\section{Comparison and recommended specification}
\label{sec:specification}

\subsection{What the models share and where they differ}

All three original stress laws admit a constant-energy, constant-stress
neutral branch. The nonlinear A93/A95 correction has zero $xz$
component on that branch, but the total Reynolds shear is
$R_{xz,0}=-\ust^2$. Each model's constitutive root fixes the energy
level and the full reference covariance. Its dissipation equation
then fixes the compatible logarithmic slope when the original
transport coefficients are retained.

\begin{table}[htbp]
\centering\small
\caption{Selected outer-model distinctions. The constitutive lower bounds
are for incompressible rank-one velocity perturbations at fixed positive
$k,\eps$ on the canonical reference.}
\label{tab:comparison}
\begin{tabularx}{\textwidth}{@{}>{\raggedright\arraybackslash}p{0.27\textwidth}YYY@{}}
\toprule
Property&A93&A95&E94/E95\\
\midrule
Stress correction & Three quadratic gradient tensors &
One strain--rotation tensor & No explicit nonlinear stress tensor\\
$q_m$ & $3.4012654174$ & $3.8220004261$ & $3.3245064563$\\
Original $\sigma_\eps$ & $1.3$ & $1.3$ & $1.2$\\
Natural $\kappa_m$ & $0.4283236287$ & $0.4040609953$ &
$0.4120874716$\\
$\sigma_\eps^{\rm eq}$ at $\kappa=0.41$ &
$1.1911514931$ & $1.3384963992$ & $1.1878733447$\\
Reference constitutive lower bound divided by $\nu_{t,0}$ &
$0.036$ (sufficient) & $0.02$ (sufficient) &
$4/q_E^2$ (sharp)\\
Negative projected diffusion away from reference &
Explicit example & Explicit example & Explicit example\\
Unrestricted covariance qualification &
Counterexamples exist & Positive on stated real domain &
Positive on stated coefficient domain\\
\bottomrule
\end{tabularx}
\end{table}

The wall-element construction is shared in form but uses the selected
model's $q_m$, $\kappa$, outer stress, scalar transport coefficients,
and pressure offset. This common structure does not make the three
coupled models dynamically equivalent. In particular, a more
favourable sufficient reference tangent bound does not rank their
terrain-flow accuracy or nonlinear stability.

\subsection{Principal formulation}

For a common prescribed atmospheric reference, use the following
specification:
\begin{enumerate}
\item Select $\kappa,z_0,\ust,H,Z,\nu$ and the outer model; calculate
its exact $q_m$ and use
$\sigma_{\eps,m}^{\rm eq}=\kappa^2q_m/D_m$.
Retain the complete original stress law on its stated domain.
Alternatively, retain the original outer coefficients and consistently
use its natural $\kappa_m$.
\item Retain the original SCI auxiliary viscosity coefficients and
roughness correlation. Declare the sampler $\mathcal K$ and all
physical heights. Calibrate $h_r$ once at $K_0=q_m\ust^2$ if a
prescribed $z_0$ is required, and verify the resistance range.
\item Compute the actual finite-viscosity production budget and
corrected cutoff. Use the same selected $\kappa$ in the amplitude,
roughness anchor, cutoff, and outer boundary data. Require a strict
admissibility margin on the reference plateau.
\item Reconstruct the inner mean velocity with numerical no slip,
retain inner $k$ transport, constrain inner $\eps$ to its declared
profile, and exchange one shared mass, full-traction, and $k$ flux
at the matching surface. Supply the outer $\eps$ trace consistently.
\item Use the all-mode conservative source redistribution. If the
bounded-rate positivity adaptation is selected, apply it to the whole
blended destruction and change the dependent dissipation trace with it.
Its inactive plateau must contain the equilibrium.
\item Use a conservative curved patch and a declared spatial/state
transition for terrain. Preserve storage, tangential transport,
pressure--normal-stress balance, and actual interface transpiration.
\item Qualify the implemented spatial and temporal operators against
the stated reference and evidence levels before terrain predictions
are interpreted physically.
\end{enumerate}

\subsection{Optional variants that must be named}

The wall overlap variant changes
$\alpha$ to $\kappa/\sqrt{q_m}$ and repeats the entire resistance,
production, cutoff, and response calculation. A further baseline-tail
variant also sets $c_\ell^w=\kappa q_m^{3/2}$ and
$c_\mu^w=q_m^{-2}$. These options provide additional asymptotic
matching properties, while the finite molecular/roughness intercept
remains. They change the original wall law's response and require
their own nonequilibrium validation.

An outer destruction-coefficient alternative is given in
\cref{out:appendixCeps2}; it replaces the chosen outer
diffusivity adjustment rather than being combined with it.
An additional constitutive regularization away from the reference,
if required by a solver, must likewise be stated and assessed as a
new model. No universal off-reference regularization or completed
terrain validation is claimed here.


\clearpage
\appendix

\section{Tensor identities, parameter roles, and dimensions}
\label{app:tensor}

For $A=S+\Omega$, $S^T=S$, and $\Omega^T=-\Omega$,
\begin{align}
 A^2+(A^2)^T&=2(S^2+\Omega^2),\\
 AA^T&=S^2-\Omega^2+\Omega S-S\Omega,\\
 A^TA&=S^2-\Omega^2-\Omega S+S\Omega .
\end{align}
Therefore the quadratic A93 tensor can equivalently be written as
\begin{align}
 &-4\dev\{A^2+(A^2)^T\}
       +13\dev(AA^T)-2\dev(A^TA)\nonumber\\
 &\hspace{12mm}=3\dev(S^2)-19\dev(\Omega^2)
                         +15(\Omega S-S\Omega).
\end{align}
The commutator $\Omega S-S\Omega$ is symmetric and trace-free.
These identities fix the rotation-sign convention and prevent
coefficient comparisons from mixing distinct tensor bases.
For $A=G\boldsymbol e_x\otimes\boldsymbol e_z$, all the quadratic
terms above are diagonal. Their $xz$ contribution is zero, while
the linear stress contribution remains $-\nu_tG$.

For a trace-free symmetric $S$ with $s=(S:S)^{1/2}>0$, its largest
eigenvalue satisfies
\begin{equation}
 \lambda_{\max}(S)
 =\sqrt{\frac23}\,s
 \cos\left[\frac13\arccos
 \left(\frac{\sqrt6\,\tr(S^3)}{s^3}\right)\right].
\end{equation}
Hence the original strain-invariant coefficient
$A_s=\sqrt6\cos[\frac13\arccos(\sqrt6\,\tr(S^3)/s^3)]$
equals $3\lambda_{\max}(S)/s$.
The identity is useful for both realizability estimates and
coefficient derivatives; repeated eigenvalues require the regularity
qualifications in the main text.

\begin{table}[htbp]
\centering\small
\caption{Dimensional check for the physical-coordinate wall equations.
All stresses and forces here are divided by constant density.}
\begin{tabular}{@{}ll@{}}
\toprule Quantity & Units\\ \midrule
$k,K,t_m,\pi,R_{ij}$ & $\mathrm{m^2\,s^{-2}}$\\
$\nu,\nu_{t,A},D_A$ & $\mathrm{m^2\,s^{-1}}$\\
$\eps,E_A,E_{c,m},P_A,p_{A,0},D_{\rm pre},C_{\rm pre}$
 & $\mathrm{m^2\,s^{-3}}$\\
$\boldsymbol F$ & $\mathrm{m\,s^{-2}}$\\
$A_m,\int_0^H p_{A,0}\dd z,j_c$ & $\mathrm{m^3\,s^{-3}}$\\
$J_0,J_1,J_2$ & $\mathrm{s\,m^{-1}},\ \mathrm{s},\ \mathrm{s^2\,m^{-3}}$\\
$B=J_1/J_0$ & $\mathrm{m}$\\
$q_m,\kappa,\alpha,b,\Lambda_m,M_{\nu,m},\chi$ & dimensionless\\
\bottomrule
\end{tabular}
\end{table}

The original outer $C_\mu(A,k,\eps)$, the auxiliary constant
$c_\mu^w$, and its optional neutral-matching value $q_m^{-2}$
are different quantities. Likewise, the original dissipation cutoff
$z_{\eps,A}$ is set by $c_\ell^w$ and $\nu$, while the corrected
cutoff $z_c$ is set by the integrated turbulent production budget.
The matching height $H$, aerodynamic roughness $z_0$, sand-grain
input $h_r$, and virtual viscosity intercept $z_v$ cannot be used
interchangeably.


\section{Closed finite-viscosity AWF integrals}
\label{app:integrals}

\subsection{Resistance and forcing primitives}

With $a_A=\alpha\sqrt K>0$ and $z_v=\nu b/\sqrt K$, the
molecular-only segment has $D_A=\nu$ and the active segment has
$D_A=\nu+a_A(z-z_v)$. Splitting each integral at the physical
breakpoints gives the following formulas. Logarithms always have
dimensionless positive ratios.

Let $v=(z_v)_+$ and $\ell(z)=\min(z,v)$.
Then
\begin{align}
 J_0(z)&=\frac{\ell(z)}{\nu}
       +\frac1{a_A}\ln\frac{D_A(z)}{D_A(\ell(z))},
       \label{awf:J0}\\
 J_2(z)&=\frac{\ell(z)}{\nu^2}
       +\frac1{a_A}\left[\frac1{D_A(\ell(z))}-\frac1{D_A(z)}\right].
       \label{awf:J2}
\end{align}
These formulas include the cases of a molecular-only whole layer and
a turbulent branch active at the ground. For $\ell_1=\max(h_r,v)$,
\begin{equation}
 \begin{aligned}
 J_1(z)&=\frac{(\min(z,v)-h_r)_+^2}{2\nu}
       +\mathbf1_{z>\ell_1}\left\{
          \frac{z-\ell_1}{a_A}
          -\frac{\nu+a_A(h_r-z_v)}{a_A^2}
             \ln\frac{D_A(z)}{D_A(\ell_1)}\right\}.
 \end{aligned}
 \label{awf:J1}
\end{equation}
The second term is evaluated only on its active interval.


\subsection{Original and corrected dissipation integrals}

The exact original mean is
\begin{equation}
 \int_0^H E_A\,\dd z=
 \begin{cases}
 \dfrac{HK^{3/2}}{c_\ell^w z_{\eps,A}},&H\le z_{\eps,A},\\[2mm]
 \dfrac{K^{3/2}}{c_\ell^w}
 [1+\ln(H/z_{\eps,A})],&H>z_{\eps,A}.
 \end{cases}
\end{equation}
For an admissible corrected cutoff $0<z_c\le H$,
\begin{equation}
 \int_0^z E_{c,m}\,\dd\xi=
 \begin{cases}
 A_m z/z_c,&0\le z\le z_c,\\
 A_m[1+\ln(z/z_c)],&z_c<z\le H.
 \end{cases}
 \label{app:Ec-primitive}
\end{equation}
These primitives yield a continuous zero-endpoint redistribution flux
when the production budget is satisfied.

\subsection{Production on a constant-forcing segment}

Consider an active-viscosity segment $[z_a,z_b]$ that does not cross
$h_r$. Write $D_a=D_A(z_a)>0$, $D_b=D_A(z_b)>0$ and
\begin{equation}
 \boldsymbol t_A+\boldsymbol F L_r(z)
 =\boldsymbol c+\boldsymbol d D_A(z).
\end{equation}
Below $h_r$ take $\boldsymbol c=\boldsymbol t_A$ and
$\boldsymbol d=0$. Above $h_r$, with
$d_0=\nu-a_Az_v$, take
\begin{equation}
 \boldsymbol d=\boldsymbol F/a_A,\qquad
 \boldsymbol c=\boldsymbol t_A-\boldsymbol F(h_r+d_0/a_A).
\end{equation}
Then the exact turbulent-production integral is
\begin{align}
 \int_{z_a}^{z_b}P_A\,\dd z
 =\frac1{a_A}\Bigg\{&
 |\boldsymbol d|^2
 \left[\frac{D_b^2-D_a^2}{2}-\nu(D_b-D_a)\right]\nonumber\\
 &+2\boldsymbol c\cdot\boldsymbol d
 \left[D_b-D_a-\nu\ln\frac{D_b}{D_a}\right]\nonumber\\
 &+|\boldsymbol c|^2
 \left[\ln\frac{D_b}{D_a}
 +\nu\left(\frac1{D_b}-\frac1{D_a}\right)\right]\Bigg\}.
 \label{app:forced-production}
\end{align}
The molecular mean-flow loss on the same segment is
\begin{align}
 \int_{z_a}^{z_b}\nu|\boldsymbol U_A'|^2\,\dd z
 =\frac{\nu}{a_A}\Bigg\{&
 |\boldsymbol d|^2(D_b-D_a)
 +2\boldsymbol c\cdot\boldsymbol d\ln\frac{D_b}{D_a}\nonumber\\
 &+|\boldsymbol c|^2\left(\frac1{D_a}-\frac1{D_b}\right)\Bigg\}.
 \label{app:forced-molecular}
\end{align}
On molecular-only segments, turbulent production is zero and the
molecular term follows by integrating
$|\boldsymbol t_A+\boldsymbol F L_r|^2/\nu$.
Adding the segment formulas reproduces the forced work identity.

For the unforced reference, define $\ell=\min(z,(z_v)_+)$,
$D_\ell=D_A(\ell)$, and $\eta_D=[D_A(z)-D_\ell]/D_\ell\ge0$.
The nonzero active contribution is
\begin{equation}
 \int_0^zp_{A,0}\,\dd\xi
 =\frac{t_m^2}{a_A}
 \left[\ln(1+\eta_D)
 -\frac{\nu}{D_\ell}\frac{\eta_D}{1+\eta_D}\right],
 \qquad z>\ell.
 \label{app:positive-production}
\end{equation}
It is zero when $z=\ell$. Evaluate small increments with
\texttt{log1p} and a short series or direct nonnegative quadrature
when cancellation is significant. The identity
$t_m^2(J_0-\nu J_2)$ is exact, but subtracting its terms is poorly
conditioned near a molecular-only state.


\section{Executable interval certificate and algebraic checks}
\label{st:certificate}

The following standalone Python program uses only the standard
library. It constructs the matrix of
\cref{st:a95-certificate-form} from the exact algebraic reference
constants, encloses every arithmetic operation at 60 decimal digits,
and checks all eight $LDL^T$ pivots. Each square-root result is expanded
by one representable unit on either side of the correctly rounded
decimal result. The matrix is obtained by polarization of the displayed
quadratic form on the stated trace-free basis; no fitted matrix or
directional search is used.

The code also encloses the A93 sufficient bound and sharp E bound,
checks the three off-reference negative acoustic values, and verifies
the accompanying covariance signs and the A93 production
counterexample. These are algebraic checks of the displayed formulas.
They are not a coupled flow solve, a proof of time-integrator stability,
or terrain validation.

Copy the listing into a Python file and run it with Python 3. A
successful run reports strictly positive lower endpoints for every
A95 pivot and ends with \texttt{PASS}. Retain the assertions when
reproducing the certificate.

\begin{lstlisting}
"""Interval certificates for A93, A95, and E94/E95.
Only the Python standard library is required.
All stress and tangent formulas are given in the main text.
"""
from decimal import Decimal, Context, ROUND_FLOOR, ROUND_CEILING
from fractions import Fraction

LOW = Context(prec=60, rounding=ROUND_FLOOR)
HIGH = Context(prec=60, rounding=ROUND_CEILING)

class I:
    """Closed interval with outward-rounded elementary arithmetic."""
    def __init__(self, lo, hi=None):
        if isinstance(lo, I):
            self.lo, self.hi = lo.lo, lo.hi
        else:
            self.lo = Decimal(str(lo))
            self.hi = Decimal(str(lo if hi is None else hi))
        assert self.lo <= self.hi

    def __neg__(self):
        return I(self.hi.copy_negate(), self.lo.copy_negate())

    def __add__(self, other):
        other = I(other)
        return I(LOW.add(self.lo, other.lo),
                 HIGH.add(self.hi, other.hi))

    __radd__ = __add__

    def __sub__(self, other):
        return self + (-I(other))

    def __rsub__(self, other):
        return I(other) + (-self)

    def __mul__(self, other):
        other = I(other)
        pairs = [(self.lo, other.lo), (self.lo, other.hi),
                 (self.hi, other.lo), (self.hi, other.hi)]
        return I(min(LOW.multiply(a, b) for a, b in pairs),
                 max(HIGH.multiply(a, b) for a, b in pairs))

    __rmul__ = __mul__

    def __truediv__(self, other):
        other = I(other)
        assert not (other.lo <= 0 <= other.hi)
        pairs = [(self.lo, other.lo), (self.lo, other.hi),
                 (self.hi, other.lo), (self.hi, other.hi)]
        return I(min(LOW.divide(a, b) for a, b in pairs),
                 max(HIGH.divide(a, b) for a, b in pairs))

    def __rtruediv__(self, other):
        return I(other) / self

    def sqrt(self):
        assert self.lo >= 0
        # Decimal sqrt is correctly nearest-rounded.
        # Expanding by one representable unit encloses the root.
        lower = max(Decimal(0), LOW.next_minus(LOW.sqrt(self.lo)))
        upper = HIGH.next_plus(HIGH.sqrt(self.hi))
        return I(lower, upper)

    def out(self):
        return "[{}, {}]".format(self.lo, self.hi)


# A95 exact neutral coefficients.
root2 = I(2).sqrt()
q95 = (3 + I(61).sqrt()) / (2 * root2)
alpha_c = 3 / (root2 * q95)
f = 1 - I("4.5") / (q95 * q95)
assert f.lo > 0
CQ = f.sqrt() / (1 + 3 * q95 * q95)
beta95 = CQ * q95 * q95 * q95
c = (1 - f) / f
e = 6 * q95 * q95 / (1 + 3 * q95 * q95)
lp = c * (1 - alpha_c) + e

def form(H):
    """The polynomial P(H), extended to trace-free matrices."""
    norm = sum(x*x for row in H for x in row)
    result = I(norm)
    result -= alpha_c * (H[0][2] - H[1][1]) * (H[0][2] + H[2][0])
    result += beta95 * (
        sum(H[0][j]*H[2][j] for j in range(3))
        - sum(H[i][0]*H[i][2] for i in range(3)))
    result += beta95 * (
        lp*H[0][2] + c*H[2][0] + c*alpha_c*H[1][1]
        ) * (H[0][0] - H[2][2])
    # These two minors vanish on every rank-one matrix.
    result += I("0.2") * (H[0][1]*H[1][2] - H[0][2]*H[1][1])
    result -= I("1.6") * (H[0][0]*H[2][2] - H[0][2]*H[2][0])
    # Desired dimensionless coercivity bound: 1/50.
    result -= I("0.02") * norm
    return result

# Coordinates: xx, xy, xz, yx, yy, yz, zx, zy; zz=-xx-yy.
basis = []
for i in range(3):
    for j in range(3):
        if (i, j) == (2, 2):
            continue
        H = [[0]*3 for _ in range(3)]
        H[i][j] = 1
        if i == j:
            H[2][2] = -1
        basis.append(H)

matrix = [[I(0) for _ in range(8)] for _ in range(8)]
for i in range(8):
    matrix[i][i] = form(basis[i])
for i in range(8):
    for j in range(i):
        H = [[basis[i][a][b] + basis[j][a][b]
              for b in range(3)] for a in range(3)]
        entry = (form(H) - matrix[i][i] - matrix[j][j]) / 2
        matrix[i][j] = matrix[j][i] = entry

L = [[I(int(i == j)) for j in range(8)] for i in range(8)]
pivots = []
for i in range(8):
    pivot = matrix[i][i] - sum(
        (L[i][j]*L[i][j]*pivots[j] for j in range(i)), I(0))
    assert pivot.lo > 0, (i + 1, pivot.out())
    pivots.append(pivot)
    for row in range(i + 1, 8):
        numerator = matrix[row][i] - sum(
            (L[row][j]*L[i][j]*pivots[j] for j in range(i)), I(0))
        L[row][i] = numerator / pivot

print("A95 neutral q:", q95.out())
for i, pivot in enumerate(pivots):
    print("A95 LDL pivot", i + 1, pivot.out())
print("A95: D >= (1/50) nu_t |a|^2 |n|^2 for a.n=0.")


# A93 all-direction sufficient bound on the exact neutral root.
q93 = (57 + I(6249).sqrt()) / 40
v1 = I("2.85") / q93
v2 = I("0.15") / q93
q93cube = q93*q93*q93
beta93 = q93cube / (1000 + q93cube)
linear_loss = 15*beta93 + 45*beta93*beta93
lower93 = 1 - v1 - linear_loss*linear_loss / (4*(v1 - v2))
assert lower93.lo > Decimal("0.036")
print("A93 reference lower bound:", lower93.out())

# E94/E95 sharp simple-shear reference bound.
qE = (3 + I(41).sqrt()) / (2*root2)
lowerE = 4 / (qE*qE)
assert lowerE.lo > Decimal("0.3619")
print("E reference lower bound:", lowerE.out())


# A93 q=8: exact negative eigenvalue and positive covariance.
bad93_exact = (
    Fraction(1)
    + Fraction(3, 25)*(Fraction(304, 329) - 4*Fraction(16, 329))
    - 6*Fraction(47, 360)
    - Fraction(18, 5)*Fraction(47, 360)*Fraction(64, 63))
assert bad93_exact == Fraction(-51301, 296100)
assert bad93_exact < 0
rxx = Fraction(86, 81)
ryy = Fraction(290, 567)
rzz = Fraction(242, 567)
rxz = Fraction(-320, 987)
detxz = rxx*rzz - rxz*rxz
assert min(rxx, ryy, rzz, detxz) > 0
assert detxz == Fraction(247166356, 710169201)
print("A93 q=8 acoustic value / nu_t:", bad93_exact)
print("A93 q=8 xz covariance determinant:", detxz)

# A95 q=1: strict radicand, negative transverse value, positive R.
a0 = I("6.5")
Cbad95 = 1 / (a0 + 3/root2)
fbad95 = 1 - I("4.5")*Cbad95*Cbad95
assert fbad95.lo > 0
CQbad95 = fbad95.sqrt() / 4
beta_bad95 = (a0*a0 + 2*a0*(3/root2)).sqrt() / 4
bad95 = 1 - beta_bad95/2
assert bad95.hi < 0
rmin95 = I(2)/3 - (CQbad95*CQbad95 + Cbad95*Cbad95).sqrt()
assert rmin95.lo > 0
assert Cbad95.lo > 0  # Here P_k=C_mu.
print("A95 q=1 acoustic value / nu_t:", bad95.out())
print("A95 q=1 minimum covariance eigenvalue:", rmin95.out())

# E pure strain: exact eigenvalue, covariance and production.
badE = Fraction(-19, 68)
assert badE < 0
assert min(Fraction(4, 51), Fraction(7, 51), Fraction(91, 51)) > 0
assert Fraction(1084, 34) > 0
print("E pure-strain acoustic value / nu_t:", badE)

# A93 pure rotation violates covariance at tau*omega=8.
rotation_ryy = I(2)/3 - I(38)*64/3000
assert rotation_ryy.hi < 0
print("A93 pure rotation R_yy/k:", rotation_ryy.out())

# A93 can have negative production while R remains positive.
root3 = I(3).sqrt()
Cprod = (I(2)/3) / (I(41)/4 + 2*root3)
Dprod = 1000 + 24*root3
r1 = I(2)/3 - 2*Cprod + I(466)/(3*Dprod)
r3 = I(2)/3 + 4*Cprod - I(932)/(3*Dprod)
production93 = 12*Cprod - I(932)/Dprod
assert r1.lo > 0 and r3.lo > 0
assert production93.hi < 0
print("A93 positive-covariance negative production:", production93.out())
print("PASS: all stated interval and exact-sign checks succeeded.")
\end{lstlisting}


\section{Alternative one-coefficient change of dissipation destruction}
\label{out:appendixCeps2}

The main construction changes only $\sigma_\eps$ and retains the
original reaction coefficients. A different one-constant choice is
possible if the original dissipation diffusivity is to be kept instead.
It changes homogeneous dissipation dynamics as well as the equilibrium
balance, so it is an alternative model choice rather than an additional
step in the main construction.

For A93/A95, retain $\sigma_\eps=1.3$ and $C_{\eps1}=1.44$ and impose
\begin{equation}
 \boxed{C_{\eps2,m}^{\rm eq}(\kappa)
 =1.44+\frac{\kappa^2q_m}{1.3},\qquad m=93,95.}
 \label{out:Ceps2A}
\end{equation}
For E94/E95, retain $\sigma_\eps=1.2$ and the original function
$C_{\eps S}(\eta)$ and choose
\begin{equation}
 \boxed{C_{\eps2,E}^{\rm eq}(\kappa)
 =0.43q_E+\frac{\kappa^2q_E}{1.2}.}
 \label{out:Ceps2E}
\end{equation}
The branch value $C_{\eps S,*}=0.43$ remains valid because the stress
law and $q_E$ are unchanged. All three positive-root constitutive
relations and complete neutral covariances are therefore the same as
in \cref{out:roots,out:R0,out:normalcomponents}.

\begin{table}[htbp]
\centering
\caption{Alternative destruction-coefficient calibration with the
original $\sigma_\eps$ retained. This table is not to be combined with
the modified diffusivities in \cref{out:sigmatable}.}
\label{out:Ceps2table}
\begin{tabular}{lrrr}
\toprule
Model & $C_{\eps2,\rm org}$
& $C_{\eps2}^{\rm eq}(0.40)$
& $C_{\eps2}^{\rm eq}(0.41)$\\
\midrule
A93 & 1.92 & 1.8586172821 & 1.8798097820\\
A95 & 1.92 & 1.9104000524 & 1.9342140551\\
E94/E95 & 1.90 & 1.8728053037 & 1.8952457223\\
\bottomrule
\end{tabular}
\end{table}

If this alternative is selected, the net reaction coefficient must be
recomputed from the coefficients actually used:
\begin{equation}
 D_m^{\rm alt}=
 \begin{cases}
 C_{\eps2,m}^{\rm eq}-1.44,&m=93,95,\\
 C_{\eps2,E}^{\rm eq}-0.43q_E,&m=E,
 \end{cases}
 \qquad
 D_m^{\rm alt}=\frac{\kappa^2q_m}{\sigma_{\eps,m,\rm org}}.
 \label{out:Dalternative}
\end{equation}
In particular, the numerical constants $0.48$ and
$1.9-0.43q_E$ from \cref{out:Dm} must not be reused in the altered
reaction balance. The scalar-flux and slab-integral formulas retain
their form with the original $\sigma_\eps$ and the updated
$D_m^{\rm alt}$. The original stress laws remain unaltered, but claims
about the reaction parameters or their calibrated homogeneous-flow
behavior must identify the destruction-coefficient change.

\section{Executable finite-viscosity and wall-coefficient verification}
\label{app:wall-verification}

The following standalone Python program uses Python~3, NumPy, and
SciPy. Copy the listing to a Python file and run it directly. It
checks the original wall coefficients, the optional velocity-overlap
adjustment, and the optional shared-coefficient two-tail adjustment
for both natural and prescribed $\kappa$. Each roughness calibration
is performed once at its own reference; no evolving-state roughness
recalibration is introduced.

Closed wall integrals are compared with independent adaptive
quadrature. Weak source moments are checked by integration by parts
with a separate breakpoint-aware Gauss rule, including nonconstant
Legendre tests through degree six. Smooth-wall, molecular-only,
forced/reversed-flow, full-traction, and matching-height controls
are included. These are local mathematical checks, not flow-solver
or terrain validation. The run used for this note passes 873
assertions over 22 reported cases.

\begin{lstlisting}[caption={Standalone wall-equation and coefficient-consistency checks.}]
#!/usr/bin/env python3
"""Independent finite-nu SCI/EP-AWF identity checks; no CFD solve.
SCI roughness correlation retained; original and explicitly adjusted wall
coefficients are checked with natural and common kappa. Each reference
calibrates h once at K0; no evolving intercept calibration.
Analytic resistance formulas are compared with adaptive quadrature;
Legendre weak moments use separate breakpoint-aware Gauss quadrature.
"""
import json
import math
from dataclasses import dataclass
import numpy as np
from numpy.polynomial.legendre import Legendre, leggauss
from scipy.integrate import quad
from scipy.optimize import brentq

ALPHA=0.2295
CLW=2.55
VS=10.7
KAPPA=0.41
ROOTS={
 'A93':(57+math.sqrt(6249))/40,
 'A95':(3+math.sqrt(61))/(2*math.sqrt(2)),
 'E94_E95':(3+math.sqrt(41))/(2*math.sqrt(2)),
}
COUNT=0

def close(name, actual, expected, atol=3e-11, rtol=3e-10):
 global COUNT
 aa=np.asarray(actual,dtype=float); ee=np.asarray(expected,dtype=float)
 err=float(np.max(np.abs(aa-ee)))
 limit=atol+rtol*float(np.max(np.abs(ee)))
 if not err<=limit: raise AssertionError((name,actual,expected,err,limit))
 COUNT+=1
 return err

def require(name, condition):
 global COUNT
 if not condition:raise AssertionError(name)
 COUNT+=1

def sci_b(h,K,nu):
 if h==0:return VS
 x=h*math.sqrt(K)/(70*nu)
 exponent=max(0.5-0.4*x**0.7,1-0.79*x**(-0.28))
 return VS*(1-x**exponent)

@dataclass
class Kernel:
 K:float
 nu:float
 H:float
 h:float
 alpha:float=ALPHA
 clw:float=CLW
 def __post_init__(self):
  self.a=self.alpha*math.sqrt(self.K)
  self.b=sci_b(self.h,self.K,self.nu)
  self.zv=self.nu*self.b/math.sqrt(self.K)
  self.ze=2*self.clw*self.nu/math.sqrt(self.K)
 def nt(self,z):return self.a*max(0.,z-self.zv)
 def D(self,z):return self.nu+self.nt(z)
 def active(self,z):
  ell=min(max(self.zv,0.),z)
  return ell,self.D(ell),self.D(z)
 def J0(self,z):
  ell,dl,dz=self.active(z)
  return ell/self.nu+(math.log1p((dz-dl)/dl)/self.a if z>ell else 0.)
 def J2(self,z):
  ell,dl,dz=self.active(z)
  return ell/self.nu**2+((1/dl-1/dz)/self.a if z>ell else 0.)
 def J1(self,z):
  if z<=self.h:return 0.
  # Integral from h to the end of the molecular-only part.
  end=min(z,max(self.zv,0.))
  result=max(end-self.h,0.)**2/(2*self.nu)
  start=max(self.h,self.zv,0.)
  if z>start:
   ds=self.D(start);dz=self.D(z)
   result+=((z-start)-(self.h-self.zv+self.nu/self.a)*math.log1p((dz-ds)/ds))/self.a
  return result
 def Ebase(self,z):return self.K**1.5/(self.clw*max(z,self.ze))
 def breaks(self,extra=()):
  return sorted(set([0.,self.H]+[float(x) for x in (self.zv,self.h,self.ze,*extra) if 0<x<self.H]))
 def integral(self,fun,extra=(),upper=None):
  upper=self.H if upper is None else upper
  points=[p for p in self.breaks(extra) if p<upper]+[upper]
  return sum(quad(fun,a,b,epsabs=2e-12,epsrel=2e-12,limit=100)[0] for a,b in zip(points[:-1],points[1:]) if b>a)
 def Pint(self,z,t0):
  # Direct antiderivative of the turbulent, rather than total, work.
  ell,dl,dz=self.active(z)
  if z<=ell:return 0.
  ratio=(dz-dl)/dl
  return t0*t0/self.a*(math.log1p(ratio)-self.nu/dl*ratio/(1+ratio))

def gauss_integral(kernel,fun,extra=(),degree=72):
 x,w=leggauss(degree)
 total=0.
 br=kernel.breaks(extra)
 for a,b in zip(br[:-1],br[1:]):
  z=(a+b)/2+(b-a)*x/2
  total+=(b-a)/2*np.dot(w,[fun(float(v)) for v in z])
 return float(total)

def calibrate_h(q,u,nu,H,z0,kappa=KAPPA,alpha=ALPHA,clw=CLW):
 K=q*u*u; target=math.log(H/z0)/(kappa*u)
 def defect(h):return Kernel(K,nu,H,h,alpha,clw).J0(H)-target
 left,right=0.,H
 dl,dr=defect(left),defect(right)
 endpoint_tol=32*math.ulp(max(target,Kernel(K,nu,H,0.,alpha,clw).J0(H)))
 require('one-time SCI h calibration bracket',
         dl>=-endpoint_tol and dr<0 and target<H/nu)
 h=0. if abs(dl)<=endpoint_tol else brentq(defect,left,right,xtol=1e-14,rtol=2e-14)
 kernel=Kernel(K,nu,H,h,alpha,clw)
 require('physical calibrated roughness and positive production',
         0<=h<H and max(kernel.zv,0.)<H)
 return kernel

def outer_stress(name,q,u):
 K=q*u*u
 if name=='A93':
  fac=q*q/(1000+q**3)
  diag=np.array([2/3+28*fac/3,2/3-11*fac/3,2/3-17*fac/3])*K
 elif name=='A95':
  cq=math.sqrt(1-4.5/q**2)/(1+3*q*q)
  diag=np.array([2/3+cq*q*q,2/3,2/3-cq*q*q])*K
 else:diag=np.full(3,2*K/3)
 R=np.diag(diag);R[0,2]=R[2,0]=-u*u
 return R

def check_reference(label,kernel,q,u,z0=None,model=None,expect_gate=True,kappa=KAPPA):
 H=kernel.H;t0=u*u;A=u**3/kappa
 # Independently integrate all three resistance functions, including partial heights.
 for z in [.13*H,.47*H,H]:
  close(label+' J0',kernel.J0(z),kernel.integral(lambda x:1/kernel.D(x),upper=z))
  close(label+' J2',kernel.J2(z),kernel.integral(lambda x:1/kernel.D(x)**2,upper=z),atol=1e-9)
  close(label+' J1',kernel.J1(z),kernel.integral(lambda x:max(x-kernel.h,0.)/kernel.D(x),upper=z))
 J0=kernel.J0(H);J2=kernel.J2(H);UH=t0*J0
 z0gen=H*math.exp(-kappa*u*J0)
 if z0 is not None:
  close(label+' prescribed z0',z0gen,z0,atol=2e-13)
  close(label+' logarithmic anchor',UH,u/kappa*math.log(H/z0))
 else:z0=z0gen
 p0=lambda z:kernel.nt(z)*(t0/kernel.D(z))**2
 prod=kernel.integral(p0)
 mol=kernel.integral(lambda z:kernel.nu*(t0/kernel.D(z))**2)
 close(label+' production identity',prod,t0*t0*(J0-kernel.nu*J2),atol=4e-12)
 close(label+' positive production antiderivative',prod,kernel.Pint(H,t0))
 close(label+' molecular work',mol,kernel.nu*t0*t0*J2)
 close(label+' reference mean work',t0*UH,prod+mol)
 # SCI's original finite dissipation cutoff and continuous plateau.
 close(label+' baseline cutoff continuity',2*kernel.nu*kernel.K/kernel.ze**2,kernel.K**1.5/(kernel.clw*kernel.ze))
 expected_base=(kernel.K**1.5/kernel.clw*(1+math.log(H/kernel.ze)) if H>=kernel.ze else H*2*kernel.nu*kernel.K/kernel.ze**2)
 close(label+' baseline integral',kernel.integral(kernel.Ebase),expected_base)
 Lambda=prod/A;muwork=kappa*u*kernel.nu*J2
 close(label+' molecular Lambda correction',Lambda,math.log(H/z0)-muwork)
 gate=Lambda>=1+math.log(2.)
 require(label+' admissibility expectation',gate==expect_gate)
 record=dict(case=label,H=H,h=kernel.h,b=kernel.b,zv=kernel.zv,z0=z0,
             J0=J0,production=prod,molecular_loss=mol,Lambda=Lambda,
             molecular_Lambda_correction=muwork,tail_gate_at_H_over_2=gate)
 if not gate:
  if kernel.zv>=H:
   close(label+' all molecular zero turbulent production',prod,0.,atol=1e-14,rtol=0)
   close(label+' all molecular power',t0*UH,mol)
  return record
 zc=H*math.exp(1-Lambda)
 close(label+' finite-nu cutoff',zc,math.e*z0*math.exp(muwork))
 require(label+' finite-nu cutoff exceeds e*z0',zc>math.e*z0)
 Ec=lambda z:A/max(z,zc)
 def Ecint(z):return A*z/zc if z<=zc else A*(1+math.log(z/zc))
 jc=lambda z:Ecint(z)-kernel.Pint(z,t0)
 close(label+' epsilon budget',kernel.integral(Ec,[zc]),prod)
 close(label+' zero redistribution lower',jc(0),0.,atol=1e-13)
 close(label+' zero redistribution upper',jc(H),0.,atol=1e-11)
 # All nonconstant moments: check integration-by-parts realization, not only
 # cancellation of identical pointwise arithmetic expressions.
 raw=[];corrected=[]
 for degree in range(7):
  phi=Legendre.basis(degree);dphi=phi.deriv()
  raw_m=gauss_integral(kernel,lambda z:phi(2*z/H-1)*(p0(z)-Ec(z)),[zc])
  flux_m=-gauss_integral(kernel,lambda z:(2/H)*dphi(2*z/H-1)*jc(z),[zc])
  corrected.append(raw_m+flux_m);raw.append(raw_m)
  close(label+' corrected weak mode '+str(degree),raw_m+flux_m,0.,atol=2e-10,rtol=0)
 require(label+' mean-only correction leaves nonconstant mode',max(abs(x) for x in raw[1:])>1e-6*max(prod,1e-30))
 close(label+' positive corrected dissipation trace',Ec(H),A/H)
 record.update(zc=zc,max_uncorrected_nonconstant_moment=max(abs(x) for x in raw[1:]),
               max_corrected_weak_moment=max(abs(x) for x in corrected))
 if model:
  # Molecular stress belongs to total stress, not Reynolds covariance.
  G=t0/kernel.D(H)
  S=np.zeros((3,3));S[0,2]=S[2,0]=G/2
  RA=2*kernel.K/3*np.eye(3)-2*kernel.nt(H)*S
  Ro=outer_stress(model,q,u)
  piw=Ro[2,2]-2*kernel.K/3
  Tw=-piw*np.eye(3)+2*kernel.nu*S-RA
  To=-Ro
  ez=np.array([0.,0.,1.])
  close(label+' full interface traction',Tw@ez,To@ez)
  close(label+' inner Reynolds shear excludes molecular',RA[0,2],-t0+kernel.nu*G)
  U=np.array([UH,0.,0.])
  close(label+' internal interface work cancels',U@(Tw@ez)+U@(To@(-ez)),0.)
  require(label+' inner reference covariance positive',np.linalg.eigvalsh(RA)[0]>0)
  record['normal_pressure_offset']=float(piw)
 return record

def check_forced(label,kernel,q,u):
 H=kernel.H;t0=u*u;J0=kernel.J0(H);B=kernel.J1(H)/J0
 require(label+' nonzero force transfer',B>0)
 UH=np.array([t0*J0,0.])
 F=np.array([1.5*t0/B,.35*t0/B])
 tw=UH/J0-B*F
 TH=tw+F*(H-kernel.h)
 require(label+' signed reversed ground shear',tw[0]<0)
 def U(z):return tw*kernel.J0(z)+F*kernel.J1(z)
 def grad(z):return (tw+F*max(z-kernel.h,0.))/kernel.D(z)
 PA=lambda z:kernel.nt(z)*np.dot(grad(z),grad(z))
 ML=lambda z:kernel.nu*np.dot(grad(z),grad(z))
 close(label+' no slip',U(0),np.zeros(2))
 close(label+' matching trace',U(H),UH)
 close(label+' total stress derivative/integral',TH-tw,F*(H-kernel.h))
 close(label+' actual upper shear',kernel.D(H)*grad(H),TH)
 prod=kernel.integral(PA);mol=kernel.integral(ML)
 forcepower=kernel.integral(lambda z:np.dot(F,U(z)) if z>kernel.h else 0.)
 upperpower=float(np.dot(UH,TH))
 close(label+' forced molecular+ turbulent work',upperpower-forcepower,prod+mol)
 p0int=kernel.Pint(H,t0)
 src=kernel.integral(lambda z:PA(z)-kernel.nt(z)*(t0/kernel.D(z))**2)
 close(label+' forced source keeps counterfactual p0',src,prod-p0int)
 close(label+' forced combined mean+k energy',upperpower-forcepower-src,p0int+mol)
 # Wrong use of ground shear at the upper interface is a detectable control.
 omitted=float(np.dot(UH,TH-tw))
 require(label+' ground/upper traction confusion detected',abs(omitted)>1e-3*max(abs(upperpower),1e-30))
 # Exact frozen-input force response.
 delta=np.array([.11,-.07]);t2=UH/J0-B*(F+delta)
 close(label+' force response',t2-tw,-B*delta)
 return dict(case=label,wall_shear=tw.tolist(),interface_shear=TH.tolist(),
             production=prod,molecular_loss=mol,forcing_work=forcepower,
             upper_work=upperpower,net_k_source=src,
             omitted_work_if_ground_shear_used=omitted)


def check_kappa_dependencies(label,kernel,q,u,kappa):
 H=kernel.H
 # Centered finite difference independently checks the derivative with H.
 step=1.e-5*H
 def log_z0(hh):return math.log(hh)-kappa*u*kernel.J0(hh)
 derivative=H*(log_z0(H+step)-log_z0(H-step))/(2*step)
 exact=1-kappa*u*H/kernel.D(H)
 close(label+' H-dependence',derivative,exact,atol=2e-9,rtol=2e-8)
 # Active-branch derivative checks the finite molecular/roughness intercept.
 if H>max(kernel.zv,0.):
  close(label+' finite-intercept diffusivity',kernel.D(H),
        kernel.alpha*math.sqrt(kernel.K)*H+kernel.nu*(1-kernel.alpha*kernel.b))
 # The leading slope is tested independently at a high observation height.
 zlarge=max(H,abs(kernel.zv),kernel.nu/kernel.a)*1e8
 effective=kernel.D(zlarge)/(u*zlarge)
 close(label+' overlap slope',effective,kernel.alpha*math.sqrt(q),atol=3e-8,rtol=0)
 return exact

def main():
 results=[];calibrated={}
 H=2.;z0=.03;u=.4;nu=1.5e-5
 for model,q in ROOTS.items():
  D=.48 if model!='E94_E95' else 1.9-.43*q
  sigma_original=1.3 if model!='E94_E95' else 1.2
  natural=math.sqrt(sigma_original*D/q)
  C=(2/3)/(1.25+1.9*q) if model=='A93' else 1/((6.5 if model=='A95' else 4.)+3*q/math.sqrt(2))
  close(model+' constitutive root',C*q*q,1.)
  for choice,kappa in [('natural',natural),('common041',.41)]:
   sigma=kappa*kappa*q/D
   close(model+' '+choice+' outer balance',kappa*kappa*q/sigma-D,0.)
   if choice=='natural':close(model+' unchanged outer sigma',sigma,sigma_original)
   for variant in ['original','slope','both_tails']:
    alpha=ALPHA if variant=='original' else kappa/math.sqrt(q)
    clw=kappa*q**1.5 if variant=='both_tails' else CLW
    kernel=calibrate_h(q,u,nu,H,z0,kappa,alpha,clw)
    label=model+'_'+choice+'_'+variant
    rec=check_reference(label,kernel,q,u,z0,model,kappa=kappa)
    rec.update(kappa=kappa,alpha=alpha,clw=clw,sigma_epsilon=sigma,
               H_log_derivative=check_kappa_dependencies(label,kernel,q,u,kappa))
    if variant!='original':close(label+' matching leading slope',alpha*math.sqrt(q),kappa)
    if variant=='both_tails':
     close(label+' matching baseline tail',kernel.Ebase(H),u**3/(kappa*H))
     close(label+' matching wall cmu',alpha/clw,1/q**2)
    results.append(rec)
    if choice=='common041' and variant=='original':calibrated[model]=kernel
 q=ROOTS['E94_E95'];K=q*u*u
 smooth=Kernel(K,nu,.05,0.)
 smooth_z0=smooth.H*math.exp(-KAPPA*u*smooth.J0(smooth.H))
 smooth_endpoint=calibrate_h(q,u,nu,smooth.H,smooth_z0)
 close('admissible smooth calibration endpoint',smooth_endpoint.h,0.,atol=0.,rtol=0.)
 require('smooth test contains molecular sublayer',0<smooth.zv<smooth.H)
 results.append(check_reference('smooth_molecular_sublayer',smooth,q,u,model='E94_E95'))
 molecular=Kernel(K,nu,5e-5,0.)
 require('all molecular test lies below viscosity intercept',molecular.zv>molecular.H)
 results.append(check_reference('all_molecular_EP_rejected',molecular,q,u,expect_gate=False))
 results.append(check_forced('rough_forced_reversed',calibrated['E94_E95'],q,u))
 results.append(check_forced('smooth_forced_reversed',smooth,q,u))
 print(json.dumps(dict(status='PASS',assertions=COUNT,
  scope='Analytical identities and independent quadrature; no PDE solver or terrain validation.',
  results=results),indent=2))

if __name__=='__main__':main()

\end{lstlisting}

\section{Implementation and review checklist}
\label{app:checklist}

The following items define the information required to reproduce and
review an implementation. A failed item identifies an unresolved
condition; it is not replaced by a visually plausible velocity profile.

\begin{enumerate}
\item \textbf{Outer model identity.} Record A93, A95, or E94/E95;
the tensor and rotation conventions; coefficient branches; original
parameters; and the declared high-Re reduction. Identify any
off-reference extension or limiter separately.
\item \textbf{Outer equilibrium choice.} Record natural or prescribed
$\kappa$, the exact root $q_m$, the active dissipation coefficients,
and whether the diffusivity or the alternative destruction constant
was changed. Verify the full outer equation balance.
\item \textbf{Wall variant.} Record original $\alpha=0.2295$,
overlap-matched $\alpha$, or both-tail matching. Record the original
or adjusted $c_\ell^w$, its baseline cutoff, and the independently
derived corrected-tail coefficient.
\item \textbf{Geometry and sampling.} Record $\nu,H,Z,h_r,z_0,z_m$,
the definition of $\mathcal K$, reference $K_0$, all one-sided
samplers, and which inputs are fixed under mesh or polynomial
refinement. For curved walls verify positive normal stretches,
positive Jacobian, and global normal-tube injectivity.
\item \textbf{Resistance calibration.} Verify the admissible range
and solve the resistance condition once at the fixed reference.
Do not silently recalibrate $h_r$ as $K$ evolves. If $H$ changes,
recheck the inferred $z_0$ and all cutoff conditions.
\item \textbf{Momentum condensation.} Verify ground no slip, the
upper velocity trace, signed reversed shear, forcing cutoff,
and the distinction between ground and upper-interface traction.
\item \textbf{Energy partition.} Use eddy viscosity only in
turbulent production. Account separately for molecular mean-flow
loss and forcing work. Compare analytic integrals with independent
quadrature over every active segment.
\item \textbf{Dependent dissipation.} Obtain means and traces from
the same profile. Apply the finite-viscosity production-budget
cutoff and its gate. Do not independently prescribe both outer
dissipation trace and flux at the matching surface.
\item \textbf{All-mode source.} Check every retained source load,
including nonconstant modes. Use either the explicit redistribution
or its equivalent weak flux, and count it once.
\item \textbf{Pressure and interfaces.} Use a common velocity trace,
actual surface measure and normal, and one shared mass, momentum,
and $k$ flux. Represent the required pressure--normal-stress offset
with one compatible global pressure gauge.
\item \textbf{Transition.} Declare spatial, state, and admissibility
gates and their derivatives. Evaluate the baseline directly where
the correction is off. Record whether the positivity adaptation
remains active at $\chi=0$.
\item \textbf{Positivity.} Verify $K>0$ and coefficient domains.
When the bounded-rate remedy is selected, compute its weighted
budget, capped normalization, destruction, dependent dissipation,
and upper trace together. Test the negative control with a
locally zero $k$ away from a positive sample.
\item \textbf{Constitutive tangent.} Differentiate all active
coefficients. Reproduce the reference certificates and the stated
off-reference negative examples. Positivity of covariance or of
$P_k$ does not replace this test.
\item \textbf{Discrete representation.} Declare the reference
representation, gradient construction, coefficient projection,
quadrature, constraints, and strong or weak boundary conditions.
Measure the actual assembled residual using those operators.
\item \textbf{Temporal and nonlinear solve.} Report one-step
stationarity independently of the spatial residual, and report
true coupled residuals independently of solver-internal estimates.
Qualify pressure, scalar, and wall constraints at every required
stage.
\item \textbf{Physical validation.} Record the actual roughness
mapping, forcing, separation/reattachment, stress and energy
comparisons, and resolution studies. Equilibrium checks alone do
not validate nonequilibrium terrain predictions.
\end{enumerate}


\clearpage
\phantomsection
\addcontentsline{toc}{section}{References}
\begin{thebibliography}{99}
\bibitem[Shih et al.(1993)]{Shih1993}
T.-H. Shih, J. Zhu, and J. L. Lumley.
\newblock \emph{A Realizable Reynolds Stress Algebraic Equation Model}.
\newblock NASA Technical Memorandum 105993, 1993.
\newblock \url{https://ntrs.nasa.gov/citations/19930007407}.

\bibitem[Shih et al.(1995a)]{Shih1995ASM}
T.-H. Shih, J. Zhu, and J. L. Lumley.
\newblock A new Reynolds stress algebraic equation model.
\newblock \emph{Computer Methods in Applied Mechanics and Engineering},
125:287--302, 1995.
\newblock \url{https://doi.org/10.1016/0045-7825(95)00796-4}.

\bibitem[Shih et al.(1994)]{Shih1994EVM}
T.-H. Shih, W. W. Liou, A. Shabbir, Z. Yang, and J. Zhu.
\newblock \emph{A New $k$--$\eps$ Eddy Viscosity Model for High Reynolds
Number Turbulent Flows---Model Development and Validation}.
\newblock NASA Technical Memorandum 106721, 1994.
\newblock \url{https://ntrs.nasa.gov/citations/19950005029}.

\bibitem[Shih et al.(1995b)]{Shih1995EVM}
T.-H. Shih, W. W. Liou, A. Shabbir, Z. Yang, and J. Zhu.
\newblock A new $k$--$\eps$ eddy viscosity model for high Reynolds number
turbulent flows.
\newblock \emph{Computers \& Fluids}, 24(3):227--238, 1995.
\newblock \url{https://doi.org/10.1016/0045-7930(94)00032-T}.

\bibitem[Richards and Hoxey(1993)]{RichardsHoxey1993}
P. J. Richards and R. P. Hoxey.
\newblock Appropriate boundary conditions for computational wind engineering
models using the $k$--$\eps$ turbulence model.
\newblock \emph{Journal of Wind Engineering and Industrial Aerodynamics},
46--47:145--153, 1993.
\newblock \url{https://doi.org/10.1016/0167-6105(93)90124-7}.

\bibitem[Richards and Norris(2011)]{RichardsNorris2011}
P. J. Richards and S. E. Norris.
\newblock Appropriate boundary conditions for computational wind engineering
models revisited.
\newblock \emph{Journal of Wind Engineering and Industrial Aerodynamics},
99(4):257--266, 2011.
\newblock \url{https://doi.org/10.1016/j.jweia.2010.12.008}.

\bibitem[Blocken et al.(2007)]{Blocken2007}
B. Blocken, T. Stathopoulos, and J. Carmeliet.
\newblock CFD simulation of the atmospheric boundary layer:
wall function problems.
\newblock \emph{Atmospheric Environment}, 41(2):238--252, 2007.
\newblock \url{https://doi.org/10.1016/j.atmosenv.2006.08.019}.

\bibitem[Suga et al.(2006)]{SCI2006}
K. Suga, T. J. Craft, and H. Iacovides.
An analytical wall-function for turbulent flows and heat transfer
over rough walls.
\emph{International Journal of Heat and Fluid Flow},
27(5):852--866, 2006.
\url{https://doi.org/10.1016/j.ijheatfluidflow.2006.03.011}.

\bibitem[Cantwell et al.(2015)]{Nektar2015}
C. D. Cantwell et al.,
\emph{Nektar++: An open-source spectral/$hp$ element framework},
Computer Physics Communications \textbf{192}, 205--219 (2015).
\href{https://doi.org/10.1016/j.cpc.2015.02.008}
{doi:10.1016/j.cpc.2015.02.008}.

\bibitem[Moxey et al.(2020)]{Nektar2020}
D. Moxey et al.,
\emph{Nektar++: Enhancing the capability and application of
high-fidelity spectral/$hp$ element methods},
Computer Physics Communications \textbf{249}, 107110 (2020).
\href{https://doi.org/10.1016/j.cpc.2019.107110}
{doi:10.1016/j.cpc.2019.107110}.

\bibitem[Nektar++ developers(2026)]{NektarTheory}
Nektar++ developers,
\emph{Developer Guide, version 5.10.0, Section 14.1:
Fundamental Theories of IncNavierStokesSolver}.
\href{https://doc.nektar.info/developerguide/5.10.0/developer-guidese49.html}
{Official versioned developer guide}.
Accessed 4 November 2026.  This reference documents the framework's
velocity-correction formulation; it is not an implementation of the
present turbulence/wall model.

\end{thebibliography}
\end{document}
